% Chapter 1, Section 1 _Linear Algebra_ Jim Hefferon
%  http://joshua.smcvt.edu/linearalgebra
%  2001-Jun-09
\chapter{Sistem Linear}\leavevmode  % because no preamble?
\section{Menyelesaikan Sistem Linear}
Sistem persamaan linear lazim dijumpai dalam sains dan matematika.
Dua contoh dari pelajaran sains sekolah menengah berikut \cite{Onan}
memberi gambaran tentang bagaimana sistem tersebut muncul.

Contoh pertama berasal dari
\hypertarget{ex:Statics}{Statika}.\index{masalah Statika}
Misalkan kita mempunyai tiga benda,
mengetahui bahwa salah satunya bermassa 2~kg,
dan ingin menentukan dua massa lainnya yang belum diketahui.
Percobaan dengan mistar satu meter menghasilkan dua keadaan setimbang berikut.
\begin{center}
  \includegraphics{gr/mp/ch1.1}
  \qquad
  \includegraphics{gr/mp/ch1.2}
\end{center}
Agar massa-massa tersebut setimbang,
jumlah momen di sebelah kiri harus sama dengan jumlah momen di sebelah
kanan, dengan momen suatu benda didefinisikan sebagai massanya dikalikan jaraknya
dari titik tumpu.
Kondisi ini menghasilkan sistem dua persamaan linear.
\begin{align*}  % odd formatting to match chem problem below
      40h + 15c &= 100  \\
            25c &= 50+50h
\end{align*}

Contoh kedua
berasal dari Kimia.\index{masalah Kimia}
Dalam kondisi terkendali, kita dapat mencampur toluena $\hbox{C}_7\hbox{H}_8$ dan
asam nitrat $\hbox{H}\hbox{N}\hbox{O}_3$ untuk menghasilkan
trinitrotoluena $\hbox{C}_7\hbox{H}_5\hbox{O}_6\hbox{N}_3$
beserta produk samping berupa air
(kondisinya harus dikendalikan dengan sangat ketat\Dash trinitrotoluena
lebih dikenal sebagai TNT).
Dalam perbandingan berapa keduanya harus dicampur?
Jumlah atom setiap unsur sebelum reaksi
\begin{equation*}
    x\, {\rm C}_7{\rm H}_8\ +\ y\,{\rm H}{\rm N}{\rm O}_3
    \quad\longrightarrow\quad
    z\,{\rm C}_7{\rm H}_5{\rm O}_6{\rm N}_3\ +\ w\,{\rm H}_2{\rm O}
\tag*{}\end{equation*}
harus sama dengan jumlahnya sesudah reaksi.
Menerapkan syarat tersebut secara berurutan pada unsur C, H, N, dan O
menghasilkan sistem berikut.
\begin{align*}  % odd formatting to state more naturally
      7x      &= 7z  \\
      8x +1y  &= 5z+2w  \\
      1y      &= 3z  \\
      3y      &= 6z+1w
\end{align*}

Kedua contoh tersebut pada akhirnya bermuara pada penyelesaian suatu sistem persamaan.
Dalam setiap sistem, persamaannya hanya melibatkan pangkat satu dari setiap variabel.
Bab ini menunjukkan cara menyelesaikan sembarang sistem persamaan semacam itu.















\subsection{Metode Gauss}
\begin{definition}\label{def:linearcombination}
%<*df:linearcombination>
Suatu \definend{kombinasi linear}\index{kombinasi linear} dari
\( x_1 \), \ldots, \( x_n \) berbentuk
\begin{equation*}
   a_1x_1+a_2x_2+a_3x_3+\cdots+a_nx_n
\end{equation*}
dengan bilangan \( a_1, \ldots ,a_n\in\Re \) disebut
\definend{koefisien}\index{persamaan linear!koefisien} kombinasi tersebut.
%</df:linearcombination>
%<*df:linearequations>
Suatu \definend{persamaan linear}\index{persamaan linear}
dalam variabel $x_1$, \ldots, $x_n$
berbentuk
$a_1x_1+a_2x_2+a_3x_3+\cdots+a_nx_n=d$
dengan
\( d\in\Re \) sebagai \definend{konstanta}\index{persamaan linear!konstanta}.

Suatu tupel-$n$ \( (s_1,s_2,\ldots ,s_n)\in\Re^n \) merupakan
\definend{solusi}\index{persamaan linear!solusi} %
dari persamaan tersebut, atau dikatakan \definend{memenuhi} persamaan tersebut,
jika penyulihan bilangan
$s_1$, \ldots, $s_n$ untuk variabel-variabelnya
menghasilkan pernyataan yang benar:
$a_1s_1+a_2s_2+\cdots+a_ns_n=d$.
Suatu \definend{sistem persamaan linear}\index{persamaan linear!sistem}%
\index{sistem persamaan linear}
\begin{equation*}
  \begin{linsys}{4}
    a_{1,1}x_1 &+ &a_{1,2}x_2  &+  &\cdots &+ &a_{1,n}x_n &=  &d_1  \\
    a_{2,1}x_1 &+ &a_{2,2}x_2  &+  &\cdots &+ &a_{2,n}x_n &=  &d_2  \\
             &  &           &   &       &  &          &\vdotswithin{=}  \\
    a_{m,1}x_1 &+ &a_{m,2}x_2  &+  &\cdots &+ &a_{m,n}x_n &=  &d_m
  \end{linsys}
\end{equation*}
mempunyai solusi
\( (s_1,s_2,\ldots ,s_n) \) jika tupel-$n$ tersebut merupakan solusi bagi
semua persamaan dalam sistem.
%</df:linearequations>
\end{definition}

\begin{example}
Kombinasi \( 3x_1 + 2x_2 \) dari $x_1$ dan $x_2$ bersifat linear.
Kombinasi \( 3x_1^2 + 2x_2 \) bukan fungsi linear dari
$x_1$ dan $x_2$; demikian pula \( 3x_1 + 2\sin(x_2) \).
\end{example}

Biasanya kita menganggap $x_1$, \ldots, $x_n$ berbeda satu sama lain karena
dalam jumlah yang memuat pengulangan, suku-sukunya dapat disusun ulang
agar setiap variabel muncul sekali, seperti $2x+3y+4x=6x+3y$.
Kadang-kadang kita menyertakan suku berkoefisien nol, seperti dalam
$x-2y+0z$, dan pada kesempatan lain menghilangkannya, sesuai kebutuhan.

\begin{example}
Pasangan terurut \( (-1,5) \) merupakan solusi sistem berikut.
\begin{equation*}
  \begin{linsys}{2}
    3x_1 &+ &2x_2 &= &7  \\
    -x_1 &+ &x_2  &= &6
  \end{linsys}
\end{equation*}
Sebaliknya, \( (5,-1) \) bukan solusi.
\end{example}

Menentukan himpunan semua solusi disebut
\definend{menyelesaikan}\index{sistem persamaan linear!menyelesaikan}
sistem tersebut.
Kita tidak memerlukan tebakan atau keberuntungan;
terdapat suatu algoritma yang selalu berhasil.
Algoritma ini disebut
\definend{Metode Gauss}\index{Metode Gauss}%
\index{sistem persamaan linear!Metode Gauss}
(atau \definend{eliminasi Gauss}\index{eliminasi Gauss}%
\index{sistem persamaan linear!eliminasi Gauss}
atau \definend{eliminasi linear}\index{eliminasi linear}%
\index{sistem persamaan linear!eliminasi linear}%
\index{sistem persamaan linear!eliminasi}%
\index{eliminasi, Gauss}).
% It transforms the system, step by step, into one
% with a form that we can easily solve.
% We will first illustrate how it goes and then we will see the 
% formal statement. 

\begin{example}
Untuk menyelesaikan sistem berikut,
\begin{equation*}
  \begin{linsys}{3}
                     &   &      &   &3x_3  &=  &9  \\
                 x_1 &+  &5x_2  &-  &2x_3  &=  &2  \\
      \frac{1}{3}x_1 &+  &2x_2  &   &      &=  &3  
  \end{linsys}
\end{equation*}
secara bertahap kita mengubahnya sampai berbentuk sistem yang
mudah diselesaikan.

Transformasi pertama
menuliskan ulang sistem dengan menukar baris pertama dan ketiga.
\begin{align*}
  \quad
  &\grstep{ \text{tukar baris 1 dengan baris 3} }
  \begin{linsys}{3}
        \frac{1}{3}x_1 &+  &2x_2  &   &      &=  &3  \\
                   x_1 &+  &5x_2  &-  &2x_3  &=  &2  \\
                       &   &      &   &3x_3  &=  &9  
    \end{linsys}                                         
\end{align*}
Transformasi kedua mengalikan baris pertama dengan faktor~$3$.
\begin{align*}
\quad
  &\grstep{ \text{kalikan baris 1 dengan 3} }
  \begin{linsys}{3}
        x_1 &+  &6x_2  &   &      &=  &9  \\
        x_1 &+  &5x_2  &-  &2x_3  &=  &2  \\
            &   &      &   &3x_3  &=  &9  
   \end{linsys}                                          
\end{align*}
Transformasi ketiga adalah satu-satunya langkah yang tidak langsung dalam contoh ini.
Secara mental, kita mengalikan kedua ruas baris pertama dengan \( -1 \),
menambahkannya pada baris kedua,
dan menuliskan hasilnya sebagai baris kedua yang baru.
\begin{align*}
  &\grstep{ \text{tambahkan \(-1\) kali baris 1 ke baris 2} }
  \begin{linsys}{3}
        x_1 &+  &6x_2  &   &      &=  &9  \\
            &   &-x_2  &-  &2x_3  &=  &-7 \\
            &   &      &   &3x_3  &=  &9  
   \end{linsys}
\end{align*}
Langkah-langkah ini telah membawa sistem ke bentuk yang
memudahkan kita menentukan nilai setiap variabel.
Persamaan terbawah menunjukkan bahwa \( x_3=3 \).
Menyulihkan $3$ untuk \( x_3 \) ke persamaan tengah menunjukkan bahwa \( x_2=1 \).
Menyulihkan kedua nilai tersebut ke persamaan teratas
menghasilkan \( x_1=3 \).
Dengan demikian, sistem mempunyai solusi tunggal;
himpunan solusinya adalah \set{(3,1,3)}.
\end{example}

Kita akan menggunakan Metode Gauss di sepanjang buku ini.
Metode ini cepat dan mudah.
Sekarang akan kita tunjukkan bahwa
metode ini juga aman:
Metode Gauss tidak pernah menghilangkan solusi dan tidak pernah
menambahkan solusi asing. Jadi,
suatu tupel merupakan solusi sistem sebelum metode diterapkan jika dan hanya jika
tupel tersebut juga merupakan solusi sesudahnya.

\begin{theorem}[Metode Gauss]
\index{persamaan linear!solusi!Metode Gauss} \label{th:GaussMethod}
%<*th:GaussMethod>
Jika suatu sistem linear diubah menjadi sistem lain melalui salah satu operasi berikut,
\begin{enumerate}
  \setlength{\itemsep}{0ex}
  \item
    suatu persamaan ditukar dengan persamaan lain
  \item
    kedua ruas suatu persamaan dikalikan dengan konstanta tak nol
  \item
    suatu persamaan diganti dengan jumlah dirinya dan kelipatan persamaan lain
\end{enumerate}
maka kedua sistem mempunyai himpunan solusi yang sama.
%</th:GaussMethod>
\end{theorem}

Masing-masing dari ketiga operasi tersebut memiliki batasan.
Mengalikan suatu baris dengan~\( 0 \) tidak diperbolehkan karena jelas dapat
mengubah himpunan solusi.
Demikian pula, menambahkan kelipatan suatu baris kepada dirinya sendiri tidak diperbolehkan,
karena menambahkan~\( -1 \) kali baris tersebut kepada dirinya sendiri sama dengan
mengalikan baris itu dengan~\( 0 \).
Kita juga tidak memperbolehkan penukaran suatu baris dengan dirinya sendiri,
agar beberapa hasil dalam bab keempat lebih mudah dinyatakan.
Selain itu, operasi tersebut tidak ada gunanya.

\begin{proof}
Di sini kita akan membahas operasi penukaran persamaan.
Dua kasus lainnya
serupa dan menjadi \nearbyexercise{ex:ProveGaussMethod}.

%<*pf:GaussMethod0>
Perhatikan suatu sistem linear.
\begin{equation*}
  \begin{linsys}{4}
    a_{1,1}x_1  &+  &a_{1,2}x_2 &+  &\cdots  &+ &a_{1,n}x_n  &=  &d_1\hfill\hbox{}  \\
                &   &           &   &        &   &     &\vdotswithin{=}  \\
    a_{i,1}x_1  &+  &a_{i,2}x_2 &+  &\cdots  &+ &a_{i,n}x_n  &=  &d_i\hfill\hbox{}  \\
                &   &           &   &        &   &     &\vdotswithin{=}  \\
    a_{j,1}x_1  &+  &a_{j,2}x_2 &+  &\cdots  &+ &a_{j,n}x_n  &=  &d_j\hfill\hbox{}  \\
                &   &           &   &        &   &     &\vdotswithin{=} \\
    a_{m,1}x_1  &+  &a_{m,2}x_2 &+  &\cdots  &+ &a_{m,n}x_n  &=  &d_m  \\
  \end{linsys}
\end{equation*} 
Tupel \( (s_1,\ldots\,,s_n) \)
memenuhi sistem ini
jika dan hanya jika penyulihan nilai-nilai tersebut untuk
variabel, yaitu $s$ untuk $x$, menghasilkan konjungsi pernyataan-pernyataan benar:
$a_{1,1}s_1+a_{1,2}s_2+\cdots+a_{1,n}s_n=d_1$
dan \ldots\ 
$a_{i,1}s_1+a_{i,2}s_2+\cdots+a_{i,n}s_n=d_i$
dan \ldots\  $a_{j,1}s_1+a_{j,2}s_2+\cdots+a_{j,n}s_n=d_j$
dan \ldots\  $a_{m,1}s_1+a_{m,2}s_2+\cdots+a_{m,n}s_n=d_m$.
%</pf:GaussMethod0>

%<*pf:GaussMethod1>
Dalam daftar pernyataan yang dihubungkan dengan kata `dan',
kita dapat mengubah urutan pernyataannya.
Jadi,
syarat ini dipenuhi jika dan hanya jika
$a_{1,1}s_1+a_{1,2}s_2+\cdots+a_{1,n}s_n=d_1$
dan \ldots\  $a_{j,1}s_1+a_{j,2}s_2+\cdots+a_{j,n}s_n=d_j$
dan \ldots\  $a_{i,1}s_1+a_{i,2}s_2+\cdots+a_{i,n}s_n=d_i$
dan \ldots\  $a_{m,1}s_1+a_{m,2}s_2+\cdots+a_{m,n}s_n=d_m$.
Inilah tepatnya syarat agar \( (s_1,\ldots\,,s_n) \)
menyelesaikan sistem setelah penukaran baris.
%</pf:GaussMethod1>
\end{proof}

\begin{definition} \label{df:GaussMethod}
%<*df:GaussMethod>
Ketiga operasi dalam
\nearbytheorem{th:GaussMethod}
disebut
\definend{operasi reduksi
elementer},\index{Metode Gauss!operasi elementer}%
\index{operasi reduksi elementer}
atau \definend{operasi baris}\index{operasi baris elementer},
atau \definend{operasi Gauss}.
Operasi-operasi tersebut ialah
\definend{penukaran}\index{operasi reduksi elementer!penukaran}%
\index{penukaran baris},
\definend{perkalian dengan skalar} (atau
\definend{penskalaan ulang}\index{operasi reduksi elementer!penskalaan ulang}%
\index{penskalaan ulang baris}), dan
\definend{kombinasi baris}\index{operasi reduksi elementer!kombinasi baris}%
\index{pengombinasian baris}\index{penjumlahan baris}.
%</df:GaussMethod>
\end{definition}

Ketika menuliskan perhitungan, kita akan
menyingkat `baris \(i\)' sebagai `\( \rho_i \)'
(ini adalah huruf Yunani rho, yang dilafalkan ``ro'').
Sebagai contoh, operasi kombinasi baris akan kita nyatakan dengan
\( k\rho_i+\rho_j \),
dengan baris yang berubah dituliskan pada urutan kedua.
Untuk menghemat penulisan, kita akan
sering menggabungkan langkah penjumlahan yang menggunakan $\rho_i$ yang sama,
seperti dalam contoh berikut.

\begin{example}
Metode Gauss menerapkan operasi baris secara sistematis untuk menyelesaikan sistem.
Berikut adalah kasus yang lazim.
\begin{equation*}
  \begin{linsys}{3}
    x  &+  &y  &   &   &=  &0  \\
   2x  &-  &y  &+  &3z &=  &3  \\
    x  &-  &2y &-  &z  &=  &3  
  \end{linsys}
\end{equation*}
Kita mulai dengan menggunakan baris pertama untuk
mengeliminasi $2x$ pada baris kedua dan $x$ pada baris ketiga.
Untuk menghilangkan $2x$, secara mental kita mengalikan seluruh baris pertama
dengan $-2$, menambahkannya pada
baris kedua, lalu menuliskan hasilnya sebagai baris kedua yang baru.
Untuk mengeliminasi~$x$ pada baris ketiga, kita mengalikan baris pertama dengan
$-1$, menambahkannya pada baris ketiga, lalu menuliskan hasilnya sebagai
baris ketiga yang baru.
\begin{align*}
  &\grstep[-\rho_1 +\rho_3]{-2\rho_1 +\rho_2}
  \begin{linsys}{3}
     x  &+  &y  &   &   &=  &0  \\
        &   &-3y&+  &3z &=  &3  \\
        &   &-3y&-  &z  &=  &3  
  \end{linsys}   
\end{align*}
% In this version of the system, the last two equations involve only two unknowns.
Kita menyelesaikannya dengan mengubah sistem kedua menjadi sistem ketiga,
yang persamaan terbawahnya hanya melibatkan satu variabel tak diketahui.
Caranya ialah menggunakan
baris kedua untuk mengeliminasi suku~$y$ dari baris ketiga.
\begin{equation*}
  \grstep{-\rho_2 +\rho_3}
  \begin{linsys}{3}
     x  &+  &y  &   &   &=  &0  \\
        &   &-3y&+  &3z &=  &3  \\
        &   &   &   &-4z&=  &0
   \end{linsys}
\end{equation*}
Sekarang solusi sistem mudah ditentukan.
Baris ketiga memberikan \( z=0 \).
Lakukan substitusi balik\index{Metode Gauss!substitusi balik}%
\index{substitusi balik}
ke baris kedua untuk memperoleh \( y=-1 \).
Kemudian lakukan substitusi balik ke baris pertama untuk memperoleh \( x=1 \).
\end{example}

\begin{example}
Untuk masalah Fisika\index{masalah Statika} pada awal
bab ini, Metode Gauss menghasilkan langkah berikut.
\begin{equation*}
   \begin{linsys}{2}
     40h  &+  &15c  &=  &100      \\
     -50h &+  &25c  &=  &50         
   \end{linsys}
   \grstep{5/4\rho_1 +\rho_2}
   \begin{linsys}{2}
      40h  &+  &15c       &=  &100      \\
           &   &(175/4)c  &=  &175 
    \end{linsys}
\end{equation*}
Jadi \( c=4 \), dan substitusi balik memberikan \( h=1 \).
(Masalah Kimia akan kita selesaikan nanti.)
\end{example}

\begin{example}
Reduksi
\begin{align*}
   \begin{linsys}{3}
        x  &+  &y  &+  &z  &=  &9  \\
       2x  &+  &4y &-  &3z &=  &1  \\
       3x  &+  &6y &-  &5z &=  &0  
   \end{linsys}
   &\grstep[-3\rho_1 +\rho_3]{-2\rho_1 +\rho_2}
   \begin{linsys}{3}
      x  &+  &y  &+  &z  &=  &9  \\
         &   &2y &-  &5z &=  &-17\\
         &   &3y &-  &8z&=  &-27
    \end{linsys}                                    \\
   &\grstep{-(3/2)\rho_2+\rho_3}
   \begin{linsys}{3}
      x  &+  &y  &+  &z            &=  &9  \\
         &   &2y &-  &5z           &=  &-17\\
         &   &   &   &-(1/2)z      &=  &-(3/2) 
    \end{linsys}
\end{align*}
menunjukkan bahwa \( z=3 \), \( y=-1 \), dan \( x=7 \).
\end{example}

Sebagaimana diilustrasikan di atas, tujuan Metode Gauss
adalah menggunakan operasi reduksi elementer
untuk menyiapkan substitusi balik.

\begin{definition} \label{df:EchelonForm}
%<*df:EchelonForm>
Pada setiap baris suatu sistem,
variabel pertama yang berkoefisien tak nol disebut
\definend{variabel utama}\index{bentuk eselon!variabel utama}%
\index{utama!variabel} baris tersebut. % 
Suatu sistem berada dalam \definend{bentuk eselon}\index{bentuk eselon}
jika setiap variabel utama
terletak di sebelah kanan variabel utama pada baris di atasnya,
kecuali variabel utama pada baris pertama,
dan semua baris yang seluruh koefisiennya nol berada di bagian bawah.
%</df:EchelonForm>
\end{definition}

\begin{example} 
Tiga contoh sebelumnya hanya menggunakan operasi kombinasi baris.
Sistem linear berikut memerlukan operasi penukaran
untuk membawanya ke bentuk eselon karena
setelah kombinasi pertama,
\begin{align*}
   \begin{linsys}{4}
                x  &-  &y  &   &   &   &   &=  &0  \\
               2x  &-  &2y &+  &z  &+  &2w &=  &4  \\
                   &   &y  &   &   &+  &w  &=  &0  \\
                   &   &   &   &2z &+  &w  &=  &5  
   \end{linsys}
   &\grstep{-2\rho_1 +\rho_2}
   \begin{linsys}{4}
      x  &-  &y  &\spaceforemptycolumn   &   &   &   &=  &0  \\
         &   &   &   &z  &+  &2w &=  &4  \\
         &   &y  &   &   &+  &w  &=  &0  \\
         &   &   &   &2z &+  &w  &=  &5  
   \end{linsys}    
\end{align*}
persamaan kedua tidak memiliki $y$ sebagai variabel utama.
Kita menukarnya dengan baris di bawahnya yang memiliki $y$ sebagai variabel utama.
\begin{align*}
   &\grstep{\rho_2 \leftrightarrow\rho_3}
   \begin{linsys}{4}
      x  &-  &y  &\spaceforemptycolumn   &   &   &   &=  &0  \\
         &   &y  &   &   &+  &w  &=  &0  \\
         &   &   &   &z  &+  &2w &=  &4  \\
         &   &   &   &2z &+  &w  &=  &5  
    \end{linsys}    
\end{align*}
(Seandainya terdapat lebih dari satu baris yang sesuai di bawah baris kedua,
kita dapat menggunakan salah satunya.)
Setelah itu, Metode Gauss dilanjutkan seperti sebelumnya.
\begin{align*}
   &\grstep{-2\rho_3 +\rho_4}
   \begin{linsys}{4}
      x  &-  &y  &\spaceforemptycolumn   &   &   &   &=  &0  \\
         &   &y  &   &   &+  &w  &=  &0  \\
         &   &   &   &z  &+  &2w &=  &4  \\
         &   &   &   &   &   &-3w&=  &-3 
    \end{linsys}
\end{align*}
Substitusi balik memberikan \( w=1 \), \( z=2 \), \( y=-1 \), dan \( x=-1 \).
\end{example}

Tepatnya, untuk menyelesaikan sistem linear kita tidak memerlukan
operasi penskalaan ulang baris.
Operasi itu diperkenalkan di sini karena memudahkan dan karena akan kita gunakan
kemudian dalam bab ini sebagai bagian dari suatu variasi Metode Gauss,
yaitu Metode Gauss--Jordan.

Semua sistem sejauh ini memiliki jumlah persamaan yang sama dengan jumlah variabel
yang belum diketahui. Semuanya memiliki solusi, dan solusi itu tunggal.
Kita akan menutup subbagian ini dengan melihat
kemungkinan-kemungkinan lain.

\begin{example} \label{ex:MoreEqsThanUnks}
Sistem ini
memiliki lebih banyak persamaan daripada variabel.
\begin{equation*}
    \begin{linsys}{2}
      x  &+  &3y  &=  &1  \\
     2x  &+  &y   &=  &-3 \\
     2x  &+  &2y  &=  &-2 
    \end{linsys}  
\end{equation*}
Metode Gauss juga membantu kita memahami sistem ini, sebab langkah berikut
\begin{align*}
    &\grstep[-2\rho_1 +\rho_3]{-2\rho_1 +\rho_2}
    \begin{linsys}{2}
       x  &+  &3y  &=  &1  \\
          &   &-5y &=  &-5 \\
          &   &-4y &=  &-4 
     \end{linsys}
\end{align*}
menunjukkan bahwa salah satu persamaannya redundan.
Bentuk eselon
\begin{equation*}
    \grstep{-(4/5)\rho_2 +\rho_3}
    \begin{linsys}{2}
       x  &+  &3y  &=  &1  \\
          &   &-5y &=  &-5 \\
          &   &0   &=  &0  
     \end{linsys}
\end{equation*}
menghasilkan \( y=1 \) dan \( x=-2 \).
Persamaan `\( 0=0 \)' mencerminkan redundansi tersebut.
\end{example}

Metode Gauss juga berguna untuk sistem yang memiliki lebih banyak variabel daripada persamaan.
Subbagian berikut memuat banyak contohnya.

Perbedaan lain antara sistem linear dan contoh-contoh di atas
ialah bahwa sebagian sistem linear tidak memiliki solusi tunggal.
Hal ini dapat terjadi dengan dua cara.
Pertama, suatu sistem mungkin sama sekali tidak memiliki solusi.

\begin{example} \label{ex:MoreEqsThanUnksInconsis}
Bandingkan sistem pada contoh terakhir dengan sistem berikut.
\begin{equation*}
    \begin{linsys}{2}
      x  &+  &3y  &=  &1  \\
     2x  &+  &y   &=  &-3 \\
     2x  &+  &2y  &=  &0  
    \end{linsys}
    \grstep[-2\rho_1 +\rho_3]{-2\rho_1 +\rho_2}
    \begin{linsys}{2}
       x  &+  &3y  &=  &1  \\
          &   &-5y &=  &-5 \\
          &   &-4y &=  &-2
     \end{linsys}
\end{equation*}
Di sini sistemnya tak konsisten:~tidak ada pasangan bilangan $(s_1,s_2)$ yang
memenuhi ketiga persamaan sekaligus.
Bentuk eselon memperlihatkan ketakkonsistenan itu dengan jelas.
\begin{equation*}
  \grstep{-(4/5)\rho_2 +\rho_3}
  \begin{linsys}{2}
     x  &+  &3y  &=  &1  \\
        &   &-5y &=  &-5 \\
        &   &0   &=  &2 
   \end{linsys}
\end{equation*}
Himpunan solusinya kosong.
\end{example}

\begin{example}
Sistem sebelumnya memiliki lebih banyak persamaan daripada variabel tak diketahui,
tetapi bukan itu penyebab ketakkonsistenannya\Dash
\nearbyexample{ex:MoreEqsThanUnks}
memiliki lebih banyak persamaan daripada variabel tak diketahui, namun tetap konsisten.
Memiliki lebih banyak persamaan daripada variabel tak diketahui juga bukan syarat
terjadinya ketakkonsistenan, sebagaimana terlihat pada sistem tak konsisten berikut yang
memiliki jumlah persamaan dan variabel tak diketahui sama banyak.
\begin{equation*}
  \begin{linsys}{2}
    x  &+  &2y  &=  &8  \\
   2x  &+  &4y  &=  &8  
  \end{linsys}
  \grstep{-2\rho_1 + \rho_2}
  \begin{linsys}{2}
     x  &+  &2y  &=  &8  \\
        &   &0   &=  &-8
   \end{linsys}
\end{equation*}
Sebaliknya,
ketakkonsistenan berkaitan dengan interaksi antara ruas kiri dan ruas kanan;
pada sistem pertama di atas, ruas kiri persamaan kedua adalah dua kali
ruas kiri persamaan pertama, tetapi konstanta kedua di ruas kanan bukan dua kali konstanta pertama.
Kelak kita akan membahas lebih lanjut ketergantungan antara bagian-bagian
suatu sistem.
\end{example}

Cara lain suatu sistem linear gagal memiliki solusi tunggal,
selain karena tidak memiliki solusi,
ialah karena memiliki banyak solusi.

\begin{example}
Dalam sistem berikut,
\begin{equation*}
  \begin{linsys}{2}
    x  &+  &y   &=  &4  \\
   2x  &+  &2y  &=  &8  
  \end{linsys}
\end{equation*}
setiap pasangan bilangan yang memenuhi persamaan pertama juga
memenuhi persamaan kedua.
Himpunan solusi
% \( \{ (x,y)\suchthat x+y=4 \} \) 
\( \set{ (x,y)\suchthat x+y=4} \) 
bersifat tak hingga; beberapa pasangan anggotanya ialah
$(0,4)$, $(-1,5)$, dan $(2.5,1.5)$.

Hasil penerapan Metode Gauss di sini berbeda dari contoh sebelumnya
karena kita tidak memperoleh persamaan yang kontradiktif.
\begin{equation*}
  \grstep{-2\rho_1 + \rho_2}
  \begin{linsys}{2}
    x  &+  &y   &=  &4  \\
       &   &0   &=  &0  
   \end{linsys}
\end{equation*}
\end{example}

Jangan terkecoh oleh contoh tersebut:~persamaan $0=0$
bukan penanda bahwa suatu sistem memiliki banyak solusi.

\begin{example}  \label{ex:NoZerosInfManySols}
Ketiadaan persamaan \( 0=0 \) tidak menghalangi suatu sistem untuk memiliki
banyak solusi yang berbeda.
Sistem berikut berada dalam bentuk eselon,
tidak memuat $0=0$, tetapi memiliki tak hingga banyak solusi,
termasuk $(0,1,-1)$,
$(0,1/2,-1/2)$, $(0,0,0)$, dan $(0,-\pi,\pi)$
(setiap tripel yang komponen pertamanya $0$ dan komponen keduanya merupakan
lawan dari komponen ketiga adalah solusi).
\begin{equation*}
  \begin{linsys}{3}
     x  &+  &y  &+  &z  &=  &0  \\
        &   &y  &+  &z  &=  &0  
   \end{linsys}
\end{equation*}

Adanya \( 0=0 \) juga tidak berarti bahwa sistem harus memiliki
banyak solusi.
Hal tersebut ditunjukkan oleh \nearbyexample{ex:MoreEqsThanUnks}.
Contoh berikut juga menunjukkannya: sistem ini sama sekali tidak memiliki
solusi, walaupun
bentuk eselonnya memuat baris $0=0$.
\begin{align*}
  \begin{linsys}{3}
    2x  &   &   &-   &2z  &=  &6  \\
        &   &y  &+   &z   &=  &1  \\
    2x  &+  &y  &-   &z   &=  &7  \\
        &   &3y &+   &3z  &=  &0  
  \end{linsys}
  &\grstep{-\rho_1 +\rho_3}
  \begin{linsys}{3}
     2x  &\spaceforemptycolumn   &   &-   &2z  &=  &6  \\
         &   &y  &+   &z   &=  &1  \\
         &   &y  &+   &z   &=  &1  \\
         &   &3y &+   &3z  &=  &0  
  \end{linsys}                                    \\
  &\grstep[-3\rho_2 +\rho_4]{-\rho_2 +\rho_3}
  \begin{linsys}{3}
     2x  &\spaceforemptycolumn    &   &-   &2z  &=  &6  \\
         &   &y  &+   &z   &=  &1  \\
         &   &   &    &0   &=  &0  \\
         &   &   &    &0   &=  &-3 
   \end{linsys}
\end{align*}
\end{example}

Singkatnya,
Metode Gauss menggunakan operasi baris untuk
menyiapkan suatu sistem bagi substitusi balik.
Jika suatu langkah memperlihatkan persamaan yang kontradiktif, kita dapat berhenti
dengan kesimpulan bahwa sistem tidak memiliki solusi.
Jika kita mencapai bentuk eselon tanpa persamaan kontradiktif,
dan setiap variabel merupakan variabel utama pada suatu
baris, maka sistem memiliki solusi tunggal yang dapat ditemukan melalui
substitusi balik.
Terakhir, jika kita mencapai bentuk eselon tanpa persamaan kontradiktif,
tetapi tidak terdapat solusi tunggal\Dash
yakni, sekurang-kurangnya satu variabel bukan variabel utama\Dash
maka sistem memiliki banyak solusi.

Subbagian berikut membahas kasus ketiga.
Kita akan melihat bahwa sistem semacam itu pasti memiliki tak hingga banyak solusi,
dan kita akan mendeskripsikan
himpunan solusinya.


\medskip
\noindent\textbf{Catatan.}\hspace*{.2em}
\textit{Dalam latihan-latihan di sini dan di seluruh bagian lain buku ini,
Anda harus memberikan alasan untuk setiap jawaban.
Misalnya, jika suatu soal menanyakan apakah sebuah sistem memiliki solusi,
jawaban ya harus dibenarkan dengan menunjukkan solusinya, sedangkan
jawaban tidak harus dibenarkan dengan menunjukkan bahwa tidak ada solusi.}
\begin{exercises}
  \recommended \item 
    Gunakan Metode Gauss untuk menentukan solusi tunggal setiap sistem.
    \begin{exparts*}
      \partsitem 
        $\begin{linsys}[t]{2}
          2x  &+  &3y  &=  &13  \\
          x   &-  &y   &=  &-1
        \end{linsys}$
      \partsitem 
        $\begin{linsys}[t]{3}
          x   &  &  &-  &z  &=  &0  \\
          3x  &+ &y &   &   &=  &1  \\
          -x  &+ &y &+  &z  &=  &4
        \end{linsys}$
    \end{exparts*}
    \begin{answer}
      \begin{exparts}
        \partsitem Metode Gauss
          \begin{equation*}
            \grstep{-(1/2)\rho_1+\rho_2}
            \begin{linsys}{2}
               2x  &+  &3y      &=  &13  \\
                   &-  &(5/2)y  &=  &-15/2              
            \end{linsys}
          \end{equation*}
          memberikan solusi $y=3$ dan $x=2$.
        \partsitem Metode Gauss di sini
          \begin{equation*}
            \grstep[\rho_1+\rho_3]{-3\rho_1+\rho_2}
            \begin{linsys}{3}
              x   &  &  &-  &z  &=  &0  \\
                  &  &y &+  &3z &=  &1  \\
                  &  &y &   &   &=  &4
            \end{linsys}
            \grstep{-\rho_2+\rho_3}
            \begin{linsys}{3}
              x   &  &  &-  &z    &=  &0  \\
                  &  &y &+  &3z   &=  &1  \\
                  &  &  &   &-3z  &=  &3
            \end{linsys}
          \end{equation*}
          memberikan $x=-1$, $y=4$, dan $z=-1$.
      \end{exparts}
    \end{answer}
  \item
    Setiap sistem berada dalam bentuk eselon.
    Untuk masing-masing sistem, nyatakan apakah terdapat solusi tunggal,
    tidak ada solusi, atau tak hingga banyak solusi.
    \begin{exparts*}
      \partsitem 
        \(\begin{linsys}[t]{2}
           -3x &+ &2y   &= &0  \\          
               &  &-2y  &= &0            
        \end{linsys}\)
      \partsitem 
        \(\begin{linsys}[t]{3}
             x &+ &y &  &   &= &4  \\          
               &  &y &- &z  &= &0  \\           
        \end{linsys}\)
      \partsitem 
        \(\begin{linsys}[t]{3}
             x &+ &y &  &   &= &4  \\          
               &  &y &  &-z  &= &0  \\ 
               &  &  &  &0  &= &0  
        \end{linsys}\)
      \partsitem 
        \(\begin{linsys}[t]{2}
             x &+ &y    &= &4  \\          
               &  &0    &= &4            
        \end{linsys}\)
      \partsitem 
        \(\begin{linsys}[t]{3}
            3x &+ &6y &+  &z  &= &-0.5  \\          
               &  &   &   &-z &= &2.5  \\           
        \end{linsys}\)
      \partsitem 
        \(\begin{linsys}[t]{2}
             x &- &3y    &= &2  \\          
               &  &0    &= &0            
        \end{linsys}\)
      \partsitem 
        \(\begin{linsys}[t]{2}
             2x &+ &2y    &= &4  \\
                &  &y     &= &1 \\          
                 &  &0    &= &4            
        \end{linsys}\)
      \partsitem 
        \(\begin{linsys}[t]{2}
              2x &+ &y   &= &0  
        \end{linsys}\)
      \partsitem 
        \(\begin{linsys}[t]{2}
              x &- &y    &= &-1  \\
                &  &0     &= &0 \\          
                 &  &0    &= &4            
        \end{linsys}\)
      \partsitem 
        \(\begin{linsys}[t]{3}
              x &+ &y  &- &3z  &= &-1  \\
                &  &y  &- &z  &= &2  \\
                &  &  &  &z    &= &0  \\
                &  &  &  &0   &= &0  
        \end{linsys}\)
    \end{exparts*}
    \begin{answer}
      Jika suatu sistem memiliki persamaan kontradiktif, maka sistem tersebut tidak memiliki solusi.
      Jika tidak, dan terdapat variabel yang bukan variabel utama pada baris mana pun,
      maka sistem memiliki tak hingga banyak solusi.
      Dalam kasus terakhir, ketika tidak ada persamaan kontradiktif dan
      setiap variabel menjadi variabel utama pada suatu baris, sistem memiliki solusi tunggal.
      \begin{exparts}
        \partsitem Solusi tunggal
        \partsitem Tak hingga banyak solusi
        \partsitem Tak hingga banyak solusi
        \partsitem Tidak ada solusi
        \partsitem Tak hingga banyak solusi
        \partsitem Tak hingga banyak solusi
        \partsitem Tidak ada solusi
        \partsitem Tak hingga banyak solusi
        \partsitem Tidak ada solusi
        \partsitem Solusi tunggal
      \end{exparts}
    \end{answer}
  \recommended \item  
    Gunakan Metode Gauss untuk menyelesaikan setiap sistem
    atau simpulkan `banyak solusi' atau `tidak ada solusi'.
    \begin{exparts*}
      \partsitem \(
               \begin{linsys}[t]{2}
                  2x  &+  &2y  &=  &5  \\
                   x  &-  &4y  &=  &0  
               \end{linsys}
             \)
      \partsitem \(
               \begin{linsys}[t]{2}
                  -x  &+  &y   &=  &1  \\
                   x  &+  &y   &=  &2  
               \end{linsys}
             \) 
      \partsitem  \(
               \begin{linsys}[t]{3}
                   x  &-  &3y  &+  &z  &=  &1  \\
                   x  &+  &y   &+  &2z &=  &14 
                \end{linsys}
             \) 
      \partsitem  \(
               \begin{linsys}[t]{2}
                  -x  &-  &y   &=  &1  \\
                 -3x  &-  &3y  &=  &2  
               \end{linsys}
             \) 
      \partsitem  \(
               \begin{linsys}[t]{3}
                      &   &4y  &+  &z  &=  &20 \\
                  2x  &-  &2y  &+  &z  &=  &0  \\
                   x  &   &    &+  &z  &=  &5  \\
                   x  &+  &y   &-  &z  &=  &10 
                \end{linsys}
             \)
      \partsitem \( \begin{linsys}[t]{4}
                 2x  &   &   &+  &z  &+  &w  &=  &5  \\
                     &   &y  &   &   &-  &w  &=  &-1 \\
                 3x  &   &   &-  &z  &-  &w  &=  &0  \\
                 4x  &+  &y  &+  &2z &+  &w  &=  &9  
               \end{linsys}
            \)
    \end{exparts*}
    \begin{answer} 
      \begin{exparts}
       \partsitem Reduksi Gauss
        \begin{equation*}
          \grstep{-(1/2)\rho_1+\rho_2}
          \begin{linsys}{2}
             2x  &+  &2y  &=  &5  \\
                 &   &-5y &=  &-5/2  
          \end{linsys}
        \end{equation*}
        menunjukkan bahwa \( y=1/2 \) dan \( x=2 \) merupakan solusi tunggal.
      \partsitem Metode Gauss
        \begin{equation*}
          \grstep{\rho_1+\rho_2}
          \begin{linsys}{2}
             -x  &+  &y   &=  &1  \\
                 &   &2y  &=  &3  
           \end{linsys}
        \end{equation*}
        memberikan \( y=3/2 \) dan \( x=1/2 \) sebagai satu-satunya solusi.
      \partsitem Reduksi baris
        \begin{equation*}
            \grstep{-\rho_1+\rho_2}
            \begin{linsys}{3}
                x  &-  &3y  &+  &z  &=  &1  \\
                   &   &4y  &+  &z  &=  &13 
             \end{linsys}
        \end{equation*}
        menunjukkan bahwa terdapat banyak solusi karena variabel $z$ bukan
        variabel utama pada baris mana pun.
      \partsitem Reduksi baris
        \begin{equation*}
          \grstep{-3\rho_1+\rho_2}
          \begin{linsys}{2}
             -x  &-  &y   &=  &1  \\
                 &   &0   &=  &-1 
           \end{linsys}
        \end{equation*}
        menunjukkan bahwa tidak ada solusi.
      \partsitem Metode Gauss
        \begin{align*}
            \grstep{\rho_1\leftrightarrow\rho_4}
            \begin{linsys}{3}
                x  &+  &y   &-  &z  &=  &10 \\
               2x  &-  &2y  &+  &z  &=  &0  \\
                x  &   &    &+  &z  &=  &5  \\
                   &   &4y  &+  &z  &=  &20 
             \end{linsys}
            &\grstep[-\rho_1+\rho_3]{-2\rho_1+\rho_2}
            \begin{linsys}{3}
                x  &+  &y   &-  &z  &=  &10 \\
                   &   &-4y &+  &3z &=  &-20\\
                   &   &-y  &+  &2z &=  &-5 \\
                   &   &4y  &+  &z  &=  &20 
             \end{linsys}                                      \\
            &\grstep[\rho_2+\rho_4]{-(1/4)\rho_2+\rho_3}
            \begin{linsys}{3}
                x  &+  &y   &-  &z      &=  &10 \\
                   &   &-4y &+  &3z     &=  &-20\\
                   &   &    &   &(5/4)z &=  &0  \\
                   &   &    &   &4z     &=  &0  
             \end{linsys}
        \end{align*}
        memberikan solusi tunggal \( (x,y,z)=(5,5,0) \).
      \partsitem Di sini Metode Gauss memberikan
         \begin{align*}
            &\grstep[-2\rho_1+\rho_4]{-(3/2)\rho_1+\rho_3}
            \begin{linsys}{4}
               2x  &\spaceforemptycolumn   &   &+  &z       &+  &w       &=  &5  \\
                   &   &y  &   &        &-  &w       &=  &-1 \\
                   &   &   &-  &(5/2)z  &-  &(5/2)w  &=  &-15/2  \\
                   &   &y  &   &        &-  &w       &=  &-1  
             \end{linsys}                                             \\ 
            &\grstep{-\rho_2+\rho_4}
            \begin{linsys}{4}
               2x  &\spaceforemptycolumn   &   &+  &z       &+  &w       &=  &5  \\
                   &   &y  &   &        &-  &w       &=  &-1 \\
                   &   &   &-  &(5/2)z  &-  &(5/2)w  &=  &-15/2  \\
                   &   &   &   &        &   &0       &=  &0 
             \end{linsys}
         \end{align*}
         yang menunjukkan bahwa terdapat banyak solusi.
      \end{exparts} 
    \end{answer}
  \item  
    Selesaikan setiap sistem
    atau simpulkan `banyak solusi' atau `tidak ada solusi'.
    Gunakan Metode Gauss.
    \begin{exparts*}
      \partsitem \(
               \begin{linsys}[t]{3}
                   x  &+  &y   &+  &z   &=  &5  \\
                   x  &-  &y   &   &    &=  &0  \\
                      &   &y   &+  &2z  &=  &7  
               \end{linsys}
             \)
      \partsitem \(
               \begin{linsys}[t]{3}
                   3x  &  &    &+  &z   &=  &7  \\
                    x  &-  &y   &+   &3z    &=  &4  \\
                    x   &+  &2y   &-  &5z  &=  &-1  
               \end{linsys}
             \)
      \partsitem \(
               \begin{linsys}[t]{3}
                   x   &+  &3y   &+  &z   &=  &0  \\
                   -x  &-  &y   &   &    &=  &2  \\
                   -x   &+  &y   &+  &2z  &=  &8  
               \end{linsys}
             \)
    \end{exparts*}
    \begin{answer} 
      \begin{exparts}
       \partsitem 
         % sage: M = matrix(RDF, [[1,1,1], [1,-1,0], [0,1,2]])
         % sage: v = vector(RDF, [5,0,7])
         % sage: M_prime = M.augment(v)
         % sage: gauss_method(M_prime)
         % [ 1.0  1.0  1.0  5.0]
         % [ 1.0 -1.0  0.0  0.0]
         % [ 0.0  1.0  2.0  7.0]
         %  take -1.0 times row 1 plus row 2
         % [ 1.0  1.0  1.0  5.0]
         % [ 0.0 -2.0 -1.0 -5.0]
         % [ 0.0  1.0  2.0  7.0]
         %  take 0.5 times row 2 plus row 3
         % [ 1.0  1.0  1.0  5.0]
         % [ 0.0 -2.0 -1.0 -5.0]
         % [ 0.0  0.0  1.5  4.5]
         Metode Gauss
         \begin{align*}
          \begin{linsys}{3}
             x  &+  &y   &+  &z   &=  &5  \\
             x  &-  &y   &   &    &=  &0  \\
                &   &y   &+  &2z  &=  &7  
          \end{linsys}
          &\grstep{-\rho_1+\rho_2} 
          \begin{linsys}{3}
             x  &+  &y   &+  &z   &=  &5  \\
             0  &   &-2y  &-   &z    &=  &-5  \\
                &   &y   &+  &2z  &=  &7  
          \end{linsys}                                     \\
          &\grstep{(1/2)\rho_2+\rho_3} 
          \begin{linsys}{3}
             x  &+  &y   &+  &z   &=  &5  \\
             0  &   &-2y  &-   &z    &=  &-5  \\
                &   &    &+  &(3/2)z  &=  &9/2  
          \end{linsys}
         \end{align*}
         yang diikuti substitusi balik memberikan $x=1$, $y=1$, dan $z=3$.
       \partsitem 
         % sage: M = matrix(QQ, [[3,0,1], [1,-1,3], [1,2,-5]])
         % sage: v = vector(RDF, [7,4,-1])
         % sage: v = vector(QQ, [7,4,-1])
         % sage: M_prime = M.augment(v)
         % sage: gauss_method(M_prime)
         % [ 3  0  1  7]
         % [ 1 -1  3  4]
         % [ 1  2 -5 -1]
         %  take -1/3 times row 1 plus row 2
         %  take -1/3 times row 1 plus row 3
         % [    3     0     1     7]
         % [    0    -1   8/3   5/3]
         % [    0     2 -16/3 -10/3]
         %  take 2 times row 2 plus row 3
         % [  3   0   1   7]
         % [  0  -1 8/3 5/3]
         % [  0   0   0   0]
         Di sini Metode Gauss
         \begin{align*}
           \begin{linsys}{3}
               3x  &  &    &+  &z   &=  &7  \\
                x  &-  &y   &+   &3z    &=  &4  \\
                x   &+  &2y   &-  &5z  &=  &-1  
           \end{linsys}
           &\grstep[-(1/3)\rho_1+\rho_3]{-(1/3)\rho_1+\rho_2} 
           \begin{linsys}{3}
               3x  &  &    &+  &z        &=  &7  \\
                   &  &-y   &+ &(8/3)z   &=  &5/3  \\
                   &  &2y   &- &(16/3)z  &=  &-(10/3)  
           \end{linsys}                                      \\
           &\grstep{2\rho_2+\rho_3} 
           \begin{linsys}{3}
               3x  &  &    &+  &z        &=  &7  \\
                   &  &-y   &+ &(8/3)z   &=  &5/3  \\
                   &  &     &  &0        &=  &0  
           \end{linsys}
       \end{align*}
       menemukan bahwa variabel~$z$ bukan variabel utama suatu baris.
       Terdapat tak hingga banyak solusi.
       \partsitem
         % sage: M = matrix(QQ, [[1,3,1], [-1,-1,0], [-1,1,2]])
         % sage: v = vector(QQ, [0,2,8])
         % sage: M_prime = M.augment(v)
         % sage: gauss_method(M_prime)
         % [ 1  3  1  0]
         % [-1 -1  0  2]
         % [-1  1  2  8]
         %  take 1 times row 1 plus row 2
         %  take 1 times row 1 plus row 3
         % [1 3 1 0]
         % [0 2 1 2]
         % [0 4 3 8]
         %  take -2 times row 2 plus row 3
         % [1 3 1 0]
         % [0 2 1 2]
         % [0 0 1 4]
         Langkah-langkah
         \begin{equation*}
           \begin{linsys}{3}
              x   &+  &3y   &+  &z   &=  &0  \\
              -x  &-  &y   &   &    &=  &2  \\
              -x   &+  &y   &+  &2z  &=  &8  
           \end{linsys}
           \grstep[\rho_1+\rho_3]{\rho_1+\rho_2}
           \begin{linsys}{3}
              x   &+  &3y   &+  &z   &=  &0  \\
                  &   &2y   &+  &z   &=  &2  \\
                  &   &4y   &+  &3z  &=  &8  
           \end{linsys}
           \grstep{-2\rho_2+\rho_3}
           \begin{linsys}{3}
              x   &+  &3y   &+  &z   &=  &0  \\
                  &   &2y   &+  &z   &=  &2  \\
                  &   &     &   &z   &=  &4  
           \end{linsys}
         \end{equation*}
         menghasilkan $(x,y,z)=(-1,-1,4)$.
      \end{exparts} 
    \end{answer}
  \recommended \item 
    Sistem linear dapat diselesaikan dengan metode selain
    Metode Gauss.
    Salah satu metode yang sering diajarkan di sekolah menengah ialah
    menyelesaikan salah satu persamaan untuk suatu variabel, lalu menyulihkan
    ekspresi yang dihasilkan ke persamaan lain.
    Langkah tersebut diulang sampai diperoleh persamaan yang hanya memuat satu
    variabel.
    Dari sana kita memperoleh bilangan pertama dalam solusi, lalu menentukan
    sisanya melalui
    substitusi balik.
    Metode ini lebih lama daripada Metode Gauss karena melibatkan
    lebih banyak operasi aritmetika, dan juga lebih
    rentan menimbulkan kekeliruan.
    Untuk menggambarkan bagaimana metode ini dapat menghasilkan kesimpulan yang salah,
    kita akan menggunakan sistem
    \begin{equation*}
      \begin{linsys}{2}
            x  &+  &3y  &=  &1  \\
            2x  &+  &y   &=  &-3 \\
            2x  &+  &2y  &=  &0  
      \end{linsys}
    \end{equation*}
    dari \nearbyexample{ex:MoreEqsThanUnksInconsis}.
    \begin{exparts}
      \partsitem Selesaikan persamaan pertama untuk $x$, lalu
        sulihkan ekspresi tersebut ke persamaan kedua.
        Tentukan nilai $y$ yang dihasilkan.
      \partsitem Sekali lagi, selesaikan persamaan pertama untuk $x$,
        tetapi kali ini sulihkan ekspresi tersebut ke persamaan ketiga.
        Tentukan nilai $y$ ini.
    \end{exparts}
    Langkah tambahan apa yang harus dilakukan pengguna metode ini agar tidak
    keliru menyimpulkan bahwa suatu sistem memiliki solusi?
    \begin{answer}
      \begin{exparts}
        \partsitem Dari $x=1-3y$ diperoleh $2(1-3y)+y=-3$, sehingga $y=1$.
        \partsitem Dari $x=1-3y$ diperoleh $2(1-3y)+2y=0$, yang menghasilkan
           kesimpulan bahwa $y=1/2$.
      \end{exparts}
      Pengguna metode ini harus memeriksa setiap calon solusi dengan
      menyulihkannya kembali ke semua persamaan.
    \end{answer}
  \recommended \item 
    Untuk nilai \( k \) berapa sistem berikut
    tidak memiliki solusi, memiliki banyak solusi, atau memiliki solusi tunggal?
    \begin{equation*}
       \begin{linsys}{2}
          x  &-  &y  &=  &1  \\
         3x  &-  &3y &=  &k  
       \end{linsys}
    \end{equation*}
    \begin{answer}
      Lakukan reduksi
      \begin{equation*}
        \grstep{-3\rho_1+\rho_2}
        \begin{linsys}{2}
          x  &-  &y  &=  &1\hfill  \\
             &   &0  &=  &-3+k\hfill  
        \end{linsys}
      \end{equation*}
      untuk menyimpulkan bahwa sistem ini tidak memiliki solusi jika \( k\neq 3 \),
      sedangkan jika \( k=3 \), sistem memiliki tak hingga banyak solusi.
      Sistem ini tidak pernah memiliki solusi tunggal.
    \end{answer}
  \item 
    Sistem ini tidak linear karena memuat $\sin\alpha$ sebagai pengganti $\alpha$
    \begin{equation*}
      \begin{linsys}{3}
         2\sin\alpha  &-  &\cos\beta  &+  &3\tan\gamma  &=  &3  \\
         4\sin\alpha  &+  &2\cos\beta &-  &2\tan\gamma  &=  &10  \\
         6\sin\alpha  &-  &3\cos\beta &+  &\tan\gamma   &=  &9  
      \end{linsys}
    \end{equation*}
    namun kita tetap dapat menerapkan Metode Gauss.
    Terapkan metode tersebut.
    Apakah sistem ini memiliki solusi?
    \begin{answer}
      Misalkan \( x=\sin\alpha \), \( y=\cos\beta \), dan \( z=\tan\gamma \):
      \begin{equation*}
        \begin{linsys}{3}
           2x  &-  &y  &+  &3z  &=  &3  \\
           4x  &+  &2y &-  &2z  &=  &10  \\
           6x  &-  &3y &+  &z   &=  &9  
        \end{linsys}
         \grstep[-3\rho_1+\rho_3]{-2\rho_1+\rho_2}
         \begin{linsys}{3}
           2x  &-  &y  &+  &3z  &=  &3  \\
               &   &4y &-  &8z  &=  &4   \\
               &   &   &   &-8z &=  &0  
         \end{linsys}
      \end{equation*}
      menghasilkan \( z=0 \), \( y=1 \), dan \( x=2 \).
      Perhatikan bahwa tidak ada \( \alpha\in\R \) yang memenuhi $\sin\alpha=2$.  
      \end{answer}
  \recommended \item 
    Syarat apa yang harus dipenuhi konstanta-konstanta $b$
    agar setiap sistem berikut memiliki solusi?
    \textit{Petunjuk.} 
    Terapkan Metode Gauss dan amati apa yang terjadi pada ruas kanan.
    \begin{exparts*}
      \partsitem  \(
        \begin{linsys}[t]{2}
           x  &-  &3y  &=  &b_1 \\
          3x  &+  &y   &=  &b_2 \\
           x  &+  &7y  &=  &b_3 \\
          2x  &+  &4y  &=  &b_4 
        \end{linsys}   \)
      \partsitem \(
        \begin{linsys}[t]{3}
           x_1  &+  &2x_2  &+  &3x_3  &=  &b_1  \\
          2x_1  &+  &5x_2  &+  &3x_3  &=  &b_2  \\
           x_1  &   &      &+  &8x_3  &=  &b_3  
        \end{linsys}   \)
    \end{exparts*}
    \begin{answer} 
      \begin{exparts}
       \partsitem Metode Gauss
         \begin{equation*}
           \grstep[-\rho_1+\rho_3 \\ -2\rho_1+\rho_4]{-3\rho_1+\rho_2}
           \begin{linsys}{2}
              x  &-  &3y  &=  &b_1\hfill \\
                 &   &10y &=  &-3b_1+b_2\hfill \\
                 &   &10y &=  &-b_1+b_3\hfill \\
                 &   &10y &=  &-2b_1+b_4\hfill 
            \end{linsys}         
           \grstep[-\rho_2+\rho_4]{-\rho_2+\rho_3}
           \begin{linsys}{2}
              x  &-  &3y  &=  &b_1\hfill \\
                 &   &10y &=  &-3b_1+b_2\hfill \\
                 &   &0   &=  &2b_1-b_2+b_3\hfill \\
                 &   &0   &=  &b_1-b_2+b_4\hfill 
            \end{linsys}
         \end{equation*}
         menunjukkan bahwa sistem ini konsisten jika dan hanya jika
         \( b_3=-2b_1+b_2 \) dan \( b_4=-b_1+b_2 \).
       \partsitem Reduksi
         \begin{align*}
            &\grstep[-\rho_1+\rho_3]{-2\rho_1+\rho_2}
            \begin{linsys}{3}
              x_1  &+  &2x_2  &+  &3x_3  &=  &b_1\hfill  \\
                   &   &x_2   &-  &3x_3  &=  &-2b_1+b_2\hfill  \\
                   &   &-2x_2 &+  &5x_3  &=  &-b_1+b_3\hfill  
             \end{linsys}                                          \\
            &\grstep{2\rho_2+\rho_3}
            \begin{linsys}{3}
              x_1  &+  &2x_2  &+  &3x_3  &=  &b_1\hfill  \\
                   &   &x_2   &-  &3x_3  &=  &-2b_1+b_2\hfill  \\
                   &   &      &   &-x_3  &=  &-5b_1+2b_2+b_3\hfill  
             \end{linsys}
         \end{align*}
         menunjukkan bahwa masing-masing \( b_1 \), \( b_2 \), dan \( b_3 \)
         boleh berupa sebarang bilangan riil\Dash sistem ini selalu memiliki
         solusi tunggal.
      \end{exparts}   
      \end{answer}
  \item 
    Benar atau salah: sistem dengan lebih banyak variabel daripada persamaan
    memiliki sedikitnya satu solusi.
    (Seperti biasa, untuk menyatakan `benar' Anda harus membuktikannya,
    sedangkan untuk menyatakan `salah' Anda harus memberikan contoh tandingan.)
    \begin{answer}
      Sistem berikut memiliki lebih banyak variabel daripada persamaan
      \begin{equation*}
        \begin{linsys}{3}
          x  &+  &y  &+  &z  &=  &0  \\
          x  &+  &y  &+  &z  &=  &1  
        \end{linsys}
      \end{equation*}
      tetapi tidak memiliki solusi.   
      \end{answer}
  \item 
    Apakah setiap soal Kimia\index{masalah Kimia} seperti
    soal yang mengawali subbagian ini\Dash
    yakni soal menyeimbangkan reaksi\Dash harus memiliki tak hingga banyak solusi?
    \begin{answer}
      Ya.
      Sebagai contoh, fakta bahwa reaksi yang sama dapat berlangsung
      dalam dua labu berbeda menunjukkan bahwa dua kali sebarang solusi
      merupakan solusi lain yang berbeda (jika suatu reaksi fisik berlangsung,
      harus ada sedikitnya satu solusi tak nol).
    \end{answer}
  \recommended \item 
    Tentukan koefisien
    \( a \), \( b \), dan \( c \) agar grafik \( f(x)=ax^2+bx+c \) 
    melalui titik \( (1,2) \), \( (-1,6) \), dan \( (2,3) \).
    \begin{answer}
      Karena \( f(1)=2 \), \( f(-1)=6 \), dan \( f(2)=3 \), kita memperoleh
      suatu sistem linear.
      \begin{equation*}
        \begin{linsys}{3}
          1a  &+  &1b  &+  &c  &=  &2  \\
          1a  &-  &1b  &+  &c  &=  &6  \\
          4a  &+  &2b  &+  &c  &=  &3  
         \end{linsys}
      \end{equation*}
      Metode Gauss
      \begin{equation*}
         \grstep[-4\rho_1+\rho_3]{-\rho_1+\rho_2}
         \begin{linsys}{3}
            a  &+  &b  &+  &c  &=  &2  \\
               &   &-2b&   &   &=  &4  \\
               &   &-2b&-  &3c &=  &-5 
          \end{linsys}               
         \grstep{-\rho_2+\rho_3}
         \begin{linsys}{3}
            a  &+  &b  &+  &c  &=  &2  \\
               &   &-2b&   &   &=  &4  \\
               &   &   &   &-3c&=  &-9 
           \end{linsys}
      \end{equation*}
      menunjukkan bahwa solusinya adalah \( f(x)=1x^2-2x+3 \).  
      \end{answer}
  \item Setelah \nearbytheorem{th:GaussMethod}, kita mencatat bahwa mengalikan
   suatu baris dengan~$0$ tidak diperbolehkan karena tindakan tersebut dapat
   mengubah himpunan solusi.
   Berikan contoh sistem dengan himpunan solusi~$S_0$ yang, setelah suatu
   barisnya dikalikan dengan~$0$, menghasilkan sistem baru dengan himpunan
   solusi~$S_1$, dengan $S_0$ merupakan himpunan bagian sejati dari $S_1$,
   yaitu $S_0\neq S_1$.
   Berikan pula contoh dengan $S_0=S_1$. 
   \begin{answer}
     Di sini $S_0=\set{(1,1)}$
     \begin{equation*}
         \begin{linsys}{2}
            x  &+  &y  &=  &2  \\
            x  &-  &y  &=  &0
          \end{linsys}               
         \grstep{0\rho_2}
         \begin{linsys}{2}
            x  &+  &y  &=  &2  \\
               &   &0  &=  &0
          \end{linsys}               
     \end{equation*}
     sedangkan $S_1$ merupakan superhimpunan sejati karena memuat
     sedikitnya dua titik: $(1,1)$ dan~$(2,0)$.
     Dalam contoh berikut, himpunan solusinya tidak berubah.  
     \begin{equation*}
         \begin{linsys}{2}
            x  &+  &y  &=  &2  \\
           2x  &+  &2y &=  &4
          \end{linsys}               
         \grstep{0\rho_2}
         \begin{linsys}{2}
            x  &+  &y  &=  &2  \\
               &   &0  &=  &0
          \end{linsys}               
     \end{equation*}
   \end{answer}
  \item 
    Metode Gauss bekerja dengan menggabungkan persamaan-persamaan dalam suatu
    sistem untuk membentuk persamaan baru.
    \begin{exparts}
      \partsitem Dapatkah kita menurunkan persamaan \( 3x-2y=5 \) melalui
        serangkaian langkah reduksi Gauss dari persamaan-persamaan dalam sistem ini?
        \begin{equation*}
          \begin{linsys}{2}
             x  &+  &y  &=  &1  \\
            4x  &-  &y  &=  &6
          \end{linsys}
        \end{equation*}
      \partsitem Dapatkah kita menurunkan persamaan \( 5x-3y=2 \) melalui
        serangkaian langkah reduksi Gauss dari persamaan-persamaan dalam sistem ini?
        \begin{equation*}
          \begin{linsys}{2}
            2x  &+  &2y &=  &5  \\
            3x  &+  &y  &=  &4
          \end{linsys}
        \end{equation*}
      \partsitem Dapatkah kita menurunkan \( 6x-9y+5z=-2 \)  
        melalui serangkaian langkah reduksi Gauss dari persamaan-persamaan
        dalam sistem ini?
        \begin{equation*}
          \begin{linsys}{3}
            2x  &+  &y  &-  &z  &=  &4  \\
            6x  &-  &3y &+  &z  &=  &5
          \end{linsys}
        \end{equation*}
    \end{exparts}
    \begin{answer} 
       \begin{exparts} 
        \partsitem Ya, melalui pemeriksaan langsung, persamaan yang diberikan
          diperoleh dari
          \( -\rho_1+\rho_2 \).
        \partsitem Tidak.
          Solusi sistem sumber adalah \( (3/4,7/4) \).
          Pasangan ini tidak memenuhi persamaan yang diberikan karena
          penyulihan ke ruas kirinya menghasilkan negatif tiga per dua,
          bukan dua.
        \partsitem Ya.
          Untuk menentukan apakah baris yang diberikan merupakan
          \( c_1\rho_1+c_2\rho_2 \), selesaikan sistem persamaan yang
          menghubungkan koefisien $x$, $y$, $z$, dan suku-suku konstanta:
          \begin{equation*}
            \begin{linsys}{2}
               2c_1  &+  &6c_2  &=  &6  \\
                c_1  &-  &3c_2  &=  &-9 \\
               -c_1  &+  &c_2   &=  &5  \\
               4c_1  &+  &5c_2  &=  &-2 
            \end{linsys}
          \end{equation*}
          dan diperoleh $c_1=-3$ serta $c_2=2$, sehingga baris yang diberikan adalah
          \( -3\rho_1+2\rho_2 \).
      \end{exparts}  
     \end{answer}
  \item 
    Buktikan bahwa, untuk bilangan riil \( a,b,c,d,e \) dengan
    \( a\neq 0 \), jika persamaan linear
    \begin{equation*}
       ax+by=c
    \end{equation*}
    memiliki himpunan solusi yang sama dengan persamaan 
    \begin{equation*}
       ax+dy=e
    \end{equation*}
    maka keduanya merupakan persamaan yang sama.
    Bagaimana jika \( a=0 \)?
    \begin{answer}
      Jika \( a\neq 0 \), himpunan solusi persamaan pertama adalah
      sebagai berikut.
      \begin{equation*}
        \set{(x,y)\in\Re^2\suchthat 
             \text{$x=(c-by)/a=(c/a)-(b/a)\cdot y$}}
        \tag{$*$} 
      \end{equation*}
      Jadi, untuk suatu~$y$, kita dapat menghitung~$x$ yang bersesuaian.
      Dengan mengambil $y=0$, diperoleh solusi $(c/a,0)$; karena persamaan
      kedua $ax+dy=e$ seharusnya memiliki himpunan solusi yang sama,
      penyulihan ke dalam persamaan tersebut memberikan
      $a(c/a)+d\cdot 0=e$, sehingga $c=e$.
      Dengan mengambil $y=1$ dalam ($*$), diperoleh
      $a((c-b)/a)+d\cdot 1=e$, sehingga $b=d$.
      Jadi, keduanya merupakan persamaan yang sama.

      Jika \( a=0 \), kedua persamaan dapat berbeda tetapi tetap memiliki
      himpunan solusi yang sama, misalnya
      \( 0x+3y=6 \) dan \( 0x+6y=12 \).   
     \end{answer}
  \item 
    Tunjukkan bahwa jika \( ad-bc\neq 0 \), maka
    \begin{equation*}
      \begin{linsys}{2}
        ax  &+  &by  &=  &j  \\
        cx  &+  &dy  &=  &k  
      \end{linsys}
    \end{equation*}
    memiliki solusi tunggal.
    \begin{answer}
      Kita meninjau tiga kasus: $a\neq 0$; $a=0$ dan
      $c\neq 0$; serta $a=0$ dan $c=0$.

      Untuk kasus pertama, kita mengasumsikan \( a\neq 0 \).
      Reduksi
      \begin{equation*}
        \grstep{-(c/a)\rho_1+\rho_2}
        \begin{linsys}{2}
          ax  &+  &by                  &=  &j \hfill\hbox{} \\
              &   &(-(cb/a)+d)y  &=  &-(cj/a)+k \hfill  
         \end{linsys}
      \end{equation*}
      menunjukkan bahwa sistem ini memiliki solusi tunggal jika dan hanya jika
      \( -(cb/a)+d\neq 0   \); ingat bahwa \( a\neq 0 \), sehingga
      substitusi balik menghasilkan \( x \) yang tunggal
      (perhatikan bahwa \( j \) dan \( k \) tidak berperan dalam kesimpulan
      bahwa terdapat solusi tunggal, meskipun jika solusi tunggal itu ada,
      keduanya turut menentukan nilainya).
      Akan tetapi, \( -(cb/a)+d = (ad-bc)/a \), dan suatu pecahan tidak sama
      dengan \( 0 \) jika dan hanya jika pembilangnya tidak sama dengan \( 0 \).
      Jadi, dalam kasus pertama ini, terdapat solusi tunggal jika dan hanya jika
      $ad-bc\neq 0$.

      Dalam kasus kedua, jika \( a=0 \) tetapi \( c\neq 0 \), kita menukar
      \begin{equation*}
        \begin{linsys}{2}
          cx  &+  &dy  &=  &k  \\
              &   &by  &=  &j  
        \end{linsys}
      \end{equation*}
      untuk menyimpulkan bahwa sistem memiliki solusi tunggal jika dan hanya jika
      \( b\neq 0 \)
      (kita menggunakan asumsi kasus \( c\neq 0 \) untuk memperoleh
      \( x \) yang tunggal melalui substitusi balik).
      Namun\Dash ketika \( a=0 \) dan \( c\neq 0 \)\Dash
      syarat ``\( b\neq 0 \)''
      setara dengan syarat ``\( ad-bc\neq 0 \)''.
      Dengan demikian, kasus kedua selesai.

      Terakhir, untuk kasus ketiga, jika \( a \) dan \( c \) sama-sama
      bernilai \( 0 \), maka sistem
      \begin{equation*}
        \begin{linsys}{2}
          0x  &+  &by  &=  &j  \\
          0x  &+  &dy  &=  &k  
        \end{linsys}
      \end{equation*}
      tidak memiliki solusi jika kedua kendala pada \(y\) tak konsisten.
      Jika kedua kendala tersebut konsisten, sedikitnya satu nilai y
      memenuhinya dan \(x\) bebas, sehingga terdapat tak hingga banyak solusi.
      Jadi, sistem ini tidak pernah memiliki solusi tunggal.
      Perhatikan bahwa \( a=0 \) dan \( c=0 \) menghasilkan \( ad-bc=0 \).  
    \end{answer}
  \recommended \item 
    Dalam sistem
    \begin{equation*}
      \begin{linsys}{2}
         ax  &+  &by  &=  &c  \\
         dx  &+  &ey  &=  &f  
      \end{linsys}
    \end{equation*}
    masing-masing persamaan menggambarkan sebuah garis pada bidang \( xy \).
    Dengan penalaran geometris, tunjukkan bahwa terdapat tiga kemungkinan:
    ada solusi tunggal, tidak ada solusi, atau ada tak hingga banyak solusi.
    \begin{answer}
      Ingat bahwa jika sepasang garis memiliki dua titik berbeda yang sama,
      maka keduanya merupakan garis yang sama.
      Hal ini karena dua titik menentukan sebuah garis; jadi, kedua titik
      tersebut menentukan masing-masing dari kedua garis itu, sehingga
      keduanya merupakan garis yang sama.

      Dengan demikian, kedua garis dapat memiliki satu titik yang sama
      (menghasilkan solusi tunggal), tidak memiliki titik yang sama
      (menghasilkan tidak ada solusi), atau memiliki sedikitnya dua titik
      yang sama (sehingga keduanya merupakan garis yang sama).  
    \end{answer}
  \item \label{ex:ProveGaussMethod}
    Selesaikan pembuktian \nearbytheorem{th:GaussMethod}.
    \begin{answer}
     Untuk operasi reduksi berupa mengalikan $\rho_i$ dengan bilangan riil
     tak nol $k$, kita memperoleh bahwa \( (s_1,\ldots,s_n) \) memenuhi
     sistem
     \begin{equation*}
       \begin{linsys}{4}
         a_{1,1}x_1  &+  &a_{1,2}x_2 &+  &\cdots  &+  &a_{1,n}x_n  &=  &d_1  \\
                   &   &          &  &        &   &           &\vdotswithin{=}   \\
        ka_{i,1}x_1  &+  &ka_{i,2}x_2 &+  &\cdots  &+  &ka_{i,n}x_n &=  &kd_i  \\
                   &   &           &   &        &   &            &\vdotswithin{=}   \\
         a_{m,1}x_1  &+  &a_{m,2}x_2 &+  &\cdots  &+  &a_{m,n}x_n  &=  &d_m  
       \end{linsys}
     \end{equation*}
     jika dan hanya jika
     % \begin{align*}
     %    a_{1,1}s_1+a_{1,2}s_2+\cdots+a_{1,n}s_n
     %    &=d_1                                              \\
     %    &\alignedvdots                                     \\
     %    \text{and\ } ka_{i,1}s_1+ka_{i,2}s_2+\cdots+ka_{i,n}s_n
     %    &=kd_i                                              \\
     %    &\alignedvdots                                      \\
     %    \text{and\ } a_{m,1}s_1+a_{m,2}s_2+\cdots+a_{m,n}s_n
     %    &=d_m
     % \end{align*}
     $a_{1,1}s_1+a_{1,2}s_2+\cdots+a_{1,n}s_n=d_1$ 
     dan \ldots{} $ka_{i,1}s_1+ka_{i,2}s_2+\cdots+ka_{i,n}s_n=kd_i$
     dan \ldots{} 
     $a_{m,1}s_1+a_{m,2}s_2+\cdots+a_{m,n}s_n=d_m$
     menurut definisi `memenuhi'.
     Karena \( k\neq 0 \), pernyataan tersebut benar jika dan hanya jika
     % \begin{align*}
     %    a_{1,1}s_1+a_{1,2}s_2+\cdots+a_{1,n}s_n
     %    &=d_1                                              \\
     %    &\alignedvdots                                     \\
     %    \text{and\ } a_{i,1}s_1+a_{i,2}s_2+\cdots+a_{i,n}s_n
     %    &=d_i                                              \\
     %    &\alignedvdots                                      \\
     %    \text{and\ } a_{m,1}s_1+a_{m,2}s_2+\cdots+a_{m,n}s_n
     %    &=d_m
     % \end{align*}
     $a_{1,1}s_1+a_{1,2}s_2+\cdots+a_{1,n}s_n=d_1$
     dan \ldots{}
     $a_{i,1}s_1+a_{i,2}s_2+\cdots+a_{i,n}s_n=d_i$
     dan \ldots{}
     $a_{m,1}s_1+a_{m,2}s_2+\cdots+a_{m,n}s_n=d_m$
     (ini hanya pencoretan langsung pada kedua ruas persamaan ke-$i$),
     yang menyatakan bahwa \( (s_1,\ldots,s_n) \) menyelesaikan
     \begin{equation*}
       \begin{linsys}{4}
         a_{1,1}x_1  &+  &a_{1,2}x_2 &+  &\cdots  &+  &a_{1,n}x_n  &=  &d_1  \\
                     &   &           &   &        &   &            &\vdotswithin{=}   \\
         a_{i,1}x_1  &+  &a_{i,2}x_2 &+  &\cdots  &+  &a_{i,n}x_n  &=  &d_i  \\
                     &   &           &   &        &   &            &\vdotswithin{=}   \\
         a_{m,1}x_1  &+  &a_{m,2}x_2 &+  &\cdots  &+  &a_{m,n}x_n  &=
              &d_m  
         \end{linsys}
     \end{equation*}
     sebagaimana dikehendaki.

     Untuk operasi kombinasi $k\rho_i+\rho_j$, tupel
     \( (s_1,\ldots,s_n) \) memenuhi
     \begin{equation*}
       \begin{linsys}{4}
         a_{1,1}x_1             &+  &\cdots  &+  &a_{1,n}x_n  &=  &d_1\hfill\hbox{} \\
                                &   &        &   &            &\vdotswithin{=}   \\
         a_{i,1}x_1             &+  &\cdots  &+  &a_{i,n}x_n  &=  &d_i\hfill\hbox{} \\
                                &   &        &   &            &\vdotswithin{=}   \\
         (ka_{i,1}+a_{j,1})x_1  &+  &\cdots  &+  &(ka_{i,n}+a_{j,n})x_n
               &=  &kd_i+d_j \hfill \\
                                &   &        &   &            &\vdotswithin{=}   \\
         a_{m,1}x_1             &+   &\cdots  &+  &a_{m,n}x_n  &=  
          &d_m\hfill\hbox{} 
        \end{linsys}
     \end{equation*}
     jika dan hanya jika
     $a_{1,1}s_1+\cdots+a_{1,n}s_n=d_1$
     dan \ldots{}
     $a_{i,1}s_1+\cdots+a_{i,n}s_n=d_i$
     dan \ldots{} 
     $(ka_{i,1}+a_{j,1})s_1+\cdots+(ka_{i,n}+a_{j,n})s_n=kd_i+d_j$
     dan \ldots{}
     $a_{m,1}s_1+a_{m,2}s_2+\cdots+a_{m,n}s_n=d_m$
     sekali lagi menurut definisi `memenuhi'.
     Kurangkan \( k \) kali persamaan~\( i \) dari persamaan~\( j \).
     (Di sinilah kita memerlukan \( i\neq j \); jika \( i=j \), baris ke-\(i\)
     semula tidak lagi tersedia sebagai persamaan terpisah.)
     Pernyataan majemuk sebelumnya berlaku jika dan hanya jika
     $a_{1,1}s_1+\cdots+a_{1,n}s_n=d_1$
     dan \ldots{}
     $a_{i,1}s_1+\cdots+a_{i,n}s_n=d_i$
     dan \ldots{} $
     (ka_{i,1}+a_{j,1})s_1+\cdots+(ka_{i,n}+a_{j,n})s_n
         -(ka_{i,1}s_1+\cdots+ka_{i,n}s_n)
         =kd_i+d_j-kd_i$
      dan \ldots{} $a_{m,1}s_1+\cdots+a_{m,n}s_n=d_m$,
     yang setelah pencoretan menyatakan bahwa \( (s_1,\ldots,s_n) \)
     menyelesaikan
     \begin{equation*}
       \begin{linsys}{4}
         a_{1,1}x_1  &+   &\cdots  &+  &a_{1,n}x_n  &=  &d_1  \\
                     &    &        &   &            &\vdotswithin{=}   \\
         a_{i,1}x_1  &+   &\cdots  &+  &a_{i,n}x_n  &=  &d_i  \\
                     &    &        &   &            &\vdotswithin{=}   \\
         a_{j,1}x_1  &+  &\cdots  &+  &a_{j,n}x_n  &=  &d_j  \\
                     &   &        &   &            &\vdotswithin{=}   \\
         a_{m,1}x_1  &+  &\cdots  &+  &a_{m,n}x_n  &=
              &d_m\hfill\hbox{}
       \end{linsys}
     \end{equation*}  
     sebagaimana dikehendaki.
   \end{answer}
  \item 
    Adakah sistem linear dengan dua variabel yang himpunan solusinya
    adalah seluruh \( \Re^2 \)?
    \begin{answer}
      Ya, sistem dengan satu persamaan berikut:
      \begin{equation*}
         0x+0y=0
      \end{equation*}
      dipenuhi oleh setiap \( (x,y)\in\Re^2 \).  
    \end{answer}
  \recommended \item 
    Apakah ada operasi yang digunakan dalam Metode Gauss yang mubazir?
    Dengan kata lain, dapatkah salah satu operasi dibentuk dari kombinasi
    operasi-operasi lainnya?
    \begin{answer}
      Ya.
      Rangkaian operasi berikut menukar baris \( i \) dan \( j \)
      \begin{equation*}
         \grstep{\rho_i+\rho_j}
         \repeatedgrstep{-\rho_j+\rho_i}
         \repeatedgrstep{\rho_i+\rho_j}
         \repeatedgrstep{-1\rho_i}
      \end{equation*}  
      sehingga operasi pertukaran baris mubazir jika dua operasi lainnya tersedia.
     \end{answer}
  \item 
    Buktikan bahwa setiap operasi dalam Metode Gauss dapat dibalik.
    Artinya, tunjukkan bahwa jika dua sistem dihubungkan oleh operasi baris
    $S_1\rightarrow S_2$, maka terdapat operasi baris untuk kembali
    $S_2\rightarrow S_1$.
    \begin{answer}
      Balik pertukaran baris $\rho_i\leftrightarrow\rho_j$ dengan
      menukarnya kembali, $\rho_j\leftrightarrow\rho_i$.
      Balik langkah $k\rho_i$, yaitu mengalikan kedua ruas suatu baris dengan
      \( k\neq 0 \), dengan menerapkan operasi~$(1/k)\rho_i$ pada baris itu.

      Kasus kombinasi baris adalah kasus yang tidak langsung.
      Operasi $k\rho_i+\rho_j$
      menghasilkan baris ke-$j$ berikut.
      \begin{equation*}
        (k\cdot a_{i,1}+a_{j,1})x_1+\cdots+
        (k\cdot a_{i,n}+a_{j,n})x_n=k\cdot d_i+d_j
      \end{equation*}
      Baris ke-$i$ tidak berubah karena pembatasan $i\neq j$.
      Karena baris ke-$i$ tidak berubah, operasi $-k\rho_i+\rho_j$
      mengembalikan baris ke-$j$ ke keadaan semula.

      (Perhatikan bahwa syarat bahwa indeks \( i,j \) berbeda pada
       $k\rho_i+\rho_j$ diperlukan; tanpa syarat itu, hal berikut dapat terjadi
       \begin{equation*}
         \begin{linsys}{2}
           3x  &+  &2y  &=  &7  
         \end{linsys}
         \grstep{2\rho_1+\rho_1}
         \begin{linsys}{2}
           9x  &+  &6y  &=  &21 
          \end{linsys}                 
         \grstep{-2\rho_1+\rho_1}
         \begin{linsys}{2}
          -9x  &-  &6y  &=  &-21 
         \end{linsys}
       \end{equation*}
       sehingga hasil tersebut tidak berlaku.)
    \end{answer}
  \puzzle \item  
    \cite{Anton}
    Sebuah kotak berisi penny, nikel, dan dime, dengan
    tiga belas keping uang logam yang nilai totalnya \( 83 \) sen.
    Berapa banyak keping dari setiap jenis yang ada di dalam kotak?
    (Ini adalah uang logam Amerika Serikat;
    satu penny bernilai $1$~sen, satu nikel $5$~sen, dan
    satu dime $10$~sen.)
    \begin{answer}
      Misalkan \( p \), \( n \), dan \( d \) berturut-turut menyatakan
      banyaknya penny, nikel, dan dime.
      Jika variabel-variabelnya berupa bilangan riil, sistem
      \begin{equation*}
         \begin{linsys}{3}
              p  &+ &n   &+  &d   &=  &13   \\
              p  &+ &5n  &+  &10d &=  &83    
         \end{linsys}
         \grstep{-\rho_1+\rho_2}
         \begin{linsys}{3}
              p  &+ &n   &+  &d   &=  &13   \\
                 &  &4n  &+  &9d  &=  &70    
          \end{linsys}
      \end{equation*}
      memiliki lebih dari satu solusi; bahkan, solusinya tak hingga banyak.
      Akan tetapi, hanya ada sejumlah terbatas solusi dengan \( p \), \( n \),
      dan \( d \) berupa bilangan bulat tak negatif.
      Dengan memeriksa \( d=0 \), \ldots, \( d=8 \), diperoleh bahwa 
      \( (p,n,d)=(3,4,6) \)
      merupakan satu-satunya solusi dalam bilangan asli.  
    \end{answer}
  \puzzle \item 
    \cite{ContestProb1955no38}
    Diberikan empat bilangan bulat positif.
    Pilih sebarang tiga bilangan, tentukan rata-rata aritmetikanya,
    lalu tambahkan hasil tersebut pada bilangan keempat.
    Dengan cara ini diperoleh bilangan 29, 23, 21, dan 17.
    Salah satu bilangan semula adalah:
    \begin{exparts*}
      \partsitem 19
      \partsitem 21
      \partsitem 23
      \partsitem 29
      \partsitem 17
    \end{exparts*}
    \begin{answer}
      Dengan menyelesaikan sistem 
      \begin{equation*}
        \begin{linsys}{2}
        (1/3)(a+b+c)  &+  &d  &=  &29  \\
        (1/3)(b+c+d)  &+  &a  &=  &23  \\
        (1/3)(c+d+a)  &+  &b  &=  &21  \\
        (1/3)(d+a+b)  &+  &c  &=  &17
        \end{linsys}
      \end{equation*}
      kita memperoleh $a=12$, $b=9$, $c=3$, $d=21$.
      Jadi, butir kedua, 21, merupakan jawaban yang benar.
     \end{answer}
  \puzzle \item  
      \cite{Monthly35p47}
      Coba tertawakan penjumlahan ini:
      \( \mbox{AHAHA}+\mbox{TEHE}=\mbox{TEHAW} \).
      Bentuk ini diperoleh dengan mengganti setiap digit dalam suatu contoh
      penjumlahan sederhana dengan huruf kode; tentukan nilai setiap huruf
      dan buktikan bahwa solusinya tunggal.
      \begin{answer}
        \answerasgiven
        Perbandingan kolom satuan dan ratusan dalam penjumlahan ini
        menunjukkan bahwa harus ada simpanan dari kolom puluhan.
        Kolom puluhan kemudian menunjukkan bahwa \( A<H \), sehingga
        tidak mungkin ada simpanan dari kolom satuan ataupun ratusan.
        Kelima kolom tersebut memberikan lima persamaan berikut.
        \begin{align*}
          A+E  &=  W  \\
          2H   &=  A+10  \\
          H    &=  W+1  \\
          H+T  &=  E+10  \\
          A+1  &=  T
        \end{align*}
        Jika diselesaikan secara simultan, kelima persamaan linear dalam
        lima variabel ini menghasilkan solusi tunggal: \( A=4 \), \( T=5 \),
        \( H=7 \), \( W=6 \), dan \( E=2 \), sehingga contoh penjumlahan
        semula adalah \( 47474+5272=52746 \).  
      \end{answer}
  \puzzle \item 
     \cite{Wohascum2}
     Dewan Komisaris Kabupaten Wohascum, yang beranggotakan 20 orang,
     baru-baru ini harus memilih seorang ketua.
     Ada tiga calon ($A$, $B$, dan $C$); pada setiap surat suara, ketiga
     calon harus dicantumkan menurut urutan pilihan, tanpa abstain.
     Ternyata 11 anggota, suatu mayoritas, lebih memilih $A$ daripada $B$
     (jadi, 9 anggota lainnya lebih memilih $B$ daripada $A$).
     Demikian pula, 12 anggota lebih memilih $C$ daripada $A$.
     Berdasarkan hasil ini, diusulkan agar $B$ mengundurkan diri sehingga
     pemilihan putaran kedua dapat diadakan antara $A$ dan $C$.
     Namun, $B$ memprotes, dan kemudian diketahui bahwa 14 anggota justru
     lebih memilih $B$ daripada $C$!
     Dewan tersebut belum juga pulih dari kebingungan yang timbul.
     Jika setiap kemungkinan urutan $A$, $B$, $C$ muncul pada sedikitnya
     satu surat suara, berapa anggota yang menempatkan $B$ sebagai pilihan pertama?
     \begin{answer}
       \answerasgiven{} 
       \textit{Sebagian materi tambahan diambil dari \cite{joriki}.}
       Delapan anggota memilih $B$ sebagai pilihan pertama.
       Untuk melihatnya, kita gunakan informasi yang diberikan guna mengkaji
       banyaknya pemilih untuk setiap urutan $A$, $B$, $C$.

       Keenam urutan pilihan adalah $ABC$, $ACB$, $BAC$, $BCA$, $CAB$,
       $CBA$; misalkan urutan-urutan tersebut berturut-turut memperoleh
       $a$, $b$, $c$, $d$, $e$, $f$ suara.
       Kita mengetahui bahwa
       \begin{equation*}
         \begin{linsys}{3}
           a  &+  &b  &+  &e  &=  &11  \\
           d  &+  &e  &+  &f  &=  &12  \\
           a  &+  &c  &+  &d  &=  &14
         \end{linsys}
       \end{equation*}
       berdasarkan banyaknya orang yang lebih memilih $A$ daripada $B$,
       yang lebih memilih $C$ daripada $A$, dan yang lebih memilih $B$
       daripada $C$.
       Karena terdapat 20 suara,
       $a+b+\cdots+f=20$.
       Mengurangkan jumlah ketiga persamaan di atas dari dua kali persamaan
       sebelumnya menghasilkan $b+c+f=3$.
       Telah ditetapkan bahwa setiap urutan pilihan memperoleh sedikitnya
       satu suara, sehingga $b=c=f=1$. 

       Dari ketiga persamaan di atas, solusi lengkapnya adalah
       $a=6$, $b=1$, $c=1$, $d=7$, $e=4$, dan~$f=1$,
       sebagaimana dapat diperoleh dengan Metode Gauss.
       Jadi, banyaknya anggota yang menempatkan $B$ sebagai pilihan pertama
       adalah $c+d=1+7=8$.

       \par\noindent {\em Catatan.}
       Jawaban soal ini tetap sama seandainya kita hanya mengetahui bahwa
       {\em sedikitnya\/} 14 anggota lebih memilih $B$ daripada $C$.

       Sifat pilihan para anggota yang tampak paradoks
       ($A$ lebih disukai daripada $B$, $B$ lebih disukai daripada $C$, dan
       $C$ lebih disukai daripada $A$), yang merupakan contoh
       ``dominasi tak transitif'', lazim terjadi ketika pilihan individu digabungkan.
     \end{answer}
  \puzzle \item   
     \cite{Monthly63p93}
    ``Sistem yang terdiri atas \( n \) persamaan linear dengan
     \( n \) variabel ini,'' kata Matematikawan Agung, ``memiliki sifat
     yang mengherankan.''

     ``Astaga!'' kata Si Malang, ``Sifat apakah itu?''

     ``Perhatikan,'' kata Matematikawan Agung, ``bahwa konstanta-konstantanya
     membentuk barisan aritmetika.''

     ``Semuanya begitu jelas ketika Anda menjelaskannya!'' kata Si Malang.
     ``Maksud Anda seperti \( 6x+9y=12 \) dan \( 15x+18y=21 \)?''

     ``Tepat sekali,'' kata Matematikawan Agung sambil mengeluarkan fagotnya.
     ``Memang, sistem itu memiliki solusi tunggal.
     Dapatkah engkau menemukannya?''

     ``Astaga!'' seru Si Malang, ``Aku benar-benar bingung.''

     Apakah Anda juga bingung?
     \begin{answer}
       \answerasgiven
       \textit{Kita belum menggunakan istilah ``bergantung'';
       di sini istilah tersebut berarti bahwa Metode Gauss menunjukkan
       tidak adanya solusi tunggal.}
       Jika \( n\geq 3 \), sistem ini bergantung dan solusinya tidak tunggal.
       Jadi, \( n<3 \).
       Namun, istilah ``sistem'' mengisyaratkan \( n>1 \).
       Jadi, \( n=2 \).
       Jika persamaan-persamaannya adalah
       \begin{equation*}
         \begin{linsys}{2}
              ax  &+ &(a+d)y  &=  &a+2d  \\
         (a+3d)x  &+ &(a+4d)y &=  &a+5d  
         \end{linsys}
       \end{equation*}
       maka \( x=-1 \), \( y=2 \).  
    \end{answer}
\end{exercises}







\subsection{Mendeskripsikan Himpunan Solusi}
Sistem linear dengan solusi tunggal memiliki himpunan solusi beranggotakan satu.
Sistem linear tanpa solusi memiliki himpunan solusi kosong.
Dalam kedua kasus tersebut, himpunan solusi mudah dideskripsikan.
Himpunan solusi baru sulit dideskripsikan apabila memuat banyak anggota.

\begin{example}
Sistem ini memiliki banyak solusi karena dalam bentuk eselon
\begin{align*}
  \begin{linsys}{3}
    2x  &   &   &+  &z  &=  &3 \\
     x  &-  &y  &-  &z  &=  &1 \\
    3x  &-  &y  &   &   &=  &4 
  \end{linsys}
  &\grstep[-(3/2)\rho_1 +\rho_3]{-(1/2)\rho_1+\rho_2}
  \begin{linsys}{3}
     2x  &   &   &+  &z      &=  &3    \\
         &   &-y &-  &(3/2)z &=  &-1/2 \\
         &   &-y &-  &(3/2)z &=  &-1/2 
   \end{linsys}                                   \\
  &\grstep{-\rho_2+\rho_3}
  \begin{linsys}{3}
     2x  &   &   &+  &z      &=  &3    \\
         &   &-y &-  &(3/2)z &=  &-1/2 \\
         &   &   &   &0      &=  &0    
   \end{linsys}
\end{align*}
tidak semua variabel merupakan variabel utama.
\nearbytheorem{th:GaussMethod} menunjukkan bahwa suatu $(x,y,z)$
memenuhi sistem pertama jika dan hanya jika memenuhi sistem ketiga.
Jadi, kita dapat mendeskripsikan himpunan solusi
\begin{equation*}
  \set{(x,y,z)\suchthat
    \begin{gathered}
      2x+z=3\text{ dan }x-y-z=1 \\
      \text{dan }3x-y=4
    \end{gathered}}
\end{equation*}
sebagai berikut.
\begin{equation*}
  \set{(x,y,z)\suchthat\text{$2x+z=3$ dan $-y-3z/2=-1/2$}}
  \tag{$*$}
\end{equation*}
Deskripsi ini lebih baik karena hanya memuat dua persamaan, bukan tiga,
tetapi belum optimal karena masih terdapat interaksi antarvariabel yang
sulit dipahami.

Untuk memperbaikinya, gunakan variabel yang tidak menjadi variabel utama
pada persamaan mana pun, yaitu $z$, untuk mendeskripsikan variabel-variabel
utama, yaitu $x$ dan $y$.
Persamaan kedua memberikan
$y=(1/2)-(3/2)z$ 
dan persamaan pertama memberikan
$x=(3/2)-(1/2)z$.
Dengan demikian, himpunan solusi dapat dideskripsikan sebagai himpunan
tupel tiga berikut.
\begin{equation*}
  \set{((3/2)-(1/2)z,\,(1/2)-(3/2)z,\,z)\suchthat z\in\Re}
  \tag{$**$}
\end{equation*}
\end{example}

Dibandingkan dengan ($*$), keunggulan ($**$) adalah bahwa $z$ dapat berupa
sebarang bilangan riil.
Hal ini sangat memudahkan penentuan tupel mana yang berada dalam himpunan
solusi.
Sebagai contoh, mengambil $z=2$ menunjukkan bahwa $(1/2,-5/2,2)$ adalah solusi.

\begin{definition} \label{df:FreeVars}
%<*df:FreeVars>
Dalam sistem linear berbentuk eselon, variabel-variabel yang bukan variabel
utama disebut
\definend{bebas}.\index{bentuk eselon!variabel bebas}\index{variabel bebas}
%</df:FreeVars>
\end{definition}

\begin{example}   \label{ex:Parametrize2}
Reduksi suatu sistem linear dapat berakhir dengan lebih dari satu variabel bebas.
Penerapan Metode Gauss pada sistem ini
\begin{align*}
   \begin{linsys}{4}
               x  &+  &y   &+  &z   &-  &w   &=  &1  \\
                  &   &y   &-  &z   &+  &w   &=  &-1 \\
              3x  &   &    &+  &6z  &-  &6w  &=  &6  \\
                  &   &-y  &+  &z   &-  &w   &=  &1  
   \end{linsys}
  &\grstep{-3\rho_1 +\rho_3}
  \begin{linsys}{4}
     x  &+  &y   &+  &z   &-  &w   &=  &1  \\
        &   &y   &-  &z   &+  &w   &=  &-1 \\
        &   &-3y &+  &3z  &-  &3w  &=  &3  \\
        &   &-y  &+  &z   &-  &w   &=  &1  
  \end{linsys}                                      \\
  &\grstep[\rho_2 +\rho_4]{3\rho_2 +\rho_3}
  \begin{linsys}{4}
     x  &+  &y   &+  &z   &-  &w   &=  &1  \\
        &   &y   &-  &z   &+  &w   &=  &-1 \\
        &   &    &   &    &   &0   &=  &0  \\
        &   &    &   &    &   &0   &=  &0  
   \end{linsys}
\end{align*}
menyisakan \( x \) dan \( y \) sebagai variabel utama, sedangkan \( z \)
dan~\( w \) bebas.
Untuk memperoleh deskripsi yang kita inginkan, kita bekerja dari bawah.
Pertama, nyatakan variabel utama $y$ dalam bentuk $z$ dan $w$, yaitu
$y=-1+z-w$.
Bergerak ke persamaan teratas, penyulihan untuk $y$ memberikan
$x+(-1+z-w)+z-w=1$, dan penyelesaian terhadap $x$ memberikan $x=2-2z+2w$.
Himpunan solusi
\begin{equation*}
   \set{(2-2z+2w,-1+z-w,z,w)\suchthat z,w\in\Re}
  \tag{$**$}
\end{equation*}
menyatakan variabel-variabel utama dalam bentuk variabel-variabel bebas.
\end{example}

\begin{example}    \label{ex:Parametrize1}
Daftar variabel utama dapat melewati beberapa kolom.
Setelah reduksi berikut
\begin{align*}
  \begin{linsys}{4}
              2x  &-  &2y  &   &    &   &    &=  &0  \\
                  &   &    &   &z   &+  &3w  &=  &2  \\
              3x  &-  &3y  &   &    &   &    &=  &0  \\
               x  &-  &y   &+  &2z  &+  &6w  &=  &4  
  \end{linsys}
  &\grstep[-(1/2)\rho_1+ \rho_4]{-(3/2)\rho_1 +\rho_3}
  \begin{linsys}{4}
     2x  &-  &2y  &\spaceforemptycolumn   &    &   &    &=  &0  \\
         &   &    &   &z   &+  &3w  &=  &2  \\
         &   &    &   &    &   &0   &=  &0  \\
         &   &    &   &2z  &+  &6w  &=  &4  
   \end{linsys}                                    \\
  &\grstep{-2\rho_2 +\rho_4}
  \begin{linsys}{4}
     2x  &-  &2y  &\spaceforemptycolumn   &    &   &    &=  &0  \\
         &   &    &   &z   &+  &3w  &=  &2  \\
         &   &    &   &    &   &0   &=  &0  \\
         &   &    &   &    &   &0   &=  &0  
   \end{linsys}
\end{align*}
$x$ dan $z$ merupakan variabel utama, bukan $x$ dan~$y$.
Variabel bebasnya adalah $y$ dan~$w$, sehingga himpunan solusi dapat
dideskripsikan sebagai
$\set{ (y,y,2-3w,w)\suchthat y,w\in\Re }$.
Sebagai contoh, \( (1,1,2,0) \) memenuhi sistem\Dash ambil
$y=1$ dan $w=0$.
Tupel empat \( (1,0,5,4) \) bukan solusi karena koordinat pertamanya
tidak sama dengan koordinat keduanya.
\end{example}

%<*df:Parameter>
Variabel yang digunakan untuk mendeskripsikan suatu keluarga solusi
disebut \definend{parameter}.\index{parameter}  
%</df:Parameter>
Kita mengatakan bahwa himpunan solusi dalam contoh sebelumnya
\definend{diparameterkan\/}\index{diparameterkan}
dengan $y$ dan $w$.

Istilah `parameter' dan `variabel bebas' tidak memiliki arti yang sama.
Dalam contoh sebelumnya, $y$ dan~$w$ bebas karena keduanya bukan variabel
utama dalam sistem berbentuk eselon.
Keduanya merupakan parameter karena kita menggunakannya untuk
mendeskripsikan himpunan solusi.
Seandainya persamaan kedua ditulis ulang sebagai $w=2/3-(1/3)z$, variabel
bebasnya tetap $y$ dan~$w$, tetapi parameternya menjadi $y$ dan~$z$.

Dalam bagian-bagian selanjutnya dari buku ini, kita akan menyelesaikan sistem linear
dengan membawanya ke bentuk eselon, lalu memparameterkannya menggunakan
variabel bebas.

\begin{example}
Berikut adalah sistem lain dengan tak hingga banyak solusi.
\begin{align*}
   \begin{linsys}{4}
               x  &+  &2y  &   &   &   &   &=  &1  \\
              2x  &   &    &+  &z  &   &   &=  &2  \\
              3x  &+  &2y  &+  &z  &-  &w  &=  &4  
   \end{linsys}
  &\grstep[-3\rho_1 +\rho_3]{-2\rho_1+\rho_2}
  \begin{linsys}{4}
     x  &+  &2y  &   &   &   &   &=  &1  \\
        &   &-4y &+  &z  &   &   &=  &0  \\
        &   &-4y &+  &z  &-  &w  &=  &1  
   \end{linsys}                                    \\
  &\grstep{-\rho_2+\rho_3}
  \begin{linsys}{4}
     x  &+  &2y  &   &   &   &   &=  &1  \\
        &   &-4y &+  &z  &   &   &=  &0  \\
        &   &    &   &   &   &-w &=  &1  
   \end{linsys}
\end{align*}
Variabel utamanya adalah \( x \), \( y \), dan \( w \).
Variabel \( z \) bebas.
Perhatikan bahwa, meskipun solusinya tak hingga banyak, nilai $w$ tidak
berubah, melainkan tetap $w=-1$.
Untuk memparameterkan, nyatakan \( w \) sebagai fungsi \( z \), yaitu \( w=-1+0z \).
Kemudian \( y=(1/4)z \).
Sulihkan \( y \) ke persamaan pertama untuk memperoleh \( x=1-(1/2)z \).
Himpunan solusinya adalah
$\set{(1-(1/2)z,(1/4)z,z,-1)\suchthat z\in\Re}$.
\end{example}

Memparameterkan himpunan solusi menunjukkan bahwa sistem dengan variabel
bebas memiliki tak hingga banyak solusi.
Sebagai contoh, di atas $z$ dapat mengambil setiap nilai riil, dan setiap
nilai tersebut berkaitan dengan solusi berbeda.

Kita menutup subbagian ini dengan mengembangkan notasi yang lebih ringkas
untuk sistem linear dan himpunan solusinya.

\begin{definition}\label{df:matrix}
%<*df:matrix>
Suatu \( \nbym{m}{n} \) \definend{matriks}\index{matriks}
adalah susunan bilangan berbentuk persegi panjang dengan
\( m \)~\definend{baris}\index{matriks!baris}\index{baris}
dan \( n \)~\definend{kolom}\index{matriks!kolom}\index{kolom}.
Setiap bilangan dalam matriks disebut
\definend{entri}\index{matriks!entri}\index{entri, matriks}.
%</df:matrix>
\end{definition}

Biasanya, matriks dinotasikan dengan huruf Latin kapital.
Sebagai contoh,
\begin{equation*}
  A=
  \begin{mat}[r]
    1  &2.2  &5  \\
    3  &4    &-7
  \end{mat}
\end{equation*}
memiliki $2$~baris dan $3$~kolom, sehingga merupakan matriks
\( \nbym{2}{3} \).
Bacalah sebagai ``dua kali tiga''; banyaknya baris selalu disebutkan lebih dahulu.
(Matriks diberi tanda kurung agar, ketika dua matriks bersebelahan, kita dapat
mengetahui tempat berakhirnya matriks yang satu dan dimulainya matriks yang lain.)
Entri matriks dinamai dengan huruf kecil yang bersesuaian, sehingga entri pada
baris kedua dan kolom pertama susunan di atas adalah \( a_{2,1}=3 \).
Perhatikan bahwa urutan subskrip penting:
$a_{1,2}\neq a_{2,1}$ karena \( a_{1,2}=2.2 \).
Himpunan semua matriks \( \nbym{m}{n} \) dinotasikan dengan
\( \matspace_{\nbym{m}{n}} \).\index{matriks!himpunan}

Kita menerapkan Metode Gauss pada matriks dengan cara yang pada dasarnya sama
seperti pada sistem persamaan:
\definend{entri utama}\index{bentuk eselon!entri utama}%
\index{utama!entri} %
suatu baris matriks adalah entri tak nol pertamanya (jika ada), dan kita
melakukan operasi baris untuk memperoleh
\definend{bentuk eselon matriks},\index{bentuk eselon!matriks}%
\index{matriks!bentuk eselon} %
yang entri utama pada baris-baris bawahnya terletak di sebelah kanan entri utama
pada baris-baris di atasnya.
Kita menyukai notasi matriks karena notasi ini meringankan pekerjaan mekanis:
menyalin variabel serta menuliskan tanda $+$ dan $=$.

\begin{example}
Kita dapat menyingkat sistem linear
\begin{equation*}
  \begin{linsys}{3}
    x  &+  &2y  &   &    &=  &4   \\
       &   &y   &-  &z   &=  &0   \\
    x  &   &    &+  &2z  &=  &4   
  \end{linsys}
\end{equation*}
dengan matriks berikut.
\begin{equation*}
    \begin{amat}[r]{3}
      1  &2  &0  &4  \\
      0  &1  &-1 &0  \\
      1  &0  &2  &4
    \end{amat}
\end{equation*}
Garis vertikal mengingatkan pembaca akan perbedaan antara koefisien pada ruas
kiri sistem dan konstanta pada ruas kanannya.
Dengan garis tersebut, matriks ini disebut matriks
\definend{teraugmentasi\/}\index{matriks!teraugmentasi}
\index{matriks teraugmentasi}.
\begin{equation*}
    \begin{amat}[r]{3}
      1  &2  &0  &4  \\
      0  &1  &-1 &0  \\
      1  &0  &2  &4
    \end{amat}
  \grstep{-\rho_1 +\rho_3}
  \begin{amat}[r]{3}
       1  &2  &0  &4  \\
       0  &1  &-1 &0  \\
       0  &-2 &2  &0
     \end{amat}                        
  \grstep{2\rho_2 +\rho_3}
  \begin{amat}[r]{3}
       1  &2  &0  &4  \\
       0  &1  &-1 &0  \\
       0  &0  &0  &0
     \end{amat}
\end{equation*}
Baris kedua menyatakan $y-z=0$ dan baris pertama menyatakan
$x+2y=4$, sehingga himpunan solusinya adalah
\( \set{(4-2z,z,z)\suchthat z\in\Re} \).
\end{example}

Notasi matriks juga memperjelas deskripsi himpunan solusi.
\nearbyexample{ex:Parametrize2} memuat himpunan
$\set{(2-2z+2w,-1+z-w,z,w)\suchthat z,w\in\Re}$ 
yang sulit dibaca.
Kita akan menulisnya ulang dengan mengelompokkan semua konstanta, semua
koefisien~\( z \), dan semua koefisien~\( w \).
Kita menuliskannya secara vertikal dalam matriks satu kolom.
\begin{equation*}
  \set{\colvec[r]{2 \\ -1 \\ 0 \\ 0}
       +\colvec[r]{-2 \\ 1 \\ 1 \\ 0}\cdot z
       +\colvec[r]{2 \\ -1 \\ 0 \\ 1}\cdot w
       \suchthat z,w\in\Re}
\end{equation*}
Sebagai contoh, baris teratas menyatakan bahwa \( x=2-2z+2w \), dan baris
kedua menyatakan bahwa \( y= -1+z-w \).
(Bagian berikutnya memberikan interpretasi geometris yang membantu kita
membayangkan himpunan solusi.)

\begin{definition}\label{df:vector}
%<*df:vector>
Suatu
\definend{vektor kolom},\index{kolom!vektor}\index{vektor!kolom}
yang sering cukup disebut \definend{vektor},\index{vektor}
adalah matriks dengan satu kolom.
Matriks dengan satu baris disebut
\definend{vektor baris}\index{baris!vektor}\index{vektor!baris}.
Entri-entri suatu vektor kadang-kadang disebut
\definend{komponen}\index{komponen vektor}\index{vektor!komponen}.
Vektor kolom atau baris yang semua komponennya nol disebut
\definend{vektor nol}.\index{vektor nol}\index{vektor!nol}
%</df:vector>
\end{definition}

Vektor merupakan pengecualian terhadap kebiasaan merepresentasikan matriks
dengan huruf Latin kapital.
Kita menggunakan huruf Latin atau Yunani kecil yang diberi panah di atasnya:
\( \vec{a} \), \( \vec{b} \), \ldots\, atau
\( \vec{\alpha} \), \( \vec{\beta} \), \ldots\
(huruf tebal juga lazim digunakan:
{\boldmath \( a \)} atau {\boldmath \( \alpha \)}).
Sebagai contoh, berikut adalah vektor kolom yang komponen ketiganya \( 7 \).
\begin{equation*}
  \vec{v}=
  \colvec[r]{ 1  \\  3  \\ 7}
\end{equation*}
Vektor nol dinotasikan dengan \( \zero \).
Ada banyak vektor nol yang berbeda\Dash vektor nol dengan satu entri, vektor
nol dengan dua entri, dan seterusnya\Dash tetapi kita tetap sering mengatakan
``vektor nol tersebut'', dengan mengharapkan ukurannya jelas dari konteks.

\begin{definition}
Persamaan linear
\( a_1x_1+a_2x_2+\,\cdots\,+a_nx_n=d \)
dengan variabel tak diketahui \( x_1,\ldots\,,x_n \)
\definend{dipenuhi}\index{vektor!memenuhi persamaan}%
\index{persamaan linear!dipenuhi oleh vektor} oleh
\begin{equation*}
  \vec{s}=\colvec{s_1 \\ \vdotswithin{s_1} \\ s_n}
\end{equation*}
jika \( a_1s_1+a_2s_2+\,\cdots\,+a_ns_n=d \).
Suatu vektor memenuhi sistem linear jika memenuhi setiap persamaan dalam sistem.
\end{definition}

Gaya deskripsi himpunan solusi yang kita gunakan melibatkan penjumlahan
vektor dan perkaliannya dengan bilangan riil.
Sebelum memberikan contoh gaya tersebut, kita perlu mendefinisikan
operasi-operasi ini.

\begin{definition} \label{df:VectorSum}
%<*df:VectorSum>
\definend{Jumlah vektor}\index{vektor!jumlah}\index{jumlah!vektor}%
\index{penjumlahan vektor}
dari \( \vec{u} \) dan \( \vec{v} \) adalah vektor yang komponennya
merupakan jumlah komponen-komponen yang bersesuaian.
\begin{equation*}
  \vec{u}+\vec{v}=
  \colvec{u_1 \\ \vdotswithin{u_1} \\ u_n}
   +
  \colvec{v_1 \\ \vdotswithin{v_1} \\ v_n}
   =
  \colvec{u_1+v_1 \\ \vdotswithin{u_1+v_1} \\ u_n+v_n}
\end{equation*}
%</df:VectorSum>
\end{definition}

Perhatikan bahwa agar penjumlahan terdefinisi, kedua vektor harus memiliki
jumlah entri yang sama.
Penjumlahan entri demi entri ini berlaku untuk sebarang pasangan matriks,
bukan hanya vektor, asalkan keduanya memiliki jumlah baris dan kolom yang sama.
\index{matriks!jumlah}\index{jumlah!matriks}

\begin{definition} \label{df:VectorScalarMultiplication}
%<*df:VectorScalarMultiplication>
\definend{Perkalian skalar\/}\index{vektor!kelipatan skalar}%
\index{kelipatan skalar!vektor} antara bilangan riil
\( r \) dan vektor \( \vec{v} \) adalah vektor yang setiap komponennya
merupakan hasil kali bilangan tersebut dengan komponen vektor yang bersesuaian.
\begin{equation*}
  r\cdot\vec{v}=
  r\cdot\colvec{v_1 \\ \vdotswithin{v_1} \\ v_n}
  =
  \colvec{rv_1 \\ \vdotswithin{rv_1} \\ rv_n}
\end{equation*}
%</df:VectorScalarMultiplication>
\end{definition}

Seperti operasi penjumlahan, operasi perkalian skalar entri demi entri
berlaku tidak hanya pada vektor, tetapi juga pada sebarang matriks.
\index{matriks!perkalian skalar}\index{perkalian skalar!matriks}

Perkalian skalar ditulis sebagai \( r\cdot\vec{v} \) atau
\( \vec{v}\cdot r \); kadang-kadang tanda `$\cdot$' bahkan dihilangkan:
~$r\vec{v}$.
(Jangan menyebut perkalian skalar sebagai `hasil kali skalar' karena nama
tersebut digunakan untuk operasi yang berbeda.)

\begin{example}
\begin{equation*}
  \colvec[r]{2 \\ 3 \\ 1}
   +
  \colvec[r]{3 \\ -1 \\ 4}
  =
  \colvec{2+3 \\ 3-1 \\ 1+4}
   =
  \colvec[r]{5 \\ 2 \\ 5}
  \qquad
  7\cdot\colvec[r]{1 \\ 4 \\ -1 \\ -3}
  =
  \colvec[r]{7 \\ 28 \\ -7 \\ -21}
\end{equation*}
\end{example}

Perhatikan bahwa definisi penjumlahan dan perkalian skalar selaras pada
kasus yang saling bertumpang tindih; misalnya,
\( \vec{v} +\vec{v} = 2\vec{v} \).

Dengan definisi-definisi ini, kita siap menggunakan notasi matriks dan vektor
untuk menyelesaikan sistem sekaligus menyatakan solusinya.

\begin{example} \label{ex:ManyParamsInfManySolsSystem}
Sistem
\begin{equation*}
   \begin{linsys}{5}
      2x  &+  &y  &  &  &-  &w  &   &   &=  &4  \\
          &   &y  &  &  &+  &w  &+  &u  &=  &4  \\
       x  &   &   &- &z &+  &2w &   &   &=  &0  
   \end{linsys}
\end{equation*}
direduksi sebagai berikut.
\begin{align*}
  \begin{amat}[r]{5}
    2  &1  &0  &-1  &0  &4  \\
    0  &1  &0  &1   &1  &4  \\
    1  &0  &-1 &2   &0  &0
  \end{amat}
  &\grstep{-(1/2)\rho_1+\rho_3}
  \begin{amat}[r]{5}
    2  &1     &0  &-1    &0  &4  \\
    0  &1     &0  &1     &1  &4  \\
    0  &-1/2  &-1 &5/2   &0  &-2
  \end{amat}                                 \\
  &\grstep{(1/2)\rho_2+\rho_3}
  \begin{amat}[r]{5}
    2  &1     &0  &-1    &0    &4  \\
    0  &1     &0  &1     &1    &4  \\
    0  &0     &-1 &3     &1/2  &0
  \end{amat}
\end{align*}
Himpunan solusinya adalah
\( \set{(w+(1/2)u,4-w-u,3w+(1/2)u,w,u)\suchthat w,u\in\Re} \).
Kita menuliskannya dalam bentuk vektor.
\begin{equation*}
  \set{\colvec{x \\ y \\ z \\ w \\ u}=
       \colvec[r]{0 \\ 4 \\ 0 \\ 0 \\ 0}+
       \colvec[r]{1 \\ -1 \\ 3 \\ 1 \\ 0}w+
       \colvec[r]{1/2 \\ -1 \\ 1/2 \\ 0 \\ 1}u
       \suchthat w,u\in\Re}
\end{equation*}
Perhatikan betapa jelas notasi vektor memisahkan koefisien setiap parameter.
Sebagai contoh, baris ketiga bentuk vektor memperlihatkan dengan jelas bahwa
jika \( u \) tetap, maka \( z \) bertambah tiga kali secepat \( w \).
Hal lain yang tampak jelas adalah bahwa menetapkan \( w \) dan \( u \)
sama dengan nol memberikan
\begin{equation*}
  \colvec{x \\ y \\ z \\ w \\ u}
  =\colvec[r]{0 \\ 4 \\ 0 \\ 0 \\ 0}
\end{equation*}
sebagai suatu solusi partikular sistem linear.
\end{example}

\begin{example}
Dengan cara yang sama, sistem
\begin{equation*}
   \begin{linsys}{3}
     x  &-  &y  &+  &z  &=  &1  \\
    3x  &   &   &+  &z  &=  &3  \\
    5x  &-  &2y &+  &3z &=  &5  
  \end{linsys}
\end{equation*}
direduksi
\begin{align*}
  \begin{amat}[r]{3}
    1  &-1  &1  &1  \\
    3  &0   &1  &3  \\
    5  &-2  &3  &5
  \end{amat}
  &\grstep[-5\rho_1+\rho_3]{-3\rho_1+\rho_2}
  \begin{amat}[r]{3}
    1  &-1  &1  &1  \\
    0  &3   &-2 &0  \\
    0  &3   &-2 &0
  \end{amat}                                    \\
  &\grstep{-\rho_2+\rho_3}
  \begin{amat}[r]{3}
    1  &-1  &1  &1  \\
    0  &3   &-2 &0  \\
    0  &0   &0  &0
  \end{amat}
\end{align*}
sehingga menghasilkan himpunan solusi dengan satu parameter.
\begin{equation*}
  \set{\colvec[r]{1 \\ 0 \\ 0}
       +\colvec[r]{-1/3 \\ 2/3 \\ 1}z
       \suchthat z\in\Re}
\end{equation*}
Seperti dalam contoh sebelumnya, vektor yang tidak dikaitkan dengan parameter
\begin{equation*}
   \colvec[r]{1 \\ 0 \\ 0}
\end{equation*}
merupakan solusi partikular sistem.
\end{example}

Sebelum mengerjakan latihan, kita akan meninjau apa yang telah dicapai dan
apa yang akan kita lakukan pada sisa bab ini.
Sejauh ini kita telah mempelajari mekanika Metode Gauss.
Kita belum berhenti untuk mempertimbangkan pertanyaan-pertanyaan yang muncul,
selain membuktikan \nearbytheorem{th:GaussMethod}\Dash yang membenarkan metode
tersebut dengan menunjukkan bahwa hasilnya benar.

Sebagai contoh, dapatkah kita selalu mendeskripsikan himpunan solusi seperti
di atas, dengan menambahkan suatu vektor solusi partikular pada kombinasi
linear beberapa vektor lain yang koefisiennya tidak dibatasi?
Kita telah mencatat bahwa himpunan solusi yang dideskripsikan dengan cara ini
memiliki tak hingga banyak anggota, sehingga menjawab pertanyaan tersebut akan
memberi tahu kita tentang ukuran himpunan solusi.
Subbagian berikut menunjukkan bahwa jawabannya adalah ``ya.''
Bagian kedua bab ini kemudian menggunakan jawaban tersebut untuk
mendeskripsikan geometri himpunan solusi.

Pertanyaan lain timbul dari pengamatan bahwa Metode Gauss dapat dilakukan
dengan lebih dari satu cara (misalnya, ketika menukar baris, mungkin ada
beberapa pilihan baris yang dapat dipertukarkan).
\nearbytheorem{th:GaussMethod} menyatakan bahwa kita harus memperoleh himpunan
solusi yang sama, apa pun langkah yang ditempuh. Namun, jika Metode Gauss
dilakukan dengan dua cara, haruskah jumlah variabel bebas dalam setiap sistem
berbentuk eselon sama?
Haruskah variabel bebasnya juga variabel yang sama? Dengan kata lain, apakah
mustahil menyelesaikan suatu soal dengan satu cara sehingga $y$ dan~$w$ bebas,
tetapi dengan cara lain sehingga $y$ dan~$z$ bebas?
Bagian ketiga bab ini menjawab ``ya'': dari sebarang sistem linear awal, semua
versi berbentuk eselon yang diturunkan memiliki variabel bebas yang sama.

Jadi, pada akhir bab kita tidak hanya memiliki landasan yang kokoh dalam
praktik Metode Gauss, tetapi juga dalam teorinya.
Kita akan mengetahui secara tepat apa yang dapat dan tidak dapat terjadi
dalam suatu reduksi.

\begin{exercises}
  \recommended \item  
    Tentukan entri matriks yang ditunjukkan, jika entri tersebut terdefinisi.
    \begin{equation*}
      A=\begin{mat}[r]
        1  &3  &1  \\
        2  &-1 &4
      \end{mat}
    \end{equation*}
    \begin{exparts*}
      \partsitem \( a_{2,1} \)
      \partsitem \( a_{1,2} \)
      \partsitem \( a_{2,2} \)
      \partsitem \( a_{3,1} \)
    \end{exparts*}
    \begin{answer}
      \begin{exparts*}
        \partsitem \( 2 \)
        \partsitem \( 3 \)
        \partsitem \(-1 \)
        \partsitem Tidak terdefinisi.
      \end{exparts*}  
    \end{answer}
  \recommended \item 
    Nyatakan ukuran setiap matriks.
    \begin{exparts*}
      \partsitem \(
        \begin{mat}[r]
          1  &0  &4  \\
          2  &1  &5
        \end{mat}  \)
      \partsitem \(
        \begin{mat}[r]
          1  &1  \\
         -1  &1  \\
          3  &-1
        \end{mat}  \)
      \partsitem \(
        \begin{mat}[r]
          5  &10 \\
         10  &5
        \end{mat}  \)
    \end{exparts*}
    \begin{answer}
      \begin{exparts*}
        \partsitem \( \nbym{2}{3} \)
        \partsitem \( \nbym{3}{2} \)
        \partsitem \( \nbym{2}{2} \)
      \end{exparts*}  
    \end{answer}
  \recommended \item 
    Lakukan operasi vektor yang ditunjukkan, jika operasi tersebut terdefinisi.
    \begin{exparts*}
      \partsitem \( \colvec[r]{2 \\ 1 \\ 1}
               +\colvec[r]{3 \\ 0 \\ 4} \)
      \partsitem \( 5\colvec[r]{4 \\ -1} \)
      \partsitem \( \colvec[r]{1 \\ 5 \\ 1}
               -\colvec[r]{3 \\ 1 \\ 1} \)
      \partsitem \( 7\colvec[r]{2 \\ 1}
               +9\colvec[r]{3 \\ 5} \)
      \partsitem \( \colvec[r]{1 \\ 2}
               +\colvec[r]{1 \\ 2 \\ 3} \)
      \partsitem \( 6\colvec[r]{3 \\ 1 \\ 1}
               -4\colvec[r]{2 \\ 0 \\ 3}
               +2\colvec[r]{1 \\ 1 \\ 5} \)
    \end{exparts*}
    \begin{answer}
      \begin{exparts*}
        \partsitem \( \colvec[r]{5 \\ 1 \\ 5} \)
        \partsitem \( \colvec[r]{20 \\ -5} \)
        \partsitem \( \colvec[r]{-2 \\ 4 \\ 0} \)
        \partsitem \( \colvec[r]{41 \\ 52} \)
        \partsitem Tidak terdefinisi.
        \partsitem \( \colvec[r]{12 \\ 8 \\ 4} \)
      \end{exparts*}  
     \end{answer}
  \recommended \item  \label{exer:SolveInMatrixNotation}
    Selesaikan setiap sistem dengan menggunakan notasi matriks.
    Nyatakan solusinya dengan menggunakan vektor.
    \begin{exparts*}
      \partsitem \( \begin{linsys}[t]{2}
                  3x  &+  &6y  &=  &18  \\
                   x  &+  &2y  &=  &6   
                   \end{linsys}  \)
      \partsitem \( \begin{linsys}[t]{2}
                   x  &+  &y   &=  &1  \\
                   x  &-  &y   &=  &-1   
                    \end{linsys}  \)
      \partsitem \( \begin{linsys}[t]{3}
                   x_1  &   &     &+  &x_3   &=  &4  \\
                   x_1  &-  &x_2  &+  &2x_3  &=  &5  \\
                  4x_1  &-  &x_2  &+  &5x_3  &=  &17  
                   \end{linsys}  \)
      \partsitem \( \begin{linsys}[t]{3}
                   2a   &+  &b    &-  &c     &=  &2  \\
                   2a   &   &     &+  &c     &=  &3  \\
                    a   &-  &b    &   &      &=  &0   
                    \end{linsys}  \)
      \partsitem \( \begin{linsys}[t]{4}
                     x  &+  &2y   &-   &z   &    &    &=  &3  \\
                    2x  &+  &y    &    &    &+   &w   &=  &4  \\
                     x  &-  &y    &+   &z   &+   &w   &=  &1  
                    \end{linsys}  \)
      \partsitem \( \begin{linsys}[t]{4}
                     x  &   &     &+   &z   &+   &w   &=  &4  \\
                    2x  &+  &y    &    &    &-   &w   &=  &2  \\
                    3x  &+  &y    &+   &z   &    &    &=  &7  
                     \end{linsys}  \)
    \end{exparts*}
    \begin{answer}
      \begin{exparts}
        \partsitem Reduksi ini
          \begin{equation*}
            \begin{amat}[r]{2}
              3  &6  &18 \\
              1  &2  &6
            \end{amat}
            \grstep{(-1/3)\rho_1+\rho_2}
            \begin{amat}[r]{2}
              3  &6  &18 \\
              0  &0  &0
            \end{amat}
          \end{equation*}
          menghasilkan \( x \) sebagai variabel utama dan \( y \) sebagai
          variabel bebas.
          Dengan memilih \( y \) sebagai parameter, diperoleh \( x=6-2y \)
          dan himpunan solusi berikut.
          \begin{equation*}
            \set{\colvec[r]{6 \\ 0}+\colvec[r]{-2 \\ 1}y
              \suchthat y\in\Re}
          \end{equation*}
        \partsitem Reduksi
          \begin{equation*}
            \begin{amat}[r]{2}
              1  &1  &1  \\
              1  &-1 &-1
            \end{amat}
            \grstep{-\rho_1+\rho_2}
            \begin{amat}[r]{2}
              1  &1  &1  \\
              0  &-2 &-2
            \end{amat}
          \end{equation*}
          menghasilkan solusi tunggal \( y=1 \), \( x=0 \).
          Himpunan solusinya adalah sebagai berikut.
          \begin{equation*}
            \set{\colvec[r]{0 \\ 1} }
          \end{equation*}
        \partsitem Metode Gauss
          \begin{equation*}
            \begin{amat}[r]{3}
              1  &0  &1  &4  \\
              1  &-1 &2  &5  \\
              4  &-1 &5  &17
            \end{amat}
            \grstep[-4\rho_1+\rho_3]{-\rho_1+\rho_2}
            \begin{amat}[r]{3}
              1  &0  &1  &4  \\
              0  &-1 &1  &1  \\
              0  &-1 &1  &1
            \end{amat}     
            \grstep{-\rho_2+\rho_3}
            \begin{amat}[r]{3}
              1  &0  &1  &4  \\
              0  &-1 &1  &1  \\
              0  &0  &0  &0
            \end{amat}
          \end{equation*}
          menyisakan \( x_1 \) dan \( x_2 \) sebagai variabel utama, dengan
          \( x_3 \) sebagai variabel bebas.
          Himpunan solusinya adalah sebagai berikut.
          \begin{equation*}
            \set{\colvec[r]{4 \\ -1 \\ 0}+\colvec[r]{-1 \\ 1 \\ 1}x_3
              \suchthat x_3\in\Re}
          \end{equation*}
        \partsitem Reduksi ini
          \begin{align*}
            \begin{amat}[r]{3}
              2  &1  &-1 &2  \\
              2  &0  &1  &3  \\
              1  &-1 &0  &0
            \end{amat}
            &\grstep[-(1/2)\rho_1+\rho_3]{-\rho_1+\rho_2}
            \begin{amat}[r]{3}
              2  &1    &-1   &2  \\
              0  &-1   &2    &1  \\
              0  &-3/2 &1/2  &-1
            \end{amat}                                          \\
            &\grstep{(-3/2)\rho_2+\rho_3}
            \begin{amat}[r]{3}
              2  &1  &-1   &2  \\
              0  &-1 &2    &1  \\
              0  &0  &-5/2 &-5/2
            \end{amat}
          \end{align*}
          menunjukkan bahwa himpunan solusinya beranggota tunggal.
          \begin{equation*}
            \set{\colvec[r]{1 \\ 1 \\ 1}}
          \end{equation*}
        \partsitem Reduksi berikut mudah dilakukan
          \begin{align*}
            \begin{amat}[r]{4}
              1  &2  &-1 &0  &3 \\
              2  &1  &0  &1  &4 \\
              1  &-1 &1  &1  &1
            \end{amat}
            &\grstep[-\rho_1+\rho_3]{-2\rho_1+\rho_2}
            \begin{amat}[r]{4}
              1  &2  &-1 &0  &3  \\
              0  &-3 &2  &1  &-2 \\
              0  &-3 &2  &1  &-2
            \end{amat}                                       \\
            &\grstep{-\rho_2+\rho_3}
            \begin{amat}[r]{4}
              1  &2  &-1 &0  &3  \\
              0  &-3 &2  &1  &-2 \\
              0  &0  &0  &0  &0
            \end{amat}
          \end{align*}
          dan berakhir dengan \( x \) dan $y$ sebagai variabel utama, sedangkan
          \( z \) dan \( w \) merupakan variabel bebas.
          Menyelesaikan persamaan terhadap \( y \) menghasilkan
          \( y=(2+2z+w)/3 \), dan substitusi menunjukkan bahwa
          \( x+2(2+2z+w)/3-z=3 \), sehingga
          \( x=(5/3)-(1/3)z-(2/3)w \) dan himpunan solusinya adalah sebagai
          berikut.
          \begin{equation*}
            \set{\colvec[r]{5/3 \\ 2/3 \\ 0 \\ 0}
                 +\colvec[r]{-1/3 \\ 2/3 \\ 1 \\ 0}z
                 +\colvec[r]{-2/3 \\ 1/3 \\ 0 \\ 1}w
                 \suchthat z,w\in\Re}
          \end{equation*}
        \partsitem Reduksi
          \begin{align*}
            \begin{amat}[r]{4}
              1  &0  &1  &1  &4 \\
              2  &1  &0  &-1 &2 \\
              3  &1  &1  &0  &7
            \end{amat}
            &\grstep[-3\rho_1+\rho_3]{-2\rho_1+\rho_2}
            \begin{amat}[r]{4}
              1  &0  &1  &1  &4 \\
              0  &1  &-2 &-3 &-6\\
              0  &1  &-2 &-3 &-5
            \end{amat}                                       \\
            &\grstep{-\rho_2+\rho_3}
            \begin{amat}[r]{4}
              1  &0  &1  &1  &4 \\
              0  &1  &-2 &-3 &-6\\
              0  &0  &0  &0  &1
            \end{amat}
          \end{align*}
          menunjukkan bahwa tidak ada solusi\Dash himpunan solusinya kosong.
      \end{exparts}  
     \end{answer}
  \item \label{exer:SlvMatNot}
    Selesaikan setiap sistem dengan menggunakan notasi matriks.
    Nyatakan setiap himpunan solusi dalam notasi vektor.
    \begin{exparts*}
      \partsitem \( \begin{linsys}[t]{3}
                  2x  &+  &y  &-  &z  &=  &1  \\
                  4x  &-  &y  &   &   &=  &3  
                \end{linsys}  \)
      \partsitem \( \begin{linsys}[t]{4}
                   x  &   &   &-  &z  &   &   &=  &1  \\
                      &   &y  &+  &2z &-  &w  &=  &3  \\
                   x  &+  &2y &+  &3z &-  &w  &=  &7  
               \end{linsys}  \)
      \partsitem \( \begin{linsys}[t]{4}
                   x  &-  &y  &+  &z  &   &   &=  &0  \\
                      &   &y  &   &   &+  &w  &=  &0  \\
                  3x  &-  &2y &+  &3z &+  &w  &=  &0  \\
                      &   &-y &   &   &-  &w  &=  &0  
               \end{linsys}  \)
      \partsitem \( \begin{linsys}[t]{5}
                   a  &+  &2b &+  &3c &+  &d  &-  &e  &=  &1  \\
                  3a  &-  &b  &+  &c  &+  &d  &+  &e  &=  &3  
               \end{linsys}  \)
    \end{exparts*}
    \begin{answer}
      \begin{exparts}
      \partsitem Reduksi
        \begin{equation*}
          \begin{amat}[r]{3}
            2  &1  &-1  &1  \\
            4  &-1 &0   &3
          \end{amat}
          \grstep{-2\rho_1+\rho_2}
          \begin{amat}[r]{3}
            2  &1  &-1  &1  \\
            0  &-3 &2   &1
          \end{amat}
        \end{equation*}
        berakhir dengan \( x \) dan \( y \) sebagai variabel utama, serta
        \( z \) sebagai variabel bebas.
        Menyelesaikan persamaan terhadap \( y \) menghasilkan
        \( y=(1-2z)/(-3) \), kemudian substitusi
        \( 2x+(1-2z)/(-3)-z=1 \) menunjukkan bahwa
        \( x=((4/3)+(1/3)z)/2 \).
        Dengan demikian, himpunan solusinya adalah sebagai berikut.
        \begin{equation*}
          \set{\colvec[r]{2/3 \\ -1/3 \\ 0}
               +\colvec[r]{1/6 \\ 2/3 \\ 1}z
              \suchthat z\in\Re}
        \end{equation*}
      \partsitem Penerapan Metode Gauss berikut
        \begin{align*}
          \begin{amat}[r]{4}
            1  &0  &-1  &0  &1 \\
            0  &1  &2   &-1 &3 \\
            1  &2  &3   &-1 &7
          \end{amat}
          &\grstep{-\rho_1+\rho_3}
          \begin{amat}[r]{4}
            1  &0  &-1  &0  &1 \\
            0  &1  &2   &-1 &3 \\
            0  &2  &4   &-1 &6
          \end{amat}                                   \\
          &\grstep{-2\rho_2+\rho_3}
          \begin{amat}[r]{4}
            1  &0  &-1  &0  &1 \\
            0  &1  &2   &-1 &3 \\
            0  &0  &0   &1  &0
          \end{amat}
        \end{align*}
        menyisakan \( x \), \( y \), dan \( w \) sebagai variabel utama.
        Himpunan solusinya adalah sebagai berikut.
        \begin{equation*}
          \set{\colvec[r]{1 \\ 3 \\ 0 \\ 0}
               +\colvec[r]{1 \\ -2 \\ 1 \\ 0}z
              \suchthat z\in\Re}
        \end{equation*}
      \partsitem Reduksi baris berikut
        \begin{align*}
          \begin{amat}[r]{4}
            1  &-1 &1   &0  &0 \\
            0  &1  &0   &1  &0 \\
            3  &-2 &3   &1  &0 \\
            0  &-1 &0   &-1 &0
          \end{amat}
          &\grstep{-3\rho_1+\rho_3}
          \begin{amat}[r]{4}
            1  &-1 &1   &0  &0 \\
            0  &1  &0   &1  &0 \\
            0  &1  &0   &1  &0 \\
            0  &-1 &0   &-1 &0
          \end{amat}                                       \\
          &\grstep[\rho_2+\rho_4]{-\rho_2+\rho_3}
          \begin{amat}[r]{4}
            1  &-1 &1   &0  &0 \\
            0  &1  &0   &1  &0 \\
            0  &0  &0   &0  &0 \\
            0  &0  &0   &0  &0
          \end{amat}
        \end{align*}
        berakhir dengan \( z \) dan \( w \) sebagai variabel bebas.
        Kita memperoleh himpunan solusi berikut.
        \begin{equation*}
          \set{\colvec[r]{0 \\ 0 \\ 0 \\ 0}
               +\colvec[r]{-1 \\ 0 \\ 1 \\ 0}z
               +\colvec[r]{-1 \\ -1 \\ 0 \\ 1}w
              \suchthat z,w\in\Re}
        \end{equation*}
      \partsitem Metode Gauss yang dilakukan dengan cara berikut
        \begin{equation*}
          \begin{amat}[r]{5}
            1  &2  &3   &1  &-1 &1  \\
            3  &-1 &1   &1  &1  &3
          \end{amat}
          \grstep{-3\rho_1+\rho_2}
          \begin{amat}[r]{5}
            1  &2  &3   &1  &-1 &1  \\
            0  &-7 &-8  &-2 &4  &0
          \end{amat}
        \end{equation*}
        berakhir dengan \( c \), \( d \), dan \( e \) sebagai variabel bebas.
        Menyelesaikan persamaan terhadap \( b \) menunjukkan bahwa
        \( b=(8c+2d-4e)/(-7) \), kemudian substitusi
        \( a+2(8c+2d-4e)/(-7)+3c+1d-1e=1 \) menunjukkan bahwa
        \( a=1-(5/7)c-(3/7)d-(1/7)e \), sehingga kita memperoleh himpunan
        solusi berikut.
        \begin{equation*}
          \set{\colvec[r]{1 \\ 0 \\ 0 \\ 0 \\ 0}
               +\colvec[r]{-5/7 \\ -8/7 \\ 1 \\ 0 \\ 0}c
               +\colvec[r]{-3/7 \\ -2/7 \\ 0 \\ 1 \\ 0}d
               +\colvec[r]{-1/7 \\ 4/7 \\ 0 \\ 0 \\ 1}e
              \suchthat c,d,e\in\Re}
        \end{equation*}
    \end{exparts}  
   \end{answer}
  \item 
    Selesaikan setiap sistem dengan menggunakan notasi matriks.
    Nyatakan himpunan solusinya dengan menggunakan vektor.
    \begin{exparts*}
      \partsitem 
        $\begin{linsys}{3}
          3x &+ &2y &+ &z &= &1 \\
          x  &- &y  &+ &z &= &2 \\
          5x &+ &5y &+ &z &= &0
        \end{linsys}$
      \partsitem
        $\begin{linsys}{4}
          x  &+ &y  &- &2z &= &0 \\
          x  &- &y  &  &   &= &-3 \\
          3x &- &y  &- &2z &= &-6  \\
             &  &2y &- &2z &= &3  
        \end{linsys}$
      \partsitem
        $\begin{linsys}{5}
          2x  &- &y  &- &z &+ &w &= &4 \\
           x  &+ &y  &+ &z &  &  &= &-1 
        \end{linsys}$
      \partsitem
         $\begin{linsys}{3}
          x  &+ &y  &- &2z &= &0 \\
          x  &- &y  &  &   &= &-3 \\
          3x &- &y  &- &2z &= &0    
        \end{linsys}$ 
    \end{exparts*}
    \begin{answer}
      \begin{exparts}
        \partsitem Reduksi ini
          \begin{align*}
            \begin{amat}{3}
              3 &2  &1  &1 \\
              1 &-1 &1  &2 \\
              5 &5  &1  &0
            \end{amat}
            &\grstep[-(5/3)\rho_1+\rho_3]{-(1/3)\rho_1+\rho_2}
            \begin{amat}{3}
              3 &2    &1     &1 \\
              0 &-5/3 &2/3   &5/3 \\
              0 &5/3  &-2/3  &-5/3
            \end{amat}                                     \\
            &\grstep{\rho_2+\rho_3}
            \begin{amat}{3}
              3 &2    &1     &1 \\
              0 &-5/3 &2/3   &5/3 \\
              0 &0    &0     &0
            \end{amat}
          \end{align*}
          menghasilkan himpunan solusi berikut.
          \begin{equation*}
            \set{\colvec{x \\ y \\ z}=\colvec{1 \\ -1 \\ 0}
                                       +\colvec{-3/5 \\ 2/5 \\ 1}z
                             \suchthat z\in\Re}
          \end{equation*}
        \partsitem
          Reduksinya adalah sebagai berikut.
          \begin{align*}
            \begin{amat}{3}
              1 &1   &-2 &0 \\
              1 &-1  &0  &-3 \\
              3 &-1  &-2 &-6 \\
              0 &2   &-2 &3
            \end{amat}
            &\grstep[-3\rho_1+\rho_3]{-\rho_1+\rho_2}
            \begin{amat}{3}
              1 &1   &-2 &0 \\
              0 &-2  &2  &-3  \\
              0 &-4  &4  &-6 \\
              0 &2   &-2 &3
            \end{amat}                                  \\
            &\grstep[\rho_2+\rho_4]{-2\rho_2+\rho_3}
            \begin{amat}{3}
              1 &1   &-2 &0 \\
              0 &-2  &2  &-3  \\
              0 &0   &0  &0 \\
              0 &0   &0  &0
            \end{amat}
          \end{align*}
          Himpunan solusinya adalah sebagai berikut.
          \begin{equation*}
            \set{\colvec{-3/2 \\ 3/2 \\ 0}
                 +\colvec{1 \\ 1 \\ 1}z
                 \suchthat z\in\Re}
          \end{equation*}
      \partsitem
         Metode Gauss
         \begin{equation*}
           \begin{amat}{4}
             2 &-1 &-1 &1 &4 \\
             1 &1  &1  &0 &-1
           \end{amat}
           \grstep{-(1/2)\rho_1+\rho_2}
           \begin{amat}{4}
             2 &-1   &-1   &1    &4 \\
             0 &3/2  &3/2  &-1/2 &-3
           \end{amat}
         \end{equation*}
         menghasilkan himpunan solusi berikut.
         \begin{equation*}
           \set{\colvec{1 \\ -2 \\ 0 \\ 0}
                  +\colvec{0 \\ -1 \\ 1 \\ 0}z
                  +\colvec{-1/3 \\ 1/3 \\ 0 \\ 1}w
                 \suchthat z,w\in\Re}
         \end{equation*}
       \partsitem
          Reduksinya adalah sebagai berikut.
          \begin{equation*}
            \begin{amat}{3}
              1 &1  &-2 &0  \\
              1 &-1 &0  &-3 \\
              3 &-1 &-2 &0
           \end{amat}
           \grstep[-3\rho_1+\rho_3]{-\rho_1+\rho_2}
           \begin{amat}{3}
             1 &1  &-2 &0  \\
             0 &-2 &2  &-3 \\
             0 &-4  &4  &0
          \end{amat}
          \grstep{-2\rho_2+\rho_3}
          \begin{amat}{3}
            1 &1  &-2 &0  \\
            0 &-2 &2  &-3 \\
            0 &0  &0  &6
          \end{amat}
        \end{equation*}
        Himpunan solusinya kosong, yaitu~$\set{}$.
      \end{exparts}
    \end{answer}
  \recommended \item 
    Vektor tersebut berada dalam himpunan.
    Nilai parameter apa yang menghasilkan vektor tersebut?
    \begin{exparts}
      \partsitem $\colvec[r]{5 \\ -5}$,
        $\set{\colvec[r]{1 \\ -1}k\suchthat k\in\Re}$
      \partsitem $\colvec[r]{-1 \\ 2 \\ 1}$,
        $\set{\colvec[r]{-2 \\ 1 \\ 0}i
           +\colvec[r]{3 \\ 0 \\ 1}j\suchthat i,j\in\Re}$
      \partsitem $\colvec[r]{0 \\ -4 \\ 2}$,
        $\set{\colvec[r]{1 \\ 1 \\ 0}m
               +\colvec[r]{2 \\ 0 \\ 1}n\suchthat m,n\in\Re}$
    \end{exparts}
    \begin{answer}
      Untuk setiap soal, kita memperoleh sistem persamaan linear dengan
      menyamakan komponen-komponen yang bersesuaian.
      \begin{exparts}
       \partsitem $k=5$
       \partsitem Komponen kedua menunjukkan bahwa $i=2$, sedangkan komponen
       ketiga menunjukkan bahwa $j=1$.
       \partsitem $m=-4$, $n=2$
      \end{exparts} 
    \end{answer}
  \item 
    Tentukan apakah vektor tersebut berada dalam himpunan.
    \begin{exparts}
      \partsitem $\colvec[r]{3 \\ -1}$,
        $\set{\colvec[r]{-6 \\ 2}k\suchthat k\in\Re}$
      \partsitem $\colvec[r]{5 \\ 4}$,
        $\set{\colvec[r]{5 \\ -4}j\suchthat j\in\Re}$
      \partsitem $\colvec[r]{2 \\ 1 \\ -1}$,
        $\set{\colvec[r]{0 \\ 3 \\ -7}
             +\colvec[r]{1 \\ -1 \\ 3}r\suchthat r\in\Re}$
      \partsitem $\colvec[r]{1 \\ 0 \\ 1}$,
        $\set{\colvec[r]{2 \\ 0 \\ 1}j
            +\colvec[r]{-3 \\ -1 \\ 1}k\suchthat j,k\in\Re}$
    \end{exparts}
    \begin{answer}
      Untuk setiap soal, kita memperoleh sistem persamaan linear dengan
      menyamakan komponen-komponen yang bersesuaian.
      \begin{exparts}
        \partsitem Ya; ambil $k=-1/2$.
        \partsitem Tidak; sistem dengan persamaan $5=5\cdot j$ dan
            $4=-4\cdot j$ tidak memiliki solusi.
        \partsitem Ya; ambil $r=2$.
        \partsitem Tidak.
           Komponen kedua menghasilkan $k=0$.
           Kemudian komponen ketiga menghasilkan $j=1$.
           Namun komponen pertama tidak terpenuhi.
      \end{exparts}
     \end{answer}
  \item \cite{Cleary}
    Seorang petani dengan lahan seluas 1200 acre mempertimbangkan untuk menanam
    tiga jenis tanaman: jagung, kedelai, dan oat.
    Petani tersebut ingin menggunakan seluruh~$1200$ acre lahannya.
    Benih jagung berharga \$$20$ per acre, sedangkan benih kedelai dan oat
    masing-masing berharga \$$50$ dan~\$$12$ per acre.
    Petani tersebut memiliki \$$40\,000$ untuk membeli benih dan bermaksud
    menghabiskan seluruhnya.
    \begin{exparts}  
      \item Gunakan informasi di atas untuk merumuskan dua persamaan linear
        dengan tiga variabel tak diketahui, lalu selesaikan sistem tersebut.
     \item Solusi sistem merepresentasikan pilihan yang dapat dibuat petani.
        Tuliskan dua solusi yang masuk akal.
     \item Misalkan ketika tanaman matang pada musim gugur, petani dapat
        memperoleh pendapatan sebesar \$$100$ per acre untuk jagung,
        \$$300$ per acre untuk kedelai, dan \$$80$ per acre untuk oat.
        Manakah dari kedua solusi pada bagian sebelumnya yang akan menghasilkan
        pendapatan lebih besar?
    \end{exparts}
    \begin{answer}
      \begin{exparts}
        \item Misalkan $c$ adalah banyaknya acre untuk jagung, $s$ banyaknya
          acre untuk kedelai, dan $a$ banyaknya acre untuk oat.
          \begin{equation*}
            \begin{linsys}{3}
              c   &+   &s   &+   &a   &=   &1200 \\ 
            20c   &+   &50s &+   &12a &=   &40\,000  
            \end{linsys}
            \grstep{-20\rho_1+\rho_2}
            \begin{linsys}{3}
              c   &+   &s   &+   &a   &=   &1200 \\ 
                  &    &30s &-   &8a  &=   &16\,000  
            \end{linsys}
          \end{equation*}
          Untuk mendeskripsikan himpunan solusi, kita dapat memparameterkannya
          dengan menggunakan $a$.
          \begin{equation*}
            \set{\colvec{c \\ s \\ a}
                 =\colvec{20\,000/30 \\ 16\,000/30 \\ 0}
                  +\colvec{-38/30 \\ 8/30 \\ 1}a
                 \suchthat a\in\Re}
          \end{equation*}
        \item Ada banyak jawaban yang mungkin.
          Sebagai contoh, kita dapat mengambil $a=0$ sehingga
          $c=20\,000/30\approx 666.66$ dan $s=16000/30\approx 533.33$.
          Contoh lain adalah mengambil $a=20\,000/38\approx 526.32$, sehingga
          $c=0$ dan $s=25\,600/38\approx 673.68$.
        \item Substitusi jawaban dari bagian sebelumnya ke dalam
          $100c+300s+80a$ memberikan pendapatan
          $680\,000/3\approx\$226\,666.67$ untuk rencana pertama dan
          $4\,640\,000/19\approx\$244\,210.53$ untuk rencana kedua.
          Jadi, rencana kedua menghasilkan pendapatan yang lebih besar.
      \end{exparts}
    \end{answer}
  \item 
    Parameterkan himpunan solusi sistem dengan satu persamaan berikut.
    \begin{equation*}
      x_1+x_2+\cdots+x_n=0
    \end{equation*}
    \begin{answer}
      Sistem ini memiliki satu persamaan.
      Variabel utamanya adalah \( x_1 \); variabel lainnya bebas.
      \begin{equation*}
        \set{\colvec{-1 \\ 1 \\ \vdotswithin{-1} \\ 0}x_2
             +\cdots+
             \colvec{-1 \\ 0 \\ \vdotswithin{-1} \\ 1}x_n
             \suchthat x_2,\ldots,x_n\in\Re}
      \end{equation*}  
     \end{answer}
  \recommended \item 
    \begin{exparts}
    \partsitem Terapkan Metode Gauss pada ruas kiri untuk menyelesaikan
      \begin{equation*}
        \begin{linsys}{4}
          x  &+  &2y  &    &    &-   &w   &=   &a   \\
         2x  &   &    &+   &z   &    &    &=   &b   \\
          x  &+  &y   &    &    &+   &2w  &=   &c   
        \end{linsys}
      \end{equation*}
      terhadap \( x \), \( y \), \( z \), dan \(  w \), dalam bentuk
      konstanta $a$, $b$, dan $c$.
    \partsitem Gunakan jawaban dari bagian sebelumnya untuk menyelesaikan
      sistem berikut.
      \begin{equation*}
        \begin{linsys}{4}
          x  &+  &2y  &    &    &-   &w   &=   &3   \\
         2x  &   &    &+   &z   &    &    &=   &1   \\
          x  &+  &y   &    &    &+   &2w  &=   &-2
        \end{linsys}
      \end{equation*}
    \end{exparts}
    \begin{answer}
      \begin{exparts}
        \partsitem Metode Gauss di sini menghasilkan
          \begin{align*}
            \begin{amat}{4}
              1  &2  &0  &-1  &a  \\
              2  &0  &1  &0   &b  \\
              1  &1  &0  &2   &c
            \end{amat}
            &\grstep[-\rho_1+\rho_3]{-2\rho_1+\rho_2}
            \begin{amat}{4}
              1  &2  &0  &-1  &a  \\
              0  &-4 &1  &2   &-2a+b  \\
              0  &-1 &0  &3   &-a+c
            \end{amat}                                  \\
            &\grstep{-(1/4)\rho_2+\rho_3}
            \begin{amat}{4}
              1  &2  &0    &-1  &a  \\
              0  &-4 &1    &2   &-2a+b  \\
              0  &0  &-1/4 &5/2 &-(1/2)a-(1/4)b+c
            \end{amat}
          \end{align*}
          sehingga \( w \) tetap bebas.
          Dari sini, \(  z=2a+b-4c+10w \), dan
          \( -4y=-2a+b-(2a+b-4c+10w)-2w \), sehingga
          \( y=a-c+3w \), serta
          \( x=a-2(a-c+3w)+w=-a+2c-5w. \)
          Oleh karena itu, himpunan solusinya adalah sebagai berikut.
          \begin{equation*}
             \set{\colvec{-a+2c \\ a-c \\ 2a+b-4c \\ 0}
                  +\colvec{-5 \\ 3 \\ 10 \\ 1}w
                  \suchthat w\in\Re}
          \end{equation*}
        \partsitem Substitusikan \( a=3 \), \( b=1 \), dan \( c=-2 \).
          \begin{equation*}
             \set{\colvec[r]{-7 \\ 5 \\ 15 \\ 0}
                  +\colvec[r]{-5 \\ 3 \\ 10 \\ 1}w
                  \suchthat w\in\Re}
          \end{equation*}
      \end{exparts}  
     \end{answer}
  \item 
    Mengapa koma diperlukan dalam notasi `\( a_{i,j} \)'
    untuk entri matriks?
    \begin{answer}
       Jika koma dihilangkan, misalnya dengan menulis \( a_{123} \),
       notasinya menjadi ambigu karena dapat berarti $a_{1,23}$ atau $a_{12,3}$.
    \end{answer}
  \recommended \item 
    Berikan matriks \( \nbyn{4} \) yang entri pada posisi
    \( i,j \) adalah
    \begin{exparts*}
      \partsitem \( i+j \);
      \partsitem \( -1 \) dipangkatkan dengan \( i+j \).
    \end{exparts*}
    \begin{answer}
      \begin{exparts*}
        \partsitem \(
           \begin{mat}[r]
             2  &3  &4  &5  \\
             3  &4  &5  &6  \\
             4  &5  &6  &7  \\
             5  &6  &7  &8
           \end{mat} \)
        \partsitem \(
           \begin{mat}[r]
             1  &-1  &1   &-1  \\
            -1  &1   &-1  &1  \\
             1  &-1  &1   &-1  \\
            -1  &1   &-1  &1
           \end{mat} \)
      \end{exparts*}  
    \end{answer}
  \item  
    Untuk setiap matriks \( A \), \definend{transpos}\index{transpos}
    \index{matriks!transpos}
    dari \( A \), yang ditulis
    \( \trans{A} \), adalah matriks yang kolom-kolomnya merupakan baris-baris
    \( A \).
    Tentukan transpos setiap matriks atau vektor berikut.
    \begin{exparts*}
      \partsitem \( \begin{mat}[r]
                  1  &2  &3  \\
                  4  &5  &6
               \end{mat}  \)
      \partsitem \( \begin{mat}[r]
                  2  &-3 \\
                  1  &1
               \end{mat}  \)
      \partsitem \( \begin{mat}[r]
                  5  &10 \\
                 10  &5
               \end{mat}  \)
      \partsitem \( \colvec[r]{1 \\ 1 \\ 0} \)
    \end{exparts*}
    \begin{answer}
      \begin{exparts*}
        \partsitem \( \begin{mat}[r]
                   1  &4  \\
                   2  &5  \\
                   3  &6
                 \end{mat}  \)
        \partsitem \( \begin{mat}[r]
                   2  &1  \\
                  -3  &1
                 \end{mat}  \)
        \partsitem \( \begin{mat}[r]
                   5  &10 \\
                  10  &5
                 \end{mat}  \)
        \partsitem \( \rowvec{1 &1 &0}  \)
      \end{exparts*}  
     \end{answer}
  \recommended \item 
    \begin{exparts}
      \partsitem Deskripsikan semua fungsi \( f(x)=ax^2+bx+c \)
        sedemikian sehingga \( f(1)=2 \) dan \( f(-1)=6 \).
      \partsitem Deskripsikan semua fungsi \( f(x)=ax^2+bx+c \)
        sedemikian sehingga \( f(1)=2 \).
    \end{exparts}
    \begin{answer}
      \begin{exparts}
        \partsitem Menyubstitusikan \( x=1 \) dan \( x=-1 \) menghasilkan
          \begin{equation*}
            \begin{linsys}{3}
              a  &+  &b   &+  &c  &=  &2  \\
              a  &-  &b   &+  &c  &=  &6  
            \end{linsys}
            \grstep{-\rho_1+\rho_2}
            \begin{linsys}{3}
              a  &+  &b   &+  &c  &=  &2  \\
                 &   &-2b &   &   &=  &4  
              \end{linsys}
          \end{equation*}
          sehingga himpunan fungsinya adalah
          \( \set{f(x)=(4-c)x^2-2x+c\suchthat c\in\Re} \).
        \partsitem Menyubstitusikan \( x=1 \) menghasilkan
          \begin{equation*}
            \begin{linsys}{3}
              a  &+  &b   &+  &c  &=  &2  
            \end{linsys}
          \end{equation*}
          sehingga himpunan fungsinya adalah
          \( \set{f(x)=(2-b-c)x^2+bx+c\suchthat b,c\in\Re} \).
      \end{exparts}  
    \end{answer}
  \item Tunjukkan bahwa sebarang himpunan lima titik pada bidang \( \Re^2 \)
    terletak pada satu irisan kerucut yang sama; artinya, semuanya memenuhi
    suatu persamaan berbentuk \( ax^2+by^2+cxy+dx+ey+f=0 \), dengan setidaknya
    satu di antara \( a,\,\ldots\,,f \) tidak nol.
    \begin{answer}
      Dengan menyubstitusikan kelima pasangan $(x,y)$, kita memperoleh suatu
      sistem dengan lima persamaan dan enam variabel tak diketahui
      $a$, \ldots, $f$.
      Karena variabel tak diketahui lebih banyak daripada persamaan, jika tidak
      ada ketakkonsistenan di antara persamaan-persamaan tersebut, maka terdapat
      tak hingga banyak solusi (setidaknya satu variabel akan tetap bebas).

      Sistem itu pasti konsisten karena $a=0$, \ldots, $f=0$ merupakan suatu
      solusi (kita hanya menggunakan solusi nol ini untuk menunjukkan bahwa
      sistem tersebut konsisten\Dash paragraf sebelumnya menunjukkan bahwa
      terdapat solusi tak nol).
    \end{answer}
  \item 
    Susunlah sistem dengan empat persamaan dan empat variabel tak diketahui
    yang memiliki
    \begin{exparts}
      \partsitem himpunan solusi dengan satu parameter;
      \partsitem himpunan solusi dengan dua parameter;
      \partsitem himpunan solusi dengan tiga parameter.
    \end{exparts}
    \begin{answer}
      \begin{exparts}
      \partsitem Berikut satu contoh\Dash persamaan keempat redundan, tetapi
        tetap dapat digunakan.
        \begin{equation*}
          \begin{linsys}{4}
             x  &+  &y  &-  &z  &+  &w  &=  &0  \\
                &   &y  &-  &z  &   &   &=  &0  \\
                &   &   &   &2z &+  &2w &=  &0  \\
                &   &   &   &z  &+  &w  &=  &0
          \end{linsys}
        \end{equation*}
      \partsitem Berikut satu contoh.
        \begin{equation*}
          \begin{linsys}{4}
             x  &+  &y  &-  &z  &+  &w  &=  &0  \\
                &   &   &   &   &   &w  &=  &0  \\
                &   &   &   &   &   &w  &=  &0  \\
                &   &   &   &   &   &w  &=  &0
          \end{linsys}
        \end{equation*}
      \partsitem Berikut satu contoh.
        \begin{equation*}
          \begin{linsys}{4}
             x  &+  &y  &-  &z  &+  &w  &=  &0  \\
             x  &+  &y  &-  &z  &+  &w  &=  &0  \\
             x  &+  &y  &-  &z  &+  &w  &=  &0  \\
             x  &+  &y  &-  &z  &+  &w  &=  &0 
          \end{linsys}
        \end{equation*}
    \end{exparts}  
   \end{answer}
  \puzzle \item 
    \cite{Shepelev}
    Teka-teki ini berasal dari situs web Rusia
    \texttt{http://www.arbuz.uz/}. Ada banyak cara untuk menyelesaikannya,
    tetapi penyelesaian saya menggunakan aljabar linear dan sangat naif.
    Ada sebuah planet yang dihuni oleh para arbuzoid
    (kaum semangka, jika diterjemahkan dari bahasa Rusia).
    Makhluk-makhluk itu memiliki tiga warna: merah, hijau, dan biru.
    Terdapat $13$~arbuzoid merah, $15$~arbuzoid hijau, dan $17$~arbuzoid biru.
    Ketika dua arbuzoid yang warnanya berbeda bertemu, keduanya berubah menjadi
    warna ketiga.

    Pertanyaannya: mungkinkah mereka semua pada akhirnya memiliki warna yang
    sama?
    \begin{answer}
       \answerasgiven
       Penyelesaian saya adalah merepresentasikan banyaknya arbuzoid sebagai
       vektor berdimensi~$3$, dan merepresentasikan semua transisi elementer
       yang mungkin sebagai vektor pula:
       \begin{center}
         \begin{tabular}{rr}
           M: &$13$  \\
           H: &$15$  \\
           B: &$17$
         \end{tabular}
         \qquad
         Operasi:
         $\colvec[r]{-1 \\ -1 \\ 2}$, 
         $\colvec[r]{-1 \\ 2 \\ -1}$, 
         dan
         $\colvec[r]{2 \\ -1 \\ -1}$
       \end{center}
       Sekarang cukup diperiksa apakah salah satu sistem persamaan linear
       berikut memiliki solusi:
       \begin{equation*}
         \colvec[r]{13 \\ 15 \\ 17}
         +x\colvec[r]{-1 \\ -1  \\ 2}
         +y\colvec[r]{-1 \\ 2 \\ -1}
         +z\colvec[r]{2 \\ -1 \\ -1}
         =\colvec[r]{0 \\ 0 \\ 45}
         \qquad
         \text{(atau $\colvec[r]{0 \\ 45 \\ 0}$ atau $\colvec[r]{45 \\ 0 \\ 0}$)}
       \end{equation*}
       Penyelesaian
       \begin{equation*}
         \begin{amat}[r]{3}
          -1  &-1 &2  &-13  \\
          -1  &2  &-1 &-15  \\
           2  &-1 &-1 &28
         \end{amat}
         \grstep[2\rho_1+\rho_3]{-\rho_1+\rho_2}
         \repeatedgrstep{\rho_2+\rho_3}
         \begin{amat}[r]{3}
          -1  &-1 &2  &-13  \\
           0  &3  &-3 &-2  \\
           0  &0  &0  &0
         \end{amat}
       \end{equation*}
       menghasilkan $y+2/3=z$, sehingga jika banyaknya transformasi $z$
       merupakan bilangan bulat, maka $y$ bukan bilangan bulat.
       Kedua sistem lainnya menghasilkan kesimpulan serupa, sehingga tidak ada
       solusi.
    \end{answer}
  \puzzle \item 
    \cite{USSROlympiad174}
    \begin{exparts}
      \partsitem Selesaikan sistem persamaan berikut.
        \begin{equation*}
          \begin{linsys}{2}
            ax  &+  &y  &=  &a^2  \\
             x  &+  &ay &=  &1
         \end{linsys}
        \end{equation*}
        Untuk nilai $a$ yang mana sistem tersebut tidak memiliki solusi, dan
        untuk nilai $a$ yang mana sistem tersebut memiliki tak hingga banyak
        solusi?
      \partsitem Jawab pertanyaan di atas untuk sistem berikut.
        \begin{equation*}
          \begin{linsys}{2}
            ax  &+  &y  &=  &a^3  \\    
             x  &+  &ay &=  &1
          \end{linsys}
        \end{equation*}
    \end{exparts}
    \begin{answer}
       \answerasgiven
       \begin{exparts}
        \partsitem Penyelesaian formal sistem menghasilkan
          \begin{equation*}
            x=\frac{a^3-1}{a^2-1}  
            \qquad
            y=\frac{-a^2+a}{a^2-1}.
          \end{equation*}
          Jika $a+1\neq 0$ dan $a-1\neq 0$, maka sistem memiliki solusi tunggal
          \begin{equation*}
            x=\frac{a^2+a+1}{a+1}
            \qquad
            y=\frac{-a}{a+1}.
          \end{equation*}
          Jika $a=-1$ atau $a=+1$, rumus-rumus tersebut tidak terdefinisi.
          Untuk kasus pertama, kita memperoleh sistem
          \begin{equation*}
            \left\{ 
            \begin{linsys}{2}
              -x &+  &y  &=  &1 \\
               x &-  &y  &=  &1
            \end{linsys}\right.
          \end{equation*}
          yang tak konsisten.
          Untuk kasus kedua, kita memperoleh
          \begin{equation*}
            \left\{
            \begin{linsys}{2}
               x &+  &y  &=  &1 \\
               x &+  &y  &=  &1
            \end{linsys}\right.
          \end{equation*}
          yang memiliki tak hingga banyak solusi (misalnya, untuk sebarang
          $x$, ambil $y=1-x$).
        \partsitem Penyelesaian sistem menghasilkan
          \begin{equation*}
            x=\frac{a^4-1}{a^2-1}
            \qquad
            y=\frac{-a^3+a}{a^2-1}.
          \end{equation*}
          Jika $a^2-1\neq 0$, sistem tersebut memiliki solusi tunggal
          $x=a^2+1$, $y=-a$.
          Untuk $a=-1$ dan $a=1$, kita memperoleh sistem-sistem
          \begin{equation*}
            \left\{
            \begin{linsys}{2}
              -x &+  &y  &=  &-1 \\
               x &-  &y  &=  &1
            \end{linsys}\right.
            \qquad
            \left\{
            \begin{linsys}{2}
               x &+  &y  &=  &1 \\
               x &+  &y  &=  &1
            \end{linsys}\right.
         \end{equation*}
         yang keduanya memiliki tak hingga banyak solusi.
      \end{exparts}
    \end{answer}
  \puzzle \item 
    \cite{MathMag52p48}
    Di udara, sebuah bola berlapis emas memiliki massa \( 7588 \) gram.
    Diketahui bahwa bola itu mungkin mengandung satu atau beberapa logam
    aluminium, tembaga, perak, atau timbal.
    Ketika ditimbang secara berurutan dalam kondisi standar di dalam air,
    benzena, alkohol, dan gliserin, massanya masing-masing adalah \( 6588 \),
    \( 6688 \), \( 6778 \), dan \( 6328 \) gram.
    Berapa banyak masing-masing logam tersebut, jika ada, yang terkandung di
    dalam bola jika massa jenis relatif zat-zat itu dianggap sebagai berikut?
    \begin{center}
       \begin{tabular}{lrclr}
         Aluminium &\( 2.7 \)
            &\makebox[3em]{\mbox{}\hfill\mbox{}} &Alkohol &0.81 \\
         Tembaga   &\( 8.9 \)  &     &Benzena  &\( 0.90 \) \\
         Emas      &\( 19.3 \) &     &Gliserin &\( 1.26 \) \\
         Timbal    &\( 11.3 \) &     &Air      &\( 1.00 \) \\
         Perak     &\( 10.5 \)
       \end{tabular}
    \end{center}
    \begin{answer}
      \answerasgiven
      Misalkan \( u \), \( v \), \( x \), \( y \), dan \( z \) berturut-turut
      adalah volume Al, Cu, Pb, Ag, dan Au yang terkandung dalam bola, dalam
      satuan \( {\rm cm}^3 \), dan anggap bola tersebut tidak berongga.
      Karena kehilangan massa di dalam air (massa jenis relatif \( 1.00 \))
      adalah \( 1000 \) gram, volume bola adalah \( 1000\mbox{ cm}^3 \).
      Data tersebut, yang sebagian berlebih tetapi tetap konsisten, hanya
      menghasilkan \( 2 \) persamaan bebas: satu mengenai volume dan satu lagi
      mengenai massa.
      \begin{equation*}
        \begin{linsys}{5}
           u  &+  &v    &+  &x     &+  &y     &+  &z     &=  &1000  \\
        2.7u  &+  &8.9v &+  &11.3x &+  &10.5y &+  &19.3z &=  &7588
        \end{linsys}
      \end{equation*}
      Jelas bahwa bola harus mengandung sejumlah aluminium agar massa jenis
      relatif rata-ratanya lebih rendah daripada massa jenis relatif semua
      logam lainnya.
      Bagian soal ini tidak memiliki hasil yang tunggal, karena banyaknya tiga
      logam dapat dipilih secara sebarang, asalkan pilihan tersebut tidak
      menghasilkan banyaknya logam yang negatif.

      Jika bola hanya mengandung aluminium dan emas, terdapat kira-kira
      \( 294.5\mbox{ cm}^3 \) emas dan \( 705.5\mbox{ cm}^3 \) aluminium.
      Kemungkinan lain adalah kira-kira masing-masing \( 124.7\mbox{ cm}^3 \)
      Cu, Au, Pb, dan Ag, serta \( 501.2\mbox{ cm}^3  \) Al.
    \end{answer}
\end{exercises}

























\subsection{\texorpdfstring{$\text{Umum}=\text{Partikular}+\text{Homogen}$}{Umum=Partikular+Homogen}}
Pada subbagian sebelumnya, semua deskripsi himpunan solusi
mengikuti satu pola.
Setiap deskripsi terdiri atas sebuah vektor yang merupakan solusi partikular
sistem, ditambah kombinasi tanpa pembatasan dari beberapa vektor lain.
Himpunan solusi dari
\nearbyexample{ex:ManyParamsInfManySolsSystem} memperlihatkan pola ini.
\begin{equation*}
  \set{
   \underbracket[.7pt]{
     \colvec[r]{0 \\ 4 \\ 0 \\ 0 \\ 0}}_{\text{\shortstack{\rule{0pt}{2ex}solusi \\
                                                    partikular}}}+
   \underbracket[.7pt]{w\colvec[r]{1 \\ -1 \\ 3 \\ 1 \\ 0}+
       u\colvec[r]{1/2 \\ -1 \\ 1/2 \\ 0 \\ 1}}_{\text{\shortstack{\rule{0pt}{2ex}kombinasi tanpa\\
                                                                pembatasan}}}
       \suchthat w,u\in\Re}
\end{equation*}
Kombinasi itu tidak dibatasi karena
$w$ dan $u$ boleh berupa sembarang bilangan riil\Dash tidak ada
syarat seperti ``sedemikian sehingga $2w-u=0\mspace{1mu}$'' yang membatasi
pasangan $w,u$ yang boleh kita gunakan.

Contoh itu menunjukkan suatu himpunan solusi tak hingga yang mengikuti pola tersebut.
Dua jenis himpunan solusi lainnya juga mengikutinya.
Himpunan solusi berelemen tunggal mengikuti pola karena memiliki
suatu solusi partikular dan bagian kombinasi tanpa pembatasannya bersifat trivial.
Artinya, alih-alih merupakan kombinasi dua vektor atau
satu vektor, bagian itu merupakan kombinasi tanpa vektor.
(Menurut konvensi, jumlah atas himpunan vektor yang kosong
adalah vektor nol.)
Himpunan solusi kosong mengikuti pola karena tidak ada
solusi partikular, sehingga tidak ada jumlah yang berbentuk demikian.

\begin{theorem} \label{th:GenEqPartPlusHomo}
%<*th:GenEqPartPlusHomo>
Himpunan solusi setiap sistem linear
berbentuk
\begin{equation*}
   \set{\vec{p}+c_1\vec{\beta}_1+\,\cdots\,+c_k\vec{\beta}_k
     \suchthat c_1,\,\ldots\,,c_k\in\Re}
\end{equation*}
di mana \( \vec{p} \) adalah sembarang solusi partikular
dan banyaknya vektor
$\vec{\beta}_1$, \ldots, $\vec{\beta}_k$ sama dengan
jumlah variabel bebas yang dimiliki sistem setelah reduksi Gauss.
%</th:GenEqPartPlusHomo>
\end{theorem}

Deskripsi solusi itu terdiri atas dua bagian:
solusi partikular $\vec{p}$
dan kombinasi linear tanpa pembatasan dari vektor-vektor $\vec{\beta}$.
Kita akan membuktikan teorema tersebut dengan dua lema yang bersesuaian.

Kita akan terlebih dahulu berfokus pada kombinasi tanpa pembatasan.
Untuk itu, kita meninjau sistem yang memiliki vektor nol
sebagai solusi partikular,
sehingga $\vec{p}+c_1\vec{\beta}_1+\dots+c_k\vec{\beta}_k$
dapat kita ringkas menjadi $c_1\vec{\beta}_1+\dots+c_k\vec{\beta}_k$.

\begin{definition} \label{df:HomogeneousEquation}
%<*df:HomogeneousEquation>
Suatu persamaan linear disebut \definend{homogen}\index{persamaan homogen}%
\index{persamaan linear!homogen} jika konstantanya nol, sehingga
dapat ditulis sebagai $a_1x_1+a_2x_2+\,\cdots\,+a_nx_n=0$.
%</df:HomogeneousEquation>
\end{definition}

\begin{example}  \label{ex:FirstExHomoSys}
Untuk sembarang sistem linear seperti
\begin{equation*}
  \begin{linsys}{2}
    3x  &+  &4y  &=  3  \\
    2x  &-  &y   &=  1  
  \end{linsys}
\end{equation*}
kita membentuk sistem persamaan homogen yang berkaitan dengan menetapkan
ruas kanannya sama dengan nol.
\begin{equation*}
  \begin{linsys}{2}
    3x  &+  &4y  &=  0  \\
    2x  &-  &y   &=  0  
  \end{linsys}
\end{equation*}
Bandingkan reduksi sistem semula
\begin{equation*}
  \begin{linsys}{2}
    3x  &+  &4y  &=  3  \\
    2x  &-  &y   &=  1  
  \end{linsys}
  \grstep{-(2/3)\rho_1+\rho_2}
  \begin{linsys}{2}
    3x  &+  &4y        &=  3  \\
        &   &-(11/3)y   &=  -1  
   \end{linsys}
\end{equation*}
dengan reduksi sistem homogen yang berkaitan.
\begin{equation*}
  \begin{linsys}{2}
    3x  &+  &4y  &=  0  \\
    2x  &-  &y   &=  0  
  \end{linsys}
  \grstep{-(2/3)\rho_1+\rho_2}
  \begin{linsys}{2}
    3x  &+  &4y        &=  0  \\
        &   &-(11/3)y   &=  0
   \end{linsys}
\end{equation*}
Jelas bahwa kedua reduksi berlangsung dengan cara yang sama.
Karena itu, kita dapat mempelajari cara mereduksi suatu sistem linear
dengan mempelajari cara mereduksi sistem homogen yang berkaitan.
\end{example}

Mempelajari sistem homogen yang berkaitan memberikan keuntungan besar
dibandingkan mempelajari sistem semula.
Sistem tak homogen dapat bersifat tak konsisten.
Namun, sistem homogen pasti konsisten karena selalu memiliki setidaknya
satu solusi, yaitu vektor nol.


\begin{example} \label{ex:HomoZeroOnlySol}
Beberapa sistem homogen hanya memiliki vektor nol sebagai solusi.
\begin{equation*}
   \begin{linsys}{3}
     3x  &+  &2y  &+  &z  &=  &0  \\
     6x  &+  &4y  &   &   &=  &0  \\
         &   &y   &+  &z  &=  &0  
   \end{linsys}
   \grstep{-2\rho_1 +\rho_2}
   \begin{linsys}{3}
      3x  &+  &2y  &+  &z  &=  &0  \\
          &   &    &   &-2z&=  &0  \\
          &   &y   &+  &z  &=  &0  
    \end{linsys}
   \grstep{\rho_2 \leftrightarrow\rho_3}
   \begin{linsys}{3}
      3x  &+  &2y  &+  &z  &=  &0  \\
          &   &y   &+  &z  &=  &0  \\
          &   &    &   &-2z&=  &0
    \end{linsys}
\end{equation*}
\end{example}

\begin{example} \label{ex:SolnChemProb}
Beberapa sistem homogen memiliki banyak solusi.
Salah satunya adalah masalah Kimia\index{masalah Kimia}
pada halaman pertama subbagian pertama.
\begin{align*}
  \begin{linsys}{4}
              7x  &   &   &-  &7z  &   &   &=  &0  \\
              8x  &+  &y  &-  &5z  &-  &2w &=  &0  \\
                  &   &y  &-  &3z  &   &   &=  &0  \\
                  &   &3y &-  &6z  &-  &w  &=  &0  
  \end{linsys}
  &\grstep{-(8/7)\rho_1+\rho_2}
  \begin{linsys}{4}
                7x &   &   &-  &7z  &   &   &=  &0  \\
                   &   &y  &+  &3z  &-  &2w &=  &0  \\
                   &   &y  &-  &3z  &   &   &=  &0  \\
                   &   &3y &-  &6z  &-  &w  &=  &0  
   \end{linsys}                                        \\
  &\grstep[-3\rho_2+\rho_4]{-\rho_2+\rho_3}
  \begin{linsys}{4}
                7x &   &   &-  &7z  &   &   &=  &0  \\
                   &   &y  &+  &3z  &-  &2w &=  &0  \\
                   &   &   &   &-6z &+  &2w &=  &0  \\
                   &   &   &   &-15z&+  &5w &=  &0  
   \end{linsys}                                        \\
  &\grstep{-(5/2)\rho_3+\rho_4}
  \begin{linsys}{4}
                7x &   &   &-  &7z  &   &   &=  &0  \\
                   &   &y  &+  &3z  &-  &2w &=  &0  \\
                   &   &   &   &-6z &+  &2w &=  &0  \\
                   &   &   &   &    &   &0  &=  &0  
   \end{linsys}
\end{align*}
Himpunan solusi
\begin{equation*}
  \set{\colvec[r]{1/3 \\ 1 \\ 1/3 \\ 1}w \suchthat w\in\Re}
\end{equation*}
memiliki banyak vektor selain vektor nol
(jika \( w \) menyatakan jumlah molekul, solusi hanya bermakna
ketika \( w \) merupakan kelipatan tak negatif dari $3$).
\end{example}

\begin{lemma} \label{le:HomoSltnSpanVecs}
%<*le:HomoSltnSpanVecs>
Untuk setiap sistem linear homogen terdapat
vektor-vektor $\vec{\beta}_1$, \ldots, $\vec{\beta}_k$ sedemikian sehingga
himpunan solusi sistem tersebut adalah
\begin{equation*}
  \set{c_1\vec{\beta}_1+\cdots+c_k\vec{\beta}_k \suchthat c_1,\ldots,c_k\in\Re}
\end{equation*}
di mana $k$ adalah jumlah variabel bebas dalam bentuk eselon sistem tersebut.
%</le:HomoSltnSpanVecs>
\end{lemma}

Sebelum pembuktian, kita bahas dua hal yang berkaitan dengannya.
Pertama, pembuktian akan menunjukkan bahwa, untuk suatu sistem homogen
dalam bentuk eselon, substitusi balik dapat digunakan untuk menyatakan
semua variabel utama dalam bentuk variabel-variabel bebas.
Untuk melihat mengapa hal ini cukup untuk membuktikan lema, kita berikan
sebuah contoh (uraian langkah demi langkah ini juga akan kita gunakan
untuk memperkenalkan notasi dalam pembuktian).

%<*ex:HomoSystemGivesSpanningSet0>
Perhatikan sistem persamaan homogen dalam bentuk eselon berikut.
\begin{equation*}
  \begin{linsys}{5}
     x  &+  &y   &+  &2z  &+  &u  &+  &v  &=  &0  \\
        &   &y   &+  &z  &+  &u  &-  &v  &=  &0  \\
        &   &    &   &   &   &u  &+  &v  &=  &0  
  \end{linsys}
\end{equation*}
Sistem ini memiliki lima variabel: $x_1=x$, \ldots{}, $x_5=v$.
Untuk melakukan substitusi balik, kita mulai dari baris paling bawah,
yaitu persamaan~$m=3$.
Kita selesaikan variabel utamanya, $x_{\ell_3}=x_4=u$, dengan menggunakan
suku utama $a_{m,\ell_m}x_{\ell_m}=a_{3,4}x_4=1\cdot u$
(lambang `$\ell$' berasal dari kata Inggris ``leading'', yang berarti utama).
Kita memperoleh $u=-v$.
%</ex:HomoSystemGivesSpanningSet0>

Selanjutnya, substitusi balik memasuki tahap iterasi.
%<*ex:HomoSystemGivesSpanningSet1>
Kita akan menggunakan variabel~$t$ untuk menghitung jarak baris yang sedang
ditinjau dari baris paling bawah. Jadi, kita ambil $t=1$ dan meninjau
baris $m-t=3-1=2$.
Saat mencapai baris ini, kita telah menyatakan variabel utama pada baris~$3$
dalam bentuk variabel bebas.
Kita mengganti $x_{\ell_3}=x_4=u$ dengan bentuknya dalam variabel bebas,
sehingga diperoleh $y+z+(-v)-v=0$.
Kemudian, dengan menyelesaikan variabel utama pada baris ini, kita memperoleh
$x_{\ell_{m-t}}=x_2=y=-z+2v$.

Berikutnya adalah baris yang berjarak $t=2$ dari baris paling bawah,
yaitu baris $m-t=1$.
Sejauh ini, kita telah menyatakan variabel utama pada baris $3$ dan~$2$
dalam bentuk variabel bebas.
Ganti $x_{\ell_3}=u$ dan $x_{\ell_2}=y$, sehingga diperoleh
$x+(-z+2v)+2z+(-v)+v=0$.
Kemudian selesaikan variabel utama dalam bentuk variabel bebas,
yakni $x_{\ell_1}=x=-z-2v$.
%</ex:HomoSystemGivesSpanningSet2>

%<*ex:HomoSystemGivesSpanningSet3>
Pokok contoh ini adalah bahwa sekarang pekerjaan kita telah selesai.
Cukup tulis solusi dalam notasi vektor
\begin{equation*}
  \colvec{x \\ y \\ z \\ u \\ v}
  =\colvec{-1 \\ -1 \\ 1 \\ 0 \\ 0}z
   +\colvec{-2 \\ 2 \\ 0 \\ -1 \\ 1}v 
  \qquad \text{dengan $z,v\in\Re$}
\end{equation*}
untuk melihat bahwa $\vec{\beta}_1$ dan $\vec{\beta}_2$
dalam pernyataan lema adalah vektor-vektor yang berkaitan dengan
variabel bebas $z$ dan~$v$.
%</ex:HomoSystemGivesSpanningSet3>

Hal kedua mengenai pembuktian juga tampak dari contoh tersebut.
Substitusi balik bergerak naik di dalam sistem dan menggunakan hasil
dari baris-baris yang lebih bawah untuk mengolah baris berikutnya.
Hal ini menyarankan strategi pembuktian dengan
induksi matematika.\index{induksi matematika}\index{teknik pembuktian!induksi}\index{induksi}%

Induksi adalah teknik pembuktian penting yang tidak serta-merta jelas dan akan
muncul beberapa kali dalam buku ini.
Pembuktian dengan induksi akan kita lakukan dalam dua langkah: langkah dasar
dan langkah induksi.
Pada langkah dasar, kita memeriksa bahwa pernyataan berlaku untuk kasus
pertama. Di sini, kita memeriksa bahwa pada sistem linear berbentuk eselon,
variabel utama dalam persamaan paling bawah dapat ditulis dalam bentuk
variabel bebas.
Pada langkah induksi, kita harus membuktikan suatu implikasi: jika pernyataan
berlaku untuk semua kasus sebelumnya, maka pernyataan itu juga berlaku
untuk kasus yang sedang ditinjau.
Secara khusus, hipotesis induksi di sini menyatakan bahwa jika variabel utama
pada baris-baris yang lebih bawah dapat kita tulis dalam bentuk variabel bebas,
maka variabel utama pada baris berikutnya di atasnya juga dapat kita tulis
dalam bentuk tersebut.

Kedua langkah itu bersama-sama membuktikan pernyataan untuk semua baris.
Menurut langkah dasar, pernyataan berlaku untuk persamaan paling bawah;
menurut langkah induksi, fakta bahwa pernyataan berlaku untuk persamaan
paling bawah menunjukkan bahwa pernyataan itu berlaku pula untuk persamaan
berikutnya di atasnya.
Penerapan langkah induksi sekali lagi kemudian menunjukkan bahwa pernyataan
berlaku untuk persamaan ketiga dari bawah, dan seterusnya.\appendrefs{induksi matematika}%

\begin{proof}
%<*pf:HomoSltnSpanVecs0>
Terapkan Metode Gauss untuk memperoleh bentuk eselon.
Mungkin terdapat beberapa persamaan \( 0=0 \); persamaan-persamaan ini kita abaikan
(jika sistem hanya terdiri atas persamaan \( 0=0 \), lema berlaku secara trivial
karena tidak ada variabel utama).
% If there are any variables that appear only with with a coefficient of
% $0$ then they are free, and will never lead a row, so we will
% ignore that possibility.  
Namun, karena sistem bersifat homogen, tidak ada persamaan yang kontradiktif.

Kita akan menggunakan induksi untuk menunjukkan bahwa
setiap variabel utama dapat ditulis dalam bentuk variabel bebas.
Hal itu akan menyelesaikan pembuktian karena variabel-variabel bebas
dapat digunakan sebagai parameter dan vektor-vektor $\vec{\beta}$
adalah vektor koefisien dari variabel-variabel bebas tersebut.
%</pf:HomoSltnSpanVecs0>

%<*pf:HomoSltnSpanVecs1>
Untuk langkah dasar, perhatikan persamaan paling bawah,
\begin{equation*}
  a_{m,\ell_m}x_{\ell_m}+a_{m,1+\ell_{m}}x_{1+\ell_{m}}+\cdots+a_{m,n}x_n=0
   % \tag{$*$}
\end{equation*}
di mana \( a_{m,\ell_m}\neq 0 \).
%</pf:HomoSltnSpanVecs1>
%<*pf:HomoSltnSpanVecs2>
Karena ini adalah baris paling bawah, semua variabel % $x_{\ell_{m}+1}$, \ldots{}
setelah variabel utama pasti merupakan variabel bebas.
Pindahkan variabel-variabel tersebut ke ruas kanan dan bagi dengan $a_{m,\ell_m}$
\begin{equation*}
  x_{\ell_m}
  =(-a_{m,1+\ell_{m}}/a_{m,\ell_m})x_{1+\ell_{m}}+\cdots+(-a_{m,n}/a_{m,\ell_m})x_n
\end{equation*}
untuk menyatakan variabel utama dalam bentuk variabel bebas.
%</pf:HomoSltnSpanVecs2>
(Satu catatan teknis:~jika pada persamaan paling bawah tidak terdapat
variabel di sebelah kanan~$x_{\ell_m}$, maka \( x_{\ell_m}=0 \).
Hal ini memenuhi pernyataan yang sedang kita buktikan karena, seperti
disinggung pada awal subbagian ini, \( x_{\ell_m} \) ditulis sebagai jumlah
sejumlah variabel bebas, yakni jumlah nol buah variabel, berdasarkan konvensi
bahwa jumlah trivial bernilai~$0$.)

%<*pf:HomoSltnSpanVecs3>
Untuk langkah induksi, andaikan bahwa pada persamaan ke-\( m \),
persamaan ke-\text{\( (m-1) \)}, dan seterusnya, hingga dan termasuk
persamaan ke-\mbox{\( m-(t-1) \)}, variabel utama dapat kita tulis
dalam bentuk variabel bebas.
Dengan kata lain, andaikan bahwa pernyataan tentang kemampuan menulis
variabel utama suatu baris dalam bentuk variabel bebas berlaku untuk
$t$~baris paling bawah, dengan $1\leq t<m$.
Kita harus menunjukkan bahwa pernyataan tersebut juga berlaku untuk
persamaan berikutnya di atasnya, yaitu persamaan ke-\( (m-t) \).

Untuk setiap \( x_{\ell_m} \), \ldots, \( x_{\ell_{m-(t-1)}} \),
substitusikan bentuknya dalam variabel bebas.
Hasilnya memiliki suku utama
$a_{m-t,\ell_{m-t}}x_{\ell_{m-t}}$ dengan \( a_{m-t,\ell_{m-t}}\neq 0 \),
sedangkan bagian lain dari ruas kiri merupakan kombinasi variabel bebas.
Pindahkan variabel-variabel bebas ke ruas kanan dan bagi dengan
\( a_{m-t,\ell_{m-t}} \), sehingga \( x_{\ell_{m-t}} \) akhirnya
tertulis dalam bentuk variabel bebas.

Kita telah menyelesaikan langkah dasar dan langkah induksi. Oleh karena itu,
berdasarkan prinsip induksi matematika, proposisi tersebut benar.
%</pf:HomoSltnSpanVecs3>
\end{proof}

Hal ini menunjukkan, seperti yang dibahas antara lema dan pembuktiannya,
bahwa himpunan solusi dapat kita parameterkan dengan menggunakan
variabel-variabel bebas.
Kita mengatakan bahwa himpunan vektor
$\set{c_1\vec{\beta}_1+\cdots+c_k\vec{\beta}_k \suchthat c_1,\ldots,c_k\in\Re}$
\definend{dibangkitkan oleh}\index{dibangkitkan oleh}\index{dibangkitkan}
atau~\definend{direntang oleh}\index{direntang oleh}\index{rentang}
himpunan
\( \set{\smash{\vec{\beta}_1},\ldots,\smash{\vec{\beta}_k}} \).

% The proof mentions a tricky point.
% We follow the convention that the sum of an empty set of vectors is the 
% zero vector.
% In particular, we needs this where 
% a homogeneous system has a unique solution because it
% takes the \( c \)'s to be the free variables
% and if there is a unique solution then there are no free variables.

Untuk menuntaskan pembuktian \nearbytheorem{th:GenEqPartPlusHomo},
lema berikutnya meninjau bagian solusi partikular dalam deskripsi
himpunan solusi.

\begin{lemma}  \label{th:GenEqPartHomo}
%<*th:GenEqPartHomo>
Untuk suatu sistem linear dan sembarang solusi partikular $\vec{p}\/$,
himpunan solusinya sama dengan
% \begin{equation*}
$
  \set{\vec{p}+\vec{h} \suchthat \text{ \( \vec{h} \) memenuhi
                                sistem homogen yang berkaitan}     }$.
%\end{equation*}
%</th:GenEqPartHomo>
\end{lemma}

% So fixing any particular solution gives the above 
% description of the solution set.

\begin{proof}
Kita akan menunjukkan inklusi timbal balik antarkedua himpunan: setiap solusi
sistem berada dalam himpunan di atas, dan setiap anggota himpunan tersebut
merupakan solusi sistem.\appendrefs{kesamaan himpunan}

%<*pf:GenEqPartHomo0>
Untuk inklusi himpunan dalam arah pertama, yakni bahwa jika suatu vektor
menyelesaikan sistem maka vektor tersebut berada dalam himpunan di atas,
andaikan \( \vec{s} \) menyelesaikan sistem.
Maka \( \vec{s}-\vec{p} \) menyelesaikan sistem homogen yang berkaitan,
karena untuk setiap indeks persamaan \( i \),
\begin{multline*}
  a_{i,1}(s_1-p_1)+\cdots+a_{i,n}(s_n-p_n) \\
  =(a_{i,1}s_1+\cdots+a_{i,n}s_n)       
  -(a_{i,1}p_1+\cdots+a_{i,n}p_n)  
  =d_i-d_i                 
  =0
\end{multline*}
di mana \( p_j \) dan \( s_j \) masing-masing merupakan komponen ke-\( j \)
dari \( \vec{p} \) dan \( \vec{s} \).
Tuliskan \( \vec{s} \) dalam bentuk \( \vec{p}+\vec{h} \) yang diperlukan
dengan menuliskan \( \vec{s}-\vec{p} \) sebagai \( \vec{h} \).
%</pf:GenEqPartHomo0>

%<*pf:GenEqPartHomo1>
Untuk inklusi himpunan dalam arah sebaliknya, ambil suatu vektor berbentuk
$\vec{p}+\vec{h}$, dengan \( \vec{p} \) menyelesaikan sistem dan
\( \vec{h} \) menyelesaikan sistem homogen yang berkaitan. Perhatikan bahwa
$\vec{p}+\vec{h}$ menyelesaikan sistem yang diberikan karena untuk setiap
indeks persamaan~$i$,
\begin{multline*}
  a_{i,1}(p_1+h_1)+\cdots+a_{i,n}(p_n+h_n)  \\
  =(a_{i,1}p_1+\cdots+a_{i,n}p_n)      
   +(a_{i,1}h_1+\cdots+a_{i,n}h_n)  
  =d_i+0                                
  =d_i
\end{multline*}
di mana, seperti sebelumnya, \( p_j \) dan \( h_j \) masing-masing merupakan
komponen ke-\( j \) dari \( \vec{p} \) dan \( \vec{h} \).
%</pf:GenEqPartHomo1>
\end{proof}

Kedua lema tersebut bersama-sama membuktikan
\nearbytheorem{th:GenEqPartPlusHomo}.
Ingat teorema itu melalui slogan
``\( \text{Umum} = \text{Partikular} + \text{Homogen} \)''.

\begin{example} \label{ex:IllusGenEqPartHomo}
Sistem berikut mengilustrasikan \nearbytheorem{th:GenEqPartPlusHomo}.
\begin{equation*}
  \begin{linsys}{3}
    x  &+  &2y  &-  &z  &=  &1  \\
    2x &+  &4y  &   &   &=  &2  \\
       &   &y   &-  &3z &=  &0
  \end{linsys}
\end{equation*}
Metode Gauss
\begin{equation*}
  \grstep{-2\rho_1+\rho_2}
  \begin{linsys}{3}
    x  &+  &2y  &-  &z  &=  &1  \\
       &   &    &   &2z &=  &0  \\
       &   &y   &-  &3z &=  &0
  \end{linsys}                           
  \grstep{\rho_2\leftrightarrow\rho_3} 
  \begin{linsys}{3}
      x  &+  &2y  &-  &z  &=  &1  \\      
         &   &y   &-  &3z &=  &0  \\
         &   &    &   &2z &=  &0
   \end{linsys}
\end{equation*}
menunjukkan bahwa solusi umumnya merupakan himpunan berelemen tunggal.
\begin{equation*}
  \set{\colvec[r]{1 \\ 0 \\ 0} }
\end{equation*}
Vektor tunggal tersebut jelas merupakan suatu solusi partikular.
Sistem homogen yang berkaitan direduksi melalui operasi baris yang sama
\begin{equation*}
  \begin{linsys}{3}
    x  &+  &2y  &-  &z  &=  &0  \\
    2x &+  &4y  &   &   &=  &0  \\
       &   &y   &-  &3z &=  &0
  \end{linsys}
  \grstep{-2\rho_1+\rho_2}
  \repeatedgrstep{\rho_2\swap\rho_3} 
  \begin{linsys}{3}
      x  &+  &2y  &-  &z  &=  &0  \\      
         &   &y   &-  &3z &=  &0  \\
         &   &    &   &2z &=  &0
   \end{linsys}
\end{equation*}
juga menghasilkan himpunan berelemen tunggal.
\begin{equation*}
  \set{\colvec[r]{0 \\ 0 \\ 0} }
\end{equation*}
Jadi, seperti dibahas pada awal subbagian ini, dalam kasus solusi tunggal
ini, solusi umum diperoleh dengan mengambil solusi partikular dan
menambahkan padanya solusi tunggal dari sistem homogen yang berkaitan.
\end{example}

\begin{example}
Pada awal subbagian ini juga dibahas bahwa kasus ketika himpunan solusi umum
kosong tetap mengikuti pola
$\text{Umum}=\text{Partikular}+\text{Homogen}$.
Sistem berikut mengilustrasikannya.
\begin{equation*}
  \begin{linsys}{4}
    x  &   &  &+  &z  &+ &w  &=  &-1  \\
   2x  &-  &y &   &   &+ &w  &=  &3   \\
    x  &+  &y &+  &3z &+ &2w &=  &1   
  \end{linsys}
  \grstep[-\rho_1+\rho_3]{-2\rho_1+\rho_2}
  \begin{linsys}{4}
    x  &   &  &+  &z  &+ &w  &=  &-1  \\
       &   &-y&-  &2z &- &w  &=  &5   \\
       &   &y &+  &2z &+ &w  &=  &2   
   \end{linsys}
\end{equation*}
Sistem tersebut tidak memiliki solusi karena dua persamaan terakhir
saling bertentangan.
Namun, sistem homogen yang berkaitan memiliki solusi, sebagaimana semua
sistem homogen.
\begin{equation*}
  \begin{linsys}{4}
    x  &   &  &+  &z  &+ &w  &=  &0   \\
   2x  &-  &y &   &   &+ &w  &=  &0   \\
    x  &+  &y &+  &3z &+ &2w &=  &0   
  \end{linsys}
  \grstep[-\rho_1+\rho_3]{-2\rho_1+\rho_2}
  \grstep{\rho_2+\rho_3}
  \begin{linsys}{4}
    x  &   &  &+  &z  &+ &w  &=  &0   \\
       &   &-y&-  &2z &- &w  &=  &0   \\
       &   &  &   &   &  &0  &=  &0         
  \end{linsys}
\end{equation*}
Bahkan, himpunan solusinya tak hingga.
\begin{equation*}
  \set{\colvec[r]{-1 \\ -2 \\ 1 \\ 0}z+\colvec[r]{-1 \\ -1 \\ 0 \\ 1}w
         \suchthat z,w\in\Re}
\end{equation*}
Meskipun demikian, karena sistem semula tidak memiliki solusi partikular,
himpunan solusi umumnya kosong\Dash tidak ada vektor berbentuk
$\vec{p}+\vec{h}$ karena tidak ada $\vec{p}\:$.
\end{example}

\begin{corollary} \label{co:ThreeKindsSolutionSets}
%<*co:ThreeKindsSolutionSets>
Himpunan solusi sistem linear dapat berupa himpunan kosong, memiliki satu
elemen, atau memiliki tak hingga banyak elemen.
%</co:ThreeKindsSolutionSets>
\end{corollary}

\begin{proof}
%<*pf:ThreeKindsSolutionSets0>
Kita telah melihat contoh ketiga kemungkinan tersebut, sehingga kita hanya
perlu membuktikan bahwa tidak ada kemungkinan lain.

Pertama, perhatikan bahwa suatu sistem homogen yang memiliki setidaknya
satu solusi yang tidak sama dengan \( \zero \), yaitu $\vec{v}$, memiliki
tak hingga banyak solusi.
Hal ini karena sembarang kelipatan skalar dari~$\vec{v}$ juga menyelesaikan
sistem homogen tersebut, dan himpunan kelipatan skalar $\vec{v}$ memuat
tak hingga banyak vektor: jika $s,t\in\Re$ berbeda, maka
$s\vec{v}\neq t\vec{v}$, sebab
$s\vec{v}-t\vec{v}=(s-t)\vec{v}$ tidak sama dengan~$\zero$. Memang, setiap
komponen $\vec{v}$ yang tidak sama dengan $0$ tetap menghasilkan nilai yang
tidak sama dengan $0$ setelah dikalikan dengan faktor $s-t$ yang tidak sama
dengan $0$.
%</pf:ThreeKindsSolutionSets0>

%<*pf:ThreeKindsSolutionSets1>
Sekarang terapkan \nearbylemma{th:GenEqPartHomo} untuk menyimpulkan bahwa
suatu himpunan solusi
\begin{equation*}
  \set{\vec{p}+\vec{h}\suchthat
    \text{\( \vec{h} \) menyelesaikan sistem homogen yang berkaitan}}
\end{equation*}
bersifat kosong (jika tidak ada solusi partikular \( \vec{p} \)), memiliki
satu elemen (jika terdapat suatu \( \vec{p} \) dan sistem homogen memiliki
solusi tunggal \( \zero \)), atau bersifat tak hingga (jika terdapat suatu
\( \vec{p} \) dan sistem homogen memiliki solusi yang tidak sama dengan
$\zero$, sehingga menurut
paragraf sebelumnya sistem itu memiliki tak hingga banyak solusi).
%</pf:ThreeKindsSolutionSets1>
\end{proof}

Tabel berikut merangkum faktor-faktor yang memengaruhi banyaknya anggota
dalam himpunan solusi umum.

\smallskip
%<*table:KindsSolutionSets>
\begin{center} % \small 
\begin{tabular}{r@{}c}
  &\hspace*{2.5em}\begin{tabular}{c} 
      \textit{banyaknya solusi} \\[-.5ex]
      \textit{sistem homogen}
    \end{tabular}                            \\[1.7ex] 
  \begin{tabular}{r@{\hspace*{.5em}}} 
     \ \\[.6ex]
     \textit{solusi} \\[-.55ex]
     \textit{partikular}   \\[-.5ex]
     \textit{ada?}
  \end{tabular}
  &\begin{tabular}{r|c@{\hspace*{1em}}c} % \cline{2-3}   
     \multicolumn{1}{c}{\ }
         &\textit{satu}    &\textit{tak hingga}                    \\
     \cline{2-3}
     \textit{ya}
        &\rule{0ex}{16pt}\begin{tabular}{@{}c@{}} solusi \\[-.5ex] tunggal \end{tabular}
        &\begin{tabular}{@{}c@{}} tak hingga banyak \\[-.5ex] solusi \end{tabular}
         \\[2ex] % \hline  
     \textit{tidak}
        &\begin{tabular}{@{}c@{}} tidak ada \\[-.5ex] solusi \end{tabular}
        &\begin{tabular}{@{}c@{}} tidak ada \\[-.5ex] solusi \end{tabular}
         \\ % \hline
   \end{tabular}
\end{tabular}
\end{center}
%</table:KindsSolutionSets>
\smallskip

Dimensi di bagian atas tabel adalah dimensi yang lebih sederhana.
Ketika kita menerapkan Metode Gauss pada suatu sistem linear dengan
mengabaikan konstanta di ruas kanan dan hanya memperhatikan koefisien
di ruas kiri, pada akhirnya setiap variabel menjadi variabel utama pada
suatu baris, atau kita menemukan variabel yang tidak menjadi variabel utama
pada baris mana pun, yakni suatu variabel bebas.
(Kita memformalkan tindakan ``mengabaikan konstanta di ruas kanan'' dengan
meninjau sistem homogen yang berkaitan.)

Kasus khusus yang penting adalah sistem yang memiliki jumlah persamaan
sama dengan jumlah variabel tak diketahui.
Sistem seperti itu memiliki solusi tunggal jika dan hanya jika sistem tersebut
tereduksi menjadi sistem berbentuk eselon yang setiap variabelnya merupakan
variabel utama pada suatu baris (karena banyaknya variabel sama dengan
banyaknya baris). Hal ini terjadi jika dan hanya jika sistem homogen yang
berkaitan memiliki solusi tunggal.

\begin{definition} \label{df:Nonsingular}
%<*df:Nonsingular>
Suatu matriks persegi disebut \definend{nonsingular}\index{nonsingular!matriks}
\index{matriks!nonsingular}
jika matriks tersebut merupakan matriks koefisien suatu sistem homogen
yang memiliki solusi tunggal.
Jika tidak, matriks itu disebut
\definend{singular}\index{singular!matriks}\index{matriks!singular};
dengan kata lain, jika matriks tersebut merupakan matriks koefisien suatu
sistem homogen yang memiliki tak hingga banyak solusi.
%</df:Nonsingular>
\end{definition}

\begin{example}
Matriks pertama berikut bersifat nonsingular, sedangkan yang kedua singular,
\begin{equation*}
  \begin{mat}[r]
    1  &2  \\
    3  &4
  \end{mat}
  \qquad
  \begin{mat}[r]
    1  &2  \\
    3  &6
  \end{mat}
\end{equation*}
karena sistem homogen pertama berikut memiliki solusi tunggal, sedangkan
yang kedua memiliki tak hingga banyak solusi.
\begin{equation*}
  \begin{linsys}[b]{2}
    x &+  &2y  &=  &0  \\
   3x &+  &4y  &=  &0  
  \end{linsys}
  \qquad
  \begin{linsys}[b]{2}
    x &+  &2y  &=  &0  \\
   3x &+  &6y  &=  &0
  \end{linsys}
\end{equation*}  
Pembedaan dalam definisi tersebut dibuat karena suatu sistem dengan jumlah
persamaan sama dengan jumlah variabel dapat berperilaku dengan salah satu
dari dua cara, bergantung pada apakah matriks koefisiennya nonsingular atau
singular.
Jika matriks koefisiennya nonsingular, sistem memiliki solusi tunggal untuk
sembarang konstanta di ruas kanan. Sebagai contoh, Metode Gauss menunjukkan
bahwa sistem berikut
\begin{equation*}
  \begin{linsys}{2}
    x  &+  &2y  &=  &a \\
    3x &+  &4y  &=  &b
  \end{linsys}
\end{equation*}
memiliki solusi tunggal $x=b-2a$ dan $y=(3a-b)/2$.
Sebaliknya, jika matriks koefisiennya singular, sistem tidak pernah memiliki
solusi tunggal\Dash sistem itu tidak memiliki solusi atau memiliki tak hingga
banyak solusi, seperti dua sistem berikut.
\begin{equation*}
  \begin{linsys}[b]{2}
    x  &+  &2y  &=   &1   \\
   3x  &+  &6y  &=   &2   
  \end{linsys}
  \qquad
  \begin{linsys}[b]{2}
    x  &+  &2y  &=   &1   \\
   3x  &+  &6y  &=   &3
  \end{linsys}
\end{equation*} 
\end{example}

Definisi tersebut menggunakan kata `singular' karena kata itu berarti
``menyimpang dari harapan umum.''
Orang sering kali secara naif mengharapkan bahwa sistem dengan jumlah variabel
sama dengan jumlah persamaan akan memiliki solusi tunggal.
Karena itu, kata tersebut dapat kita anggap berkonotasi ``bermasalah,'' atau
setidaknya ``tidak ideal.''
(Agak ironis bahwa `singular' diterapkan pada sistem yang tidak pernah memiliki
tepat satu solusi, tetapi itulah istilah bakunya.)

\begin{example}
Sistem-sistem dari \nearbyexample{ex:FirstExHomoSys},
\nearbyexample{ex:HomoZeroOnlySol},
dan \nearbyexample{ex:IllusGenEqPartHomo}
masing-masing memiliki sistem homogen terkait dengan solusi tunggal.
Dengan demikian, matriks-matriks berikut bersifat nonsingular.
\begin{equation*}
  \begin{mat}[r]
    3  &4  \\
    2  &-1
  \end{mat}
  \qquad
  \begin{mat}[r]
    3  &2   &1  \\
    6  &4   &0  \\
    0  &1   &1
  \end{mat}
  \qquad
  \begin{mat}[r]
    1  &2  &-1 \\
    2  &4  &0  \\
    0  &1  &-3
  \end{mat}
\end{equation*}
Masalah Kimia dari \nearbyexample{ex:SolnChemProb}
merupakan sistem homogen dengan lebih dari satu solusi, sehingga matriksnya
bersifat singular.
\begin{equation*}
  \begin{mat}[r]
    7  &0  &-7 &0  \\
    8  &1  &-5 &-2 \\
    0  &1  &-3 &0  \\
    0  &3  &-6 &-1
  \end{mat}
\end{equation*}
\end{example}

Tabel di atas memiliki dua dimensi.
Kita telah meninjau dimensi di bagian atas:~kita dapat menentukan kolom bagi
suatu sistem linear hanya dengan meninjau ruas kiri sistem; konstanta di ruas
kanan tidak berperan dalam penentuan ini.

Dimensi lain pada tabel, yaitu menentukan apakah terdapat solusi partikular,
lebih sulit.
Perhatikan dua sistem berikut dengan ruas kiri yang sama tetapi ruas kanan
yang berbeda.
\begin{equation*}
  \begin{linsys}[b]{2}
    3x &+ &2y &= &5  \\
    3x &+ &2y &= &5
  \end{linsys}
  \qquad
  \begin{linsys}[b]{2}
    3x &+ &2y &= &5  \\
    3x &+ &2y &= &4
  \end{linsys}
\end{equation*}
Sistem pertama memiliki solusi, sedangkan sistem kedua tidak. Jadi, dalam
kasus ini, konstanta di ruas kanan menentukan apakah sistem memiliki solusi.
Kita mungkin menduga bahwa ruas kiri suatu sistem linear menentukan banyaknya
solusi, sedangkan ruas kanan menentukan keberadaan solusi, tetapi dugaan itu
tidak benar.
Bandingkan dua sistem berikut, yang memiliki ruas kanan sama tetapi ruas kiri
berbeda.
\begin{equation*}
  \begin{linsys}[b]{2}
    3x &+ &2y &= &5  \\
    4x &+ &2y &= &4
  \end{linsys}
  \qquad
  \begin{linsys}[b]{2}
    3x &+ &2y &= &5  \\
    3x &+ &2y &= &4
  \end{linsys}
\end{equation*}
Sistem pertama memiliki solusi, tetapi sistem kedua tidak.
Jadi, konstanta di ruas kanan suatu sistem tidak dengan sendirinya menentukan
apakah solusi ada. Keberadaan solusi justru bergantung pada suatu interaksi
antara ruas kiri dan ruas kanan.

Untuk memperoleh intuisi mengenai interaksi tersebut, perhatikan sistem
berikut yang salah satu koefisiennya dibiarkan tak ditentukan sebagai
variabel~$c$.
\begin{equation*}
  \begin{linsys}{3}
    x  &+  &2y  &+  &3z  &=  &1  \\
    x  &+  &y   &+  &z   &=  &1  \\
   cx  &+  &3y  &+  &4z  &=  &0
  \end{linsys}
\end{equation*}
Jika \( c=2 \), sistem ini tidak memiliki solusi karena baris ketiga di ruas
kiri merupakan jumlah dua baris pertama, sedangkan konstanta ketiga di ruas
kanan bukan jumlah dua konstanta pertama.
Jika \( c\neq 2 \), sistem ini memiliki solusi tunggal (cobalah dengan \( c=1 \)).
Agar suatu sistem memiliki solusi, jika satu baris matriks koefisien di ruas
kiri merupakan kombinasi linear dari baris-baris lain, maka konstanta pada
baris yang bersesuaian di ruas kanan harus merupakan kombinasi yang sama dari
konstanta-konstanta milik baris-baris lain tersebut.

Intuisi lebih lanjut mengenai interaksi tersebut diperoleh dengan mempelajari
kombinasi linear.
Itulah fokus kita dalam bab kedua, setelah kita menuntaskan pembahasan
Metode Gauss pada sisa bab ini.

\begin{exercises}
  \item Selesaikan sistem berikut.
    Kemudian selesaikan sistem homogen yang berkaitan.
    \begin{equation*}
      \begin{linsys}{4}
        x  &+ &y  &- &2z &= &0 \\
        x  &- &y  &  &   &= &-3 \\
        3x &- &y  &- &2z &= &-6  \\
           &  &2y &- &2z &= &3  
      \end{linsys}
   \end{equation*}
   \begin{answer}
     Reduksi berikut menyelesaikan sistem tersebut.
     \begin{align*}
       \begin{amat}{3}
         1 &1   &-2 &0 \\
         1 &-1  &0  &-3  \\
         3 &-1  &-2 &-6 \\
         0 &2   &-2 &3
       \end{amat}
       &\grstep[-3\rho_1+\rho_3]{-\rho_1+\rho_2}
       \begin{amat}{3}
         1 &1   &-2 &0 \\
         0 &-2  &2  &-3  \\
         0 &-4  &4  &-6 \\
         0 &2   &-2 &3
       \end{amat}                                 \\
       &\grstep[\rho_2+\rho_4]{-2\rho_2+\rho_3}
       \begin{amat}{3}
         1 &1   &-2 &0 \\
         0 &-2  &2  &-3  \\
         0 &0   &0  &0 \\
         0 &0   &0  &0
       \end{amat}
     \end{align*}
     Himpunan solusinya adalah sebagai berikut.
     \begin{equation*}
       \set{\colvec{-3/2 \\ 3/2 \\ 0}
            +\colvec{1 \\ 1 \\ 1}z
            \suchthat z\in\Re}
     \end{equation*}
     Dengan cara yang sama, kita dapat mereduksi sistem homogen yang berkaitan
     \begin{align*}
       \begin{amat}{3}
         1 &1   &-2 &0 \\
         1 &-1  &0  &0  \\
         3 &-1  &-2 &0 \\
         0 &2   &-2 &0
       \end{amat}
       &\grstep[-3\rho_1+\rho_3]{-\rho_1+\rho_2}
       \begin{amat}{3}
         1 &1   &-2 &0 \\
         0 &-2  &2  &0  \\
         0 &-4  &4  &0 \\
         0 &2   &-2 &0
       \end{amat}                                         \\
       &\grstep[\rho_2+\rho_4]{-2\rho_2+\rho_3}
       \begin{amat}{3}
         1 &1   &-2 &0 \\
         0 &-2  &2  &0  \\
         0 &0   &0  &0 \\
         0 &0   &0  &0
       \end{amat}
     \end{align*}
     untuk memperoleh himpunan solusinya.
     \begin{equation*}
       \set{
            \colvec{1 \\ 1 \\ 1}z
            \suchthat z\in\Re}
     \end{equation*}
   \end{answer}
  \recommended \item 
    Selesaikan setiap sistem.
    Nyatakan himpunan solusinya menggunakan vektor.
    Tentukan suatu solusi partikular dan himpunan solusi sistem homogen.
    % (These systems also appear in 
    % Exercise~2.\nearbyexercise{exer:SolveInMatrixNotation}.)
    \begin{exparts*}
      \partsitem \( \begin{linsys}[t]{2}
                  3x  &+  &6y  &=  &18  \\
                   x  &+  &2y  &=  &6   
                   \end{linsys}  \)
      \partsitem \( \begin{linsys}[t]{2}
                   x  &+  &y   &=  &1  \\
                   x  &-  &y   &=  &-1   
                    \end{linsys}  \)
      \partsitem \( \begin{linsys}[t]{3}
                   x_1  &   &     &+  &x_3   &=  &4  \\
                   x_1  &-  &x_2  &+  &2x_3  &=  &5  \\
                  4x_1  &-  &x_2  &+  &5x_3  &=  &17  
                    \end{linsys}  \)
      \partsitem \( \begin{linsys}[t]{3}
                   2a   &+  &b    &-  &c     &=  &2  \\
                   2a   &   &     &+  &c     &=  &3  \\
                    a   &-  &b    &   &      &=  &0   
                    \end{linsys}  \)
      \partsitem \( \begin{linsys}[t]{4}
                     x  &+  &2y   &-   &z   &    &    &=  &3  \\
                    2x  &+  &y    &    &    &+   &w   &=  &4  \\
                     x  &-  &y    &+   &z   &+   &w   &=  &1  
                    \end{linsys}  \)
      \partsitem \( \begin{linsys}[t]{4}
                     x  &   &     &+   &z   &+   &w   &=  &4  \\
                    2x  &+  &y    &    &    &-   &w   &=  &2  \\
                    3x  &+  &y    &+   &z   &    &    &=  &7  
                    \end{linsys}  \)
    \end{exparts*}
    \begin{answer}
      Untuk perhitungan aritmetikanya, lihat jawaban pada subbagian sebelumnya.

      \begin{exparts}
        \partsitem
          Berikut adalah himpunan solusinya.
          \begin{equation*}
            S=\set{\colvec[r]{6 \\ 0}+\colvec[r]{-2 \\ 1}y
              \suchthat y\in\Re}
          \end{equation*}
          Berikut adalah solusi partikular dan himpunan solusi sistem homogen
          yang berkaitan.
          \begin{equation*}
            \colvec[r]{6 \\ 0}
              \quad\text{dan}\quad
            \set{\colvec[r]{-2 \\ 1}y
              \suchthat y\in\Re}
          \end{equation*}
          \textit{Catatan.}
          Ada dua hal yang mungkin membingungkan di sini.
          Pertama, himpunan $S$ di atas sama dengan himpunan berikut
          \begin{equation*}
            T=\set{\colvec[r]{4 \\ 1}+\colvec[r]{-2 \\ 1}y
              \suchthat y\in\Re}
          \end{equation*}
          karena kedua himpunan memuat anggota yang sama.
          Semua vektor berikut merupakan jawaban yang benar untuk pertanyaan
          ``Apa salah satu solusi partikularnya?''
          \begin{equation*}
            \colvec{6 \\ 0},\quad
            \colvec{4 \\ 1},\quad
            \colvec{2 \\ 2},\quad
            \colvec{1 \\ 2.5}
          \end{equation*}
          Hal kedua yang mungkin membingungkan adalah bahwa huruf yang kita
          gunakan dalam himpunan tidak berpengaruh.
          Himpunan berikut juga sama dengan $S$.
          \begin{equation*}
            U=\set{\colvec[r]{6 \\ 0}+\colvec[r]{-2 \\ 1}u
                   \suchthat u\in\Re}
          \end{equation*}
        \partsitem
          Berikut adalah himpunan solusinya.
          \begin{equation*}
            \set{\colvec[r]{0 \\ 1} }
          \end{equation*}
          Berikut adalah suatu solusi partikular dan himpunan solusi sistem
          homogen yang berkaitan.
          \begin{equation*}
            \colvec[r]{0 \\ 1}
              \qquad
            \set{\colvec[r]{0 \\ 0} }
          \end{equation*}
        \partsitem
          Himpunan solusinya tak hingga.
          \begin{equation*}
            \set{\colvec[r]{4 \\ -1 \\ 0}+\colvec[r]{-1 \\ 1 \\ 1}x_3
              \suchthat x_3\in\Re}
          \end{equation*}
          Berikut adalah suatu solusi partikular dan himpunan solusi sistem
          homogen yang berkaitan.
          \begin{equation*}
            \colvec[r]{4 \\ -1 \\ 0}
              \qquad
            \set{\colvec[r]{-1 \\ 1 \\ 1}x_3
              \suchthat x_3\in\Re}
          \end{equation*}
        \partsitem
          Himpunan solusinya berelemen tunggal.
          \begin{equation*}
            \set{\colvec[r]{1 \\ 1 \\ 1}}
          \end{equation*}
          Berikut adalah suatu solusi partikular dan himpunan solusi sistem
          homogen yang berkaitan.
          \begin{equation*}
            \colvec[r]{1 \\ 1 \\ 1}
              \qquad
            \set{\colvec[r]{0 \\ 0 \\ 0}}
          \end{equation*}
        \partsitem
          Himpunan solusinya tak hingga.
          \begin{equation*}
            \set{\colvec[r]{5/3 \\ 2/3 \\ 0 \\ 0}
                 +\colvec[r]{-1/3 \\ 2/3 \\ 1 \\ 0}z
                 +\colvec[r]{-2/3 \\ 1/3 \\ 0 \\ 1}w
                 \suchthat z,w\in\Re}
          \end{equation*}
          Berikut adalah suatu solusi partikular dan himpunan solusi sistem
          homogen yang berkaitan.
          \begin{equation*}
            \colvec[r]{5/3 \\ 2/3 \\ 0 \\ 0}
              \qquad
            \set{\colvec[r]{-1/3 \\ 2/3 \\ 1 \\ 0}z
                 +\colvec[r]{-2/3 \\ 1/3 \\ 0 \\ 1}w
                 \suchthat z,w\in\Re}
          \end{equation*}
        \partsitem Himpunan solusi sistem ini kosong.
          Jadi, tidak terdapat solusi partikular.
          Berikut adalah himpunan solusi sistem homogen yang berkaitan.
          \begin{equation*}
            \set{\colvec[r]{-1 \\ 2 \\ 1 \\ 0}z
                 +\colvec[r]{-1 \\ 3 \\ 0 \\ 1}w
                 \suchthat z,w\in\Re}
          \end{equation*}
      \end{exparts}  
    \end{answer}
  \item 
    Selesaikan setiap sistem dan berikan himpunan solusinya dalam notasi vektor.
    Tentukan suatu solusi partikular dan himpunan solusi sistem homogen.
    \begin{exparts*}
      \partsitem \( \begin{linsys}[t]{3}
                  2x  &+  &y  &-  &z  &=  &1  \\
                  4x  &-  &y  &   &   &=  &3  
                  \end{linsys}  \)
      \partsitem \( \begin{linsys}[t]{4}
                   x  &   &   &-  &z  &   &   &=  &1  \\
                      &   &y  &+  &2z &-  &w  &=  &3  \\
                   x  &+  &2y &+  &3z &-  &w  &=  &7  
                   \end{linsys}  \)
      \partsitem \( \begin{linsys}[t]{4}
                   x  &-  &y  &+  &z  &   &   &=  &0  \\
                      &   &y  &   &   &+  &w  &=  &0  \\
                  3x  &-  &2y &+  &3z &+  &w  &=  &0  \\
                      &   &-y &   &   &-  &w  &=  &0  
                  \end{linsys}  \)
      \partsitem \( \begin{linsys}[t]{5}
                   a  &+  &2b &+  &3c &+  &d  &-  &e  &=  &1  \\
                  3a  &-  &b  &+  &c  &+  &d  &+  &e  &=  &3  
                  \end{linsys}  \)
    \end{exparts*}
    \begin{answer}
    Jawaban pada subbagian sebelumnya memperlihatkan operasi barisnya.
    Setiap jawaban di sini hanya mencantumkan himpunan solusi, solusi partikular,
    dan solusi homogen.
    \begin{exparts}
      \partsitem
        Berikut adalah himpunan solusinya.
        \begin{equation*}
          \set{\colvec[r]{2/3 \\ -1/3 \\ 0}
               +\colvec[r]{1/6 \\ 2/3 \\ 1}z
              \suchthat z\in\Re}
        \end{equation*}
        Berikut adalah suatu solusi partikular dan himpunan solusi sistem
        homogen yang berkaitan.
        \begin{equation*}
          \colvec[r]{2/3 \\ -1/3 \\ 0}
            \qquad
        \set{\colvec[r]{1/6 \\ 2/3 \\ 1}z
            \suchthat z\in\Re}
        \end{equation*}
      \partsitem
        Himpunan solusinya tak hingga.
        \begin{equation*}
          \set{\colvec[r]{1 \\ 3 \\ 0 \\ 0}
               +\colvec[r]{1 \\-2 \\ 1 \\ 0}z
              \suchthat z\in\Re}
        \end{equation*}
        Berikut adalah suatu solusi partikular dan himpunan solusi sistem
        homogen yang berkaitan.
        \begin{equation*}
          \colvec[r]{1 \\ 3 \\ 0 \\ 0}
            \qquad
          \set{\colvec[r]{1 \\ -2 \\ 1 \\ 0}z
              \suchthat z\in\Re}
        \end{equation*}
      \partsitem
        Berikut adalah himpunan solusinya.
        \begin{equation*}
          \set{\colvec[r]{0 \\ 0 \\ 0 \\ 0}
               +\colvec[r]{-1 \\ 0 \\ 1 \\ 0}z
               +\colvec[r]{-1 \\ -1 \\ 0 \\ 1}w
              \suchthat z,w\in\Re}
        \end{equation*}
        Berikut adalah suatu solusi partikular dan himpunan solusi sistem
        homogen yang berkaitan.
        \begin{equation*}
          \colvec[r]{0 \\ 0 \\ 0 \\ 0}
            \qquad
          \set{\colvec[r]{-1 \\ 0 \\ 1 \\ 0}z
               +\colvec[r]{-1 \\ -1 \\ 0 \\ 1}w
              \suchthat z,w\in\Re}
        \end{equation*}
      \partsitem
        Berikut adalah himpunan solusinya.
        \begin{equation*}
          \set{\colvec[r]{1 \\ 0 \\ 0 \\ 0 \\ 0}
               +\colvec[r]{-5/7 \\ -8/7 \\ 1 \\ 0 \\ 0}c
               +\colvec[r]{-3/7 \\ -2/7 \\ 0 \\ 1 \\ 0}d
               +\colvec[r]{-1/7 \\ 4/7 \\ 0 \\ 0 \\ 1}e
              \suchthat c,d,e\in\Re}
        \end{equation*}
        Berikut adalah suatu solusi partikular dan himpunan solusi sistem
        homogen yang berkaitan.
        \begin{equation*}
          \colvec[r]{1 \\ 0 \\ 0 \\ 0 \\ 0}
            \qquad
          \set{\colvec[r]{-5/7 \\ -8/7 \\ 1 \\ 0 \\ 0}c
               +\colvec[r]{-3/7 \\ -2/7 \\ 0 \\ 1 \\ 0}d
               +\colvec[r]{-1/7 \\ 4/7 \\ 0 \\ 0 \\ 1}e
              \suchthat c,d,e\in\Re}
        \end{equation*}
    \end{exparts}  
   \end{answer}
  \recommended \item 
    Untuk sistem
    \begin{equation*}
      \begin{linsys}{4}
       2x  &-  &y  &   &    &-  &w  &=  &3  \\
           &   &y  &+  &z   &+  &2w &=  &2  \\
        x  &-  &2y &-  &z   &   &   &=  &-1
      \end{linsys}
    \end{equation*}
    manakah di antara vektor berikut yang dapat digunakan sebagai bagian
    solusi partikular dari suatu solusi umum?
    \begin{exparts*}
      \partsitem   \( \colvec[r]{0 \\ -3 \\ 5 \\ 0} \)
      \partsitem   \( \colvec[r]{2 \\ 1 \\ 1 \\ 0} \)
      \partsitem   \( \colvec[r]{-1 \\ -4 \\ 8 \\ -1} \)
    \end{exparts*}
    \begin{answer}
      Substitusikan saja setiap vektor dan periksa apakah ketiga persamaan
      terpenuhi.
      \begin{exparts}
        \partsitem Tidak.
        \partsitem Ya.
        \partsitem Ya.
      \end{exparts}  
    \end{answer}
  \recommended \item  
    \nearbylemma{th:GenEqPartHomo} menyatakan bahwa kita dapat menggunakan
    sembarang solusi partikular sebagai $\vec{p}$.
    Jika mungkin, carilah solusi umum sistem berikut
    \begin{equation*}
      \begin{linsys}{4}
        x  &-  &y  &   &    &+  &w  &=  &4  \\
       2x  &+  &3y &-  &z   &   &   &=  &0  \\
           &   &y  &+  &z   &+  &w  &=  &4  
      \end{linsys}
    \end{equation*}
    yang menggunakan vektor yang diberikan sebagai solusi partikularnya.
    \begin{exparts*}
      \partsitem   \( \colvec[r]{0 \\ 0 \\ 0 \\ 4} \)
      \partsitem   \( \colvec[r]{-5 \\ 1 \\ -7 \\ 10} \)
      \partsitem   \( \colvec[r]{2 \\ -1 \\ 1 \\ 1} \)
    \end{exparts*}
    \begin{answer}
      Penerapan Metode Gauss pada sistem homogen yang berkaitan
      \begin{align*}
        \begin{amat}[r]{4}
           1  &-1  &0  &1  &0  \\
           2  &3   &-1 &0  &0  \\
           0  &1   &1  &1  &0
        \end{amat}
        &\grstep{-2\rho_1+\rho_2}
        \begin{amat}[r]{4}
           1  &-1  &0  &1  &0  \\
           0  &5   &-1 &-2 &0  \\
           0  &1   &1  &1  &0
        \end{amat}                                \\
        &\grstep{-(1/5)\rho_2+\rho_3}
        \begin{amat}[r]{4}
           1  &-1  &0  &1  &0  \\
           0  &5   &-1 &-2 &0  \\
           0  &0   &6/5&7/5&0
        \end{amat}
      \end{align*}
      menghasilkan himpunan solusi sistem homogen berikut.
      \begin{equation*}
        \set{\colvec[r]{-5/6 \\ 1/6 \\ -7/6 \\ 1}w\suchthat w\in\Re}
      \end{equation*}
      \begin{exparts}
        \partsitem Vektor tersebut memang merupakan solusi partikular, sehingga
          solusi umum yang diminta adalah sebagai berikut.
          \begin{equation*}
            \set{\colvec[r]{0 \\ 0 \\ 0 \\ 4}+
                 \colvec[r]{-5/6 \\ 1/6 \\ -7/6 \\ 1}w\suchthat w\in\Re}
          \end{equation*}
        \partsitem Vektor tersebut merupakan solusi partikular, sehingga solusi
          umum yang diminta adalah sebagai berikut.
          \begin{equation*}
            \set{\colvec[r]{-5 \\ 1 \\ -7 \\ 10}+
                 \colvec[r]{-5/6 \\ 1/6 \\ -7/6 \\ 1}w\suchthat w\in\Re}
          \end{equation*}
        \partsitem Vektor tersebut bukan solusi sistem karena tidak memenuhi
          persamaan ketiga.
          Tidak ada solusi umum yang demikian.
      \end{exparts} 
    \end{answer}
  \item 
     Salah satu matriks berikut nonsingular, sedangkan yang lain singular.
     Tentukan masing-masing.
     \begin{exparts*}
       \partsitem $\begin{mat}[r]
           1  &3   \\
           4  &-12    
         \end{mat}$
       \partsitem $\begin{mat}[r]
           1  &3  \\
           4  &12  
         \end{mat}$
     \end{exparts*}
     \begin{answer}
       Matriks pertama nonsingular, sedangkan matriks kedua singular.
       Terapkan Metode Gauss dan periksa apakah hasil berbentuk eselon memiliki
       bilangan yang tidak sama dengan $0$ pada setiap entri diagonal.
     \end{answer}
  \recommended \item 
    Singular atau nonsingular?
    \begin{exparts*}
      \partsitem \(
        \begin{mat}[r]
          1  &2  \\
          1  &3
        \end{mat}   \)
      \partsitem \(
        \begin{mat}[r]
          1  &2  \\
         -3  &-6
        \end{mat}   \)
      \partsitem \(
        \begin{mat}[r]
          1  &2  &1  \\
          1  &3  &1
        \end{mat}   \)
      \partsitem \(
        \begin{mat}[r]
          1  &2  &1  \\
          1  &1  &3  \\
          3  &4  &7
        \end{mat}   \)
      \partsitem \(
        \begin{mat}[r]
          2  &2  &1  \\
          1  &0  &5  \\
         -1  &1  &4
        \end{mat}   \)
    \end{exparts*}
    \begin{answer}
      \begin{exparts}
      \partsitem Nonsingular:
        \begin{equation*}
          \grstep{-\rho_1+\rho_2}
          \begin{mat}[r]
            1  &2  \\
            0  &1
          \end{mat}
        \end{equation*}
        berakhir dengan setiap baris memuat satu entri utama.
      \partsitem Singular:
        \begin{equation*}
          \grstep{3\rho_1+\rho_2}
          \begin{mat}[r]
            1  &2  \\
            0  &0
          \end{mat}
        \end{equation*}
        berakhir dengan baris \( 2 \) tanpa entri utama.
      \partsitem Bukan keduanya.
        Matriks harus persegi agar salah satu istilah tersebut berlaku.
      \partsitem Singular.
      \partsitem Nonsingular.
     \end{exparts}  
    \end{answer}
  \recommended \item 
    Apakah vektor yang diberikan berada dalam himpunan yang dibangkitkan oleh
    himpunan yang diberikan?
      \begin{exparts}
        \partsitem \( \colvec[r]{2 \\ 3}, \)\
          \( \set{\colvec[r]{1 \\ 4},
                \colvec[r]{1 \\ 5}} \)
        \partsitem \( \colvec[r]{-1 \\ 0 \\ 1}, \)\
          \( \set{\colvec[r]{2 \\ 1 \\ 0},
                \colvec[r]{1 \\ 0 \\ 1}} \)
        \partsitem \( \colvec[r]{1 \\ 3 \\ 0}, \)\
          \( \set{\colvec[r]{1 \\ 0 \\ 4},
                \colvec[r]{2 \\ 1 \\ 5},
                \colvec[r]{3 \\ 3 \\ 0},
                \colvec[r]{4 \\ 2 \\ 1}} \)
        \partsitem \( \colvec[r]{1 \\ 0 \\ 1 \\ 1}, \)\
          \( \set{\colvec[r]{2 \\ 1 \\ 0 \\ 1},
                \colvec[r]{3 \\ 0 \\ 0 \\ 2}} \)
      \end{exparts}
      \begin{answer}
        Dalam setiap kasus, kita harus menentukan apakah vektor tersebut
        merupakan kombinasi linear dari vektor-vektor dalam himpunan.
        \begin{exparts}
          \partsitem Ya.
            Selesaikan
            \begin{equation*}
              c_1\colvec[r]{1 \\ 4}+c_2\colvec[r]{1 \\ 5}=\colvec[r]{2 \\ 3}
            \end{equation*}
            dengan
            \begin{equation*}
              \begin{amat}[r]{2}
                1  &1  &2  \\
                4  &5  &3
              \end{amat}
              \grstep{-4\rho_1+\rho_2}
              \begin{amat}[r]{2}
                1  &1  &2  \\
                0  &1  &-5
              \end{amat}
            \end{equation*}
            untuk menyimpulkan bahwa terdapat $c_1$ dan $c_2$ yang menghasilkan
            kombinasi tersebut.
          \partsitem Tidak.
            Reduksi
            \begin{equation*}
              \begin{amat}[r]{2}
                2  &1  &-1 \\
                1  &0  &0  \\
                0  &1  &1
              \end{amat}
              \grstep{-(1/2)\rho_1+\rho_2}
              \begin{amat}[r]{2}
                2  &1     &-1 \\
                0  &-1/2  &1/2  \\
                0  &1     &1
              \end{amat}
              \grstep{2\rho_2+\rho_3}
              \begin{amat}[r]{2}
                2  &1     &-1 \\
                0  &-1/2  &1/2  \\
                0  &0     &2
              \end{amat}
            \end{equation*}
            menunjukkan bahwa
            \begin{equation*}
              c_1\colvec[r]{2 \\ 1 \\ 0}+c_2\colvec[r]{1 \\ 0 \\ 1}
                =\colvec[r]{-1 \\ 0 \\ 1}
            \end{equation*}
            tidak memiliki solusi.
          \partsitem Ya.
            Reduksi
            \begin{align*}
              \begin{amat}[r]{4}
                1  &2  &3  &4  &1  \\
                0  &1  &3  &2  &3  \\
                4  &5  &0  &1  &0
              \end{amat}
              &\grstep{-4\rho_1+\rho_3}
              \begin{amat}[r]{4}
                1  &2  &3  &4  &1  \\
                0  &1  &3  &2  &3  \\
                0  &-3 &-12&-15&-4
              \end{amat}                   \\
              &\grstep{3\rho_2+\rho_3}
              \begin{amat}[r]{4}
                1  &2  &3  &4  &1  \\
                0  &1  &3  &2  &3  \\
                0  &0  &-3 &-9 &5
              \end{amat}
            \end{align*}
            menunjukkan bahwa terdapat tak hingga banyak cara
            \begin{equation*}
              \set{\colvec[r]{c_1 \\ c_2 \\ c_3 \\ c_4}=
                   \colvec[r]{-10 \\ 8 \\ -5/3 \\ 0}+
                   \colvec[r]{-9 \\ 7 \\ -3 \\ 1}c_4
                    \suchthat c_4\in\Re}
            \end{equation*}
            untuk menuliskan kombinasi.
            \begin{equation*}
              \colvec[r]{1 \\ 3 \\ 0}=
              c_1\colvec[r]{1 \\ 0 \\ 4}+
              c_2\colvec[r]{2 \\ 1 \\ 5}+
              c_3\colvec[r]{3 \\ 3 \\ 0}+
              c_4\colvec[r]{4 \\ 2 \\ 1}
            \end{equation*}
          \partsitem Tidak.
            Perhatikan komponen ketiganya.
        \end{exparts} 
      \end{answer}
  \item 
     Buktikan bahwa setiap sistem linear dengan matriks koefisien nonsingular
     memiliki solusi, dan bahwa solusi tersebut tunggal.
    \begin{answer}
        Karena matriks koefisiennya nonsingular, Metode Gauss berakhir pada
        suatu bentuk eselon yang setiap variabelnya merupakan variabel utama
        dalam suatu persamaan.
        Substitusi balik menghasilkan solusi tunggal.

      (Cara lain untuk melihat bahwa solusi tersebut tunggal adalah dengan
      memperhatikan bahwa, menurut definisi, jika matriks koefisiennya
      nonsingular maka sistem homogen yang berkaitan memiliki solusi tunggal.
      Karena solusi umum merupakan jumlah suatu solusi partikular dengan
      setiap solusi homogen, solusi umum memiliki paling banyak satu elemen.)
     \end{answer}
  \item 
    Dalam pembuktian \nearbylemma{le:HomoSltnSpanVecs}, apa yang terjadi jika
    tidak ada persamaan selain \( 0=0 \)?
    \begin{answer}
      Dalam kasus ini, himpunan solusinya adalah seluruh \( \Re^n \), dan
      himpunan tersebut dapat kita nyatakan dalam bentuk yang diminta.
      \begin{equation*}
        \set{c_1\colvec[r]{1 \\ 0 \\ \vdotswithin{1} \\ 0}
             +c_2\colvec[r]{0 \\ 1 \\ \vdotswithin{0} \\ 0}
             +\cdots
             +c_n\colvec[r]{0 \\ 0 \\ \vdotswithin{0} \\ 1}
             \suchthat c_1,\ldots,c_n\in\Re}
      \end{equation*}  
     \end{answer}
  \recommended \item 
    Buktikan bahwa jika \( \vec{s} \) dan \( \vec{t} \) memenuhi suatu sistem
    homogen, maka vektor-vektor berikut juga memenuhinya.
    \begin{exparts*}
      \partsitem \( \vec{s}+\vec{t} \)
      \partsitem \( 3\vec{s} \)
      \partsitem \( k\vec{s}+m\vec{t} \) untuk \( k,m\in\Re \)
    \end{exparts*}
    Apa yang salah dengan argumen berikut: ``Ketiga hasil ini menunjukkan bahwa
    jika suatu sistem homogen memiliki satu solusi, maka sistem itu memiliki
    banyak solusi\Dash setiap kelipatan suatu solusi merupakan solusi lain,
    dan setiap jumlah solusi juga merupakan solusi\Dash sehingga tidak ada
    sistem homogen yang memiliki tepat satu solusi.''?
    \begin{answer}
      Andaikan \( \vec{s},\vec{t}\in\Re^n \) dan tuliskan keduanya sebagai berikut.
      \begin{equation*}
        \vec{s}=\colvec{s_1 \\ \vdotswithin{s_1} \\ s_n}
          \qquad
        \vec{t}=\colvec{t_1 \\ \vdotswithin{t_1} \\ t_n}
      \end{equation*}
      Misalkan pula \( a_{i,1}x_1+\cdots+a_{i,n}x_n=0 \) adalah persamaan ke-\( i \)
      dalam sistem homogen tersebut.
      \begin{exparts}
        \partsitem Pemeriksaannya mudah.
          \begin{multline*}
            a_{i,1}(s_1+t_1)+\cdots+a_{i,n}(s_n+t_n)      \\
            =
            (a_{i,1}s_1+\cdots+a_{i,n}s_n)
            +(a_{i,1}t_1+\cdots+a_{i,n}t_n)           
            =
            0+0
          \end{multline*}
        \partsitem Kasus ini serupa dengan kasus sebelumnya.
          \begin{equation*}
            a_{i,1}(3s_1)+\cdots+a_{i,n}(3s_n)
            =3(a_{i,1}s_1+\cdots+a_{i,n}s_n)
            =3\cdot 0=0
          \end{equation*}
        \partsitem Kasus ini tidak jauh lebih sulit.
          \begin{multline*}
            a_{i,1}(ks_1+mt_1)+\cdots+a_{i,n}(ks_n+mt_n)  \\
            =
            k(a_{i,1}s_1+\cdots+a_{i,n}s_n)
            +m(a_{i,1}t_1+\cdots+a_{i,n}t_n)         
            =
            k\cdot 0+m\cdot 0
          \end{multline*}
      \end{exparts}  
     Kekeliruan dalam argumen tersebut adalah bahwa setiap kombinasi linear
     yang hanya melibatkan vektor nol tetap menghasilkan vektor nol.
   \end{answer}
  \item
    Buktikan bahwa jika suatu sistem yang seluruh koefisien dan konstantanya
    rasional memiliki solusi, maka sistem itu memiliki setidaknya satu solusi
    yang seluruh komponennya rasional.
    Apakah sistem tersebut harus memiliki tak hingga banyak solusi demikian?
    \begin{answer}
      Pertama, pembuktiannya.

      Jika semua koefisien sistem rasional, maka Metode Gauss beserta substitusi
      balik hanya akan menggunakan bilangan rasional (misalnya,
      \( -(m/n)\rho_i+\rho_j \)).
      Jadi, kita dapat menyatakan himpunan solusi dengan hanya menggunakan
      bilangan rasional sebagai komponen setiap vektor.
      Oleh karena itu, seluruh komponen solusi partikular bersifat rasional.

      Mengenai pertanyaan tentang tak hingga banyak solusi, terdapat tak hingga
      banyak solusi vektor rasional jika dan hanya jika sistem homogen yang
      berkaitan memiliki tak hingga banyak solusi vektor riil.
      Hal ini karena menetapkan sembarang nilai rasional pada parameter akan
      menghasilkan solusi yang seluruh komponennya rasional.
   \end{answer}
\end{exercises}

\endinput




\subsection{Membandingkan Deskripsi Himpunan}
\emph{Subbagian ini bersifat opsional.
Materi selanjutnya tidak memerlukan pembahasan di sini.}

Suatu himpunan dapat dideskripsikan dengan banyak cara berbeda.
Berikut adalah dua deskripsi berbeda untuk satu himpunan yang sama:
\begin{equation*}
  \set{\colvec[r]{1 \\ 2 \\ 3}z\suchthat z\in\Re}
  \quad\text{dan}\quad
  \set{\colvec[r]{2 \\ 4 \\ 6}w\suchthat w\in\Re}.
\end{equation*}
Sebagai contoh, himpunan ini memuat
\begin{equation*}
  \colvec[r]{5 \\ 10 \\ 15}
\end{equation*}
(ambil $z=5$ dan $w=5/2$), tetapi tidak memuat
\begin{equation*}
  \colvec[r]{4 \\ 8 \\ 11}
\end{equation*}
(komponen pertama menghasilkan $z=4$, tetapi ini bertentangan dengan komponen
ketiga; demikian pula, komponen pertama menghasilkan $w=2$, tetapi komponen
ketiga menghasilkan nilai yang berbeda).
Berikut adalah deskripsi ketiga untuk himpunan yang sama:
\begin{equation*}
  \set{\colvec[r]{3 \\ 6 \\ 9}+\colvec[r]{-1 \\ -2 \\ -3}y\suchthat y\in\Re}.
\end{equation*}

Kita perlu menentukan kapan dua deskripsi menjelaskan himpunan yang sama.
Secara lebih praktis, bagaimana seseorang dapat mengetahui bahwa jawaban suatu
soal pekerjaan rumah mendeskripsikan himpunan yang sama dengan himpunan yang
dideskripsikan pada bagian jawaban di belakang buku?

Dua himpunan sama jika dan hanya jika keduanya memiliki anggota yang sama.
Cara umum untuk menunjukkan bahwa dua himpunan, $S_1$ dan $S_2$, sama adalah
dengan menunjukkan inklusi timbal balik:\index{himpunan!inklusi timbal balik}\index{inklusi timbal balik}
setiap anggota $S_1$ juga berada dalam $S_2$, dan setiap anggota $S_2$ juga
berada dalam $S_1$.\appendrefs{kesamaan himpunan}

\begin{example}
Untuk menunjukkan bahwa
\begin{equation*}
  S_1=
  \set{\colvec[r]{1 \\ -1 \\ 0}c+\colvec[r]{1 \\ 1 \\ 0}d\suchthat c,d\in\Re}
\end{equation*}
sama dengan
\begin{equation*}
  S_2=
  \set{\colvec[r]{4 \\ 1 \\ 0}m+\colvec[r]{-1 \\ -3 \\ 0}n\suchthat m,n\in\Re}
\end{equation*}
pertama-tama kita tunjukkan bahwa $S_1\subseteq S_2$, lalu bahwa
$S_2\subseteq S_1$.

Untuk bagian pertama, kita harus memeriksa bahwa setiap vektor dari
\( S_1 \) juga berada dalam \( S_2 \).
Kita terlebih dahulu meninjau dua contoh untuk dijadikan model bagi argumen umum.
Jika kita membentuk suatu anggota $S_1$ dengan mencoba \( c=1 \) dan \( d=1 \),
maka untuk menunjukkan bahwa anggota itu berada dalam $S_2$, kita memerlukan
\( m \) dan $n$ sedemikian sehingga
\begin{equation*}
  \colvec[r]{4 \\ 1 \\ 0}m
  +\colvec[r]{-1 \\ -3 \\ 0}n
  =\colvec[r]{2 \\ 0 \\ 0}
\end{equation*}
yakni, hubungan berikut berlaku antara $m$ dan $n$.
\begin{equation*}
  \begin{linsys}{2}
    4m  &-  &n  &=  &2  \\
    1m  &-  &3n &=  &0  \\
        &   &0  &=  &0 
  \end{linsys}  
\end{equation*}
Dengan cara yang sama, jika kita mencoba \( c=2 \) dan \( d=-1 \), maka untuk
menunjukkan bahwa anggota $S_1$ yang dihasilkan berada dalam $S_2$, kita
memerlukan \( m \) dan $n$ sedemikian sehingga
\begin{equation*}
  \colvec[r]{4 \\ 1 \\ 0}m
  +\colvec[r]{-1 \\ -3 \\ 0}n
  =\colvec[r]{1 \\ -3 \\ 0}
\end{equation*}
yakni, hubungan berikut berlaku.
\begin{equation*}
  \begin{linsys}{2}
    4m  &-  &n  &=  &1  \\
    1m  &-  &3n &=  &-3 \\
        &   &0  &=  &0 
   \end{linsys}
\end{equation*}
Dalam kasus umum, untuk menunjukkan bahwa setiap vektor dari \( S_1 \)
merupakan anggota \( S_2 \), kita harus menunjukkan bahwa untuk sembarang
\( c \) dan \( d \) terdapat \( m \) dan \( n \) yang sesuai.
Dengan mengikuti pola pada contoh, tetapkan
\begin{equation*}
  \colvec{c+d \\ -c+d \\ 0}\in S_1
\end{equation*}
dan carilah \( m \) dan \( n \) sedemikian sehingga
\begin{equation*}
  \colvec[r]{4 \\ 1 \\ 0}m
  +\colvec[r]{-1 \\ -3 \\ 0}n
  =\colvec{c+d \\ -c+d \\ 0}
\end{equation*}
yakni, hubungan berikut berlaku.
\begin{equation*}
  \begin{linsys}{2}
    4m  &-  &n  &=  &c+d\hfill  \\
     m  &-  &3n &=  &-c+d\hfill  \\
        &   &0  &=  &0\hfill  
  \end{linsys}
\end{equation*}
Penerapan Metode Gauss
\begin{equation*}
  \begin{linsys}{2}
    4m  &-  &n  &=  &c+d\hfill  \\
     m  &-  &3n &=  &-c+d\hfill  
  \end{linsys}
  \grstep{-(1/4)\rho_1+\rho_2}
  \begin{linsys}{2}
    4m  &-  &n        &=  &c+d\hfill            \\
        &   &-(11/4)n &=  &-(5/4)c+(3/4)d\hfill  
   \end{linsys}
\end{equation*}
menghasilkan \( n=(5/11)c-(3/11)d \) dan \( m=(4/11)c+(2/11)d \).
Hal ini menunjukkan bahwa untuk setiap pilihan $c$ dan $d$ terdapat $m$ dan
$n$ yang sesuai.
Kita menyimpulkan bahwa setiap anggota $S_1$ merupakan anggota $S_2$ karena
anggota itu dapat ditulis ulang sebagai berikut:
\begin{equation*}
   \colvec{c+d \\ -c+d \\ 0}
   =\colvec[r]{4 \\ 1 \\ 0}((4/11)c+(2/11)d)+
   \colvec[r]{-1 \\ -3 \\ 0}((5/11)c-(3/11)d).
\end{equation*}

Untuk inklusi sebaliknya, \( S_2\subseteq S_1 \), kita melakukan kebalikannya.
Kita ingin menunjukkan bahwa untuk setiap pilihan $m$ dan $n$ terdapat
$c$ dan $d$ yang sesuai.
Jadi, tetapkan $m$ dan $n$, lalu selesaikan \( c \) dan \( d \):
\begin{equation*}
   \begin{linsys}{2}
     c  &+ &d  &= &4m-n\hfill \\
    -c  &+ &d  &= &m-3n\hfill 
   \end{linsys}
   \grstep{\rho_1+\rho_2}
   \begin{linsys}{2}
     c  &+ &d  &= &4m-n\hfill \\
        &  &2d &= &5m-4n\hfill 
    \end{linsys}
\end{equation*}
menunjukkan bahwa \( d=(5/2)m-2n \) dan \( c=(3/2)m+n \).
Jadi, setiap vektor dari \( S_2 \)
\begin{equation*}
  \colvec[r]{4 \\ 1 \\ 0}m+\colvec[r]{-1 \\ -3 \\ 0}n
\end{equation*}
juga memiliki bentuk yang tepat untuk \( S_1 \),
\begin{equation*}
  \colvec[r]{1 \\ -1 \\ 0}((3/2)m+n)
    +\colvec[r]{1 \\ 1 \\ 0}((5/2)m-2n).
\end{equation*}
\end{example}

\begin{example}
Tentu saja, kadang-kadang dua himpunan tidak sama.
Metode pada contoh sebelumnya akan membantu kita melihat hubungan antara
kedua himpunan tersebut.
Himpunan-himpunan berikut
\begin{equation*}
  P=
  \set{\colvec{x+y \\ 2x \\ y}\suchthat x,y\in\Re}
  \quad\text{dan}\quad
  R=
  \set{\colvec{m+p \\ n \\ p}\suchthat m,n,p\in\Re}
\end{equation*}
tidak sama.
Meskipun $P$ merupakan himpunan bagian dari $R$, himpunan itu adalah himpunan
bagian sejati dari $R$ karena $R$ bukan himpunan bagian dari $P$.

Untuk melihat hal ini, pertama-tama perhatikan bahwa suatu vektor dari \( P \)
dapat kita nyatakan dalam bentuk yang digunakan untuk \( R \)\Dash jika kita menetapkan
$x$ dan $y$, kita dapat menentukan $m$, $n$, dan $p$ yang sesuai:
\begin{equation*}
  \begin{linsys}{3}
     m  &   &   &+  &p  &=  &x+y\hfill  \\
        &   &n  &   &   &=  &2x\hfill   \\
        &   &   &   &p  &=  &y\hfill    
  \end{linsys}
\end{equation*}
menunjukkan bahwa kita dapat menyatakan sembarang
\begin{equation*}
  \vec{v}=
  \colvec[r]{1 \\ 2 \\ 0}x+
  \colvec[r]{1 \\ 0 \\ 1}y
\end{equation*}
sebagai anggota \( R \) dengan
\( m=x \), \( n=2x \), dan \( p=y \):
\begin{equation*}
  \vec{v}=
  \colvec[r]{1 \\ 0 \\ 0}x+
  \colvec[r]{0 \\ 1 \\ 0}2x+
  \colvec[r]{1 \\ 0 \\ 1}y.
\end{equation*}
Jadi, \( P\subseteq R \).

Namun, untuk arah sebaliknya, reduksi yang diperoleh dengan menetapkan
$m$, $n$, dan $p$, lalu mencari $x$ dan $y$,
\begin{align*}
  \begin{linsys}{2}
     x  &+  &y  &=  &m+p\hfill  \\
    2x  &   &   &=  &n\hfill    \\
        &   &y  &=  &p\hfill    
  \end{linsys}
  &\grstep{-2\rho_1+\rho_2}
  \begin{linsys}{2}
     x  &+  &y  &=  &m+p\hfill  \\
        &   &-2y&=  &-2m+n-2p\hfill \\
        &   &y  &=  &p\hfill    
   \end{linsys}                                  \\
  &\grstep{(1/2)\rho_2+\rho_3}
  \begin{linsys}{2}
     x  &+  &y  &=  &m+p\hfill  \\
        &   &-2y&=  &-2m+n-2p\hfill \\
     &   &0  &=  &-m+(1/2)n\hfill
    \end{linsys}
\end{align*}
menunjukkan bahwa satu-satunya vektor
\begin{equation*}
  \colvec{m+p \\ n \\ p}\in R
\end{equation*}
yang dapat dinyatakan dalam bentuk
\begin{equation*}
  \colvec{x+y \\ 2x \\ y}
\end{equation*}
adalah vektor yang memenuhi \( 0=-m+(1/2)n \).
Sebagai contoh,
\begin{equation*}
  \colvec[r]{0 \\ 1 \\ 0}
\end{equation*}
berada dalam \( R \), tetapi tidak dalam \( P \).
\end{example}

\begin{exercises}
  \item 
    Tentukan apakah vektor tersebut merupakan anggota himpunan.
    \begin{exparts}
      \partsitem $\colvec[r]{2 \\ 3}$, 
         $\set{\colvec[r]{1 \\ 2}k\suchthat k\in\Re}$
      \partsitem $\colvec[r]{-3 \\ 3}$, 
         $\set{\colvec[r]{1 \\ -1}k\suchthat k\in\Re}$
      \partsitem $\colvec[r]{-3 \\ 3 \\ 4}$, 
             $\set{\colvec[r]{1 \\ -1 \\ 2}k\suchthat k\in\Re}$
      \partsitem $\colvec[r]{-3 \\ 3 \\ 4}$, 
             $\set{\colvec[r]{1 \\ -1 \\ 2}k+\colvec[r]{0 \\ 0 \\ 2}m
                \suchthat k,m\in\Re}$
      \partsitem $\colvec[r]{1 \\ 4 \\ 14}$, 
             $\set{\colvec[r]{2 \\ 2 \\ 5}k+\colvec[r]{-1 \\ 0 \\ 2}m
                \suchthat k,m\in\Re}$
      \partsitem $\colvec[r]{1 \\ 4 \\ 6}$, 
             $\set{\colvec[r]{2 \\ 2 \\ 5}k+\colvec[r]{-1 \\ 0 \\ 2}m
                \suchthat k,m\in\Re}$
    \end{exparts}
    \begin{answer}
      \begin{exparts}
        \partsitem Tidak.
        \partsitem Ya.
        \partsitem Tidak.
        \partsitem Ya.
        \partsitem Tidak; dua komponen pertama mengharuskan $k=2$ dan $m=3$,
          tetapi komponen ketiga kemudian bernilai 16, bukan 14.
        \partsitem Tidak; gunakan Metode Gauss untuk menyimpulkan bahwa sistem
          tersebut tidak memiliki solusi.
      \end{exparts}
    \end{answer}
  \item 
     Berikan dua deskripsi lain untuk himpunan ini. 
     \begin{equation*}
       \set{\colvec[r]{2 \\ -5}k\suchthat k\in\Re}
     \end{equation*}
     \begin{answer}
       Salah satu cara sederhana adalah mengalikan vektor tersebut dengan dua
       dan dengan tiga:
       \begin{equation*}
         \set{\colvec[r]{4 \\ -10}k\suchthat k\in\Re}
         \quad\text{dan}\quad
         \set{\colvec[r]{6 \\ -15}k\suchthat k\in\Re}.
       \end{equation*}
      \end{answer}
  \recommended \item 
    Tunjukkan bahwa ketiga deskripsi pada awal subbagian ini semuanya
    mendeskripsikan himpunan yang sama.
    \begin{answer}
      Alih-alih menunjukkan ketiga kesamaan secara terpisah, kita dapat
      menunjukkan bahwa himpunan pertama sama dengan himpunan kedua dan bahwa
      himpunan kedua sama dengan himpunan ketiga.
      Kedua kesamaan itu mudah ditunjukkan dengan metode dalam subbagian ini.
    \end{answer}
  \recommended \item 
    Tunjukkan bahwa kedua himpunan berikut sama
    \begin{equation*}
      \set{\colvec[r]{1 \\ 4 \\ 1 \\ 1}
           +\colvec[r]{-1 \\ 0 \\ 1 \\ 0}z\suchthat z\in\Re  }
      \quad\text{dan}\quad
      \set{\colvec[r]{0 \\ 4 \\ 2 \\ 1}
           +\colvec[r]{-1 \\ 0 \\ 1 \\ 0}k\suchthat k\in\Re  },
    \end{equation*}
    dan bahwa keduanya mendeskripsikan himpunan solusi sistem berikut.
    \begin{equation*}  
       \begin{linsys}{4}
         x  &-  &y  &+  &z  &+  &w  &=  &-1  \\
            &   &y  &   &   &-  &w  &=  &3   \\
         x  &   &   &+  &z  &+  &2w &=  &4
       \end{linsys}
    \end{equation*}
    \begin{answer}
      Sistem tersebut direduksi sebagai berikut:
      \begin{align*}
         &\grstep{-\rho_1+\rho_3}
         \begin{linsys}{4}
           x  &-  &y  &+  &z  &+  &w  &=  &-1  \\
              &   &y  &   &   &-  &w  &=  &3   \\
              &   &y  &   &   &+  &w  &=  &5   
           \end{linsys}                              \\
         &\grstep{-\rho_2+\rho_3}
         \begin{linsys}{4}
           x  &-  &y  &+  &z  &+  &w  &=  &-1  \\
              &   &y  &   &   &-  &w  &=  &3   \\
              &   &   &   &   &   &2w &=  &2   
          \end{linsys}
      \end{align*}
      Ini menunjukkan bahwa \( w=1 \), \( y=4 \), dan \( x=2-z \).   
    \end{answer}
  \recommended \item 
    Tentukan apakah kedua himpunan dalam setiap bagian sama.
    \begin{exparts}
      \partsitem \( \set{\colvec[r]{1 \\ 2}
                        +\colvec[r]{0 \\ 3}t
                     \suchthat t\in\Re} \)
            dan
            \( \set{\colvec[r]{1 \\ 8}
                        +\colvec[r]{0 \\ -1}s
                     \suchthat s\in\Re} \)
      \partsitem \( \set{\colvec[r]{1 \\ 3 \\ 1}t
                        +\colvec[r]{2 \\ 1 \\ 5}s
                     \suchthat t,s\in\Re} \)
            dan
            \( \set{\colvec[r]{4 \\ 7 \\ 7}m
                        +\colvec[r]{-4 \\ -2 \\ -10}n
                     \suchthat m,n\in\Re} \)
      \partsitem \( \set{\colvec[r]{1 \\ 2}t
                     \suchthat t\in\Re} \)
            dan
            \( \set{\colvec[r]{2 \\ 4}m
                        +\colvec[r]{4 \\ 8}n
                     \suchthat m,n\in\Re} \)
      \partsitem \( \set{\colvec[r]{1 \\ 0 \\ 2}s
                        +\colvec[r]{-1 \\ 1 \\ 0}t
                     \suchthat s,t\in\Re} \)
            dan
            \( \set{\colvec[r]{-1 \\ 1 \\ 1}m
                        +\colvec[r]{0 \\ 1 \\ 3}n
                     \suchthat m,n\in\Re} \)
      \partsitem \( \set{\colvec[r]{1 \\ 3 \\ 1}t
                        +\colvec[r]{2 \\ 4 \\ 6}s
                     \suchthat t,s\in\Re} \)
            dan
            \( \set{\colvec[r]{3 \\ 7 \\ 7}t
                        +\colvec[r]{1 \\ 3 \\ 1}s
                     \suchthat t,s\in\Re} \)
    \end{exparts}
    \begin{answer}
      Untuk setiap bagian, kita sebut himpunan pertama \( S_1 \) dan himpunan
      lainnya \( S_2 \).
     \begin{exparts}
      \partsitem Kedua himpunan sama.

        Untuk melihat bahwa \( S_1\subseteq S_2 \), kita harus menunjukkan
        bahwa setiap anggota himpunan pertama berada dalam himpunan kedua.
        Artinya, untuk setiap vektor berbentuk
        \begin{equation*}
          \vec{v}=\colvec[r]{1 \\ 2}
                  +\colvec[r]{0 \\ 3}t
        \end{equation*}
        terdapat \( s \) yang sesuai sedemikian sehingga
        \begin{equation*}
          \vec{v}=\colvec[r]{1 \\ 8}
                  +\colvec[r]{0 \\ -1}s.
        \end{equation*}
        Dengan kata lain, untuk setiap \( t \) yang diberikan, kita harus
        menemukan \( s \) yang membuat sistem berikut berlaku.
        \begin{equation*}
          \begin{linsys}{2}
            1  &+  &0s  &=  &1+0t\hfill  \\
            8  &-  &1s  &=  &2+3t\hfill  
          \end{linsys}
        \end{equation*}
        Sistem tersebut tereduksi menjadi
        \begin{equation*}
          \begin{linsys}{1}
            1  &= &1 \hfill \\
            s  &= &6-3t
          \end{linsys}
        \end{equation*}
        Artinya,
        \begin{equation*}
          \colvec[r]{1 \\ 2}
          +\colvec[r]{0 \\ 3}t
          =\colvec[r]{1 \\ 8}
          +\colvec[r]{0 \\ -1}(6-3t)
        \end{equation*}
        Dengan demikian, setiap vektor yang berbentuk seperti deskripsi
        \( S_1 \) juga dapat ditulis dalam bentuk yang diperlukan agar berada
        dalam \( S_2 \).

        Untuk menunjukkan \( S_2\subseteq S_1 \), kita mencari \( t \) yang
        membuat persamaan-persamaan berikut berlaku.
        \begin{equation*}
          \begin{linsys}{2}
            1  &+  &0t  &=  &1+0s\hfill  \\
            2  &+  &3t  &=  &8-1s\hfill  
          \end{linsys}
        \end{equation*}
        Sistem tersebut dapat ditulis ulang sebagai
        \begin{equation*}
          \begin{linsys}{1}
            1 &= &1\hfill   \\
            t &= &2-(1/3)s
          \end{linsys}
        \end{equation*}
        sehingga
        \begin{equation*}
          \colvec[r]{1 \\ 8}
          +\colvec[r]{0 \\ -1}s
          =\colvec[r]{1 \\ 2}
          +\colvec[r]{0 \\ 3}(2-(1/3)s).
        \end{equation*}
      \partsitem Kedua himpunan ini sama.

        Untuk menunjukkan bahwa \( S_1\subseteq S_2 \), kita periksa bahwa
        untuk setiap \( t,s \) kita dapat menemukan \( m,n \) yang sesuai
        sehingga persamaan-persamaan berikut berlaku.
        \begin{equation*}
          \begin{linsys}{2}
           4m  &-  &4n   &=  &1t+2s\hfill  \\
           7m  &-  &2n   &=  &3t+1s\hfill  \\
           7m  &-  &10n  &=  &1t+5s\hfill  
          \end{linsys}
        \end{equation*}
        Gunakan Metode Gauss
        \begin{align*}
          \begin{amat}{2}
            4  &-4  &1t+2s  \\
            7  &-2  &3t+1s  \\
            7  &-10 &1t+5s
          \end{amat}
          &\grstep[(-7/4)\rho_1+\rho_3]{(-7/4)\rho_1+\rho_2}
          \begin{amat}{2}
            4  &-4  &1t+2s           \\
            0  &5   &(5/4)t-(10/4)s  \\
            0  &-3  &-(3/4)t+(6/4)s
          \end{amat}                              \\
          &\grstep{(3/5)\rho_2+\rho_3}
          \begin{amat}{2}
            4  &-4  &1t+2s           \\
            0  &5   &(5/4)t-(10/4)s  \\
            0  &0   &0
          \end{amat}
        \end{align*}
        untuk menyimpulkan bahwa
        \begin{equation*}
          \colvec[r]{1 \\ 3 \\ 1}t
          +\colvec[r]{2 \\ 1 \\ 5}s
          =\colvec[r]{4 \\ 7 \\ 7}((1/2)t)
          +\colvec[r]{-4 \\ -2 \\ -10}((1/4)t-(1/2)s)
        \end{equation*}
        sehingga \( S_1\subseteq S_2 \).

        Untuk menunjukkan \( S_2\subseteq S_1 \), selesaikan
        \begin{equation*}
          \begin{linsys}{2}
           1t  &+  &2s   &=  &4m-4n\hfill  \\
           3t  &+  &1s   &=  &7m-2n\hfill  \\
           1t  &+  &5s   &=  &7m-10n\hfill  
          \end{linsys}
        \end{equation*}
        dengan reduksi Gauss
        \begin{align*}
          \begin{amat}{2}
            1  &2   &4m-4n  \\
            3  &1   &7m-2n  \\
            1  &5   &7m-10n
          \end{amat}
          &\grstep[-\rho_1+\rho_3]{-3\rho_1+\rho_2}
          \begin{amat}{2}
            1  &2   &4m-4n  \\
            0  &-5  &-5m+10n\\
            0  &3   &3m-6n
          \end{amat}                                    \\
          &\grstep{(3/5)\rho_2+\rho_3}
          \begin{amat}{2}
            1  &2   &4m-4n  \\
            0  &-5  &-5m+10n\\
            0  &0   &0
          \end{amat}
        \end{align*}
        untuk memperoleh
        \begin{equation*}
          \colvec[r]{4 \\ 7 \\ 7}m
          +\colvec[r]{-4 \\ -2 \\ -10}n
          =\colvec[r]{1 \\ 3 \\ 1}(2m)
          +\colvec[r]{2 \\ 1 \\ 5}(m-2n)
        \end{equation*}
        Dengan demikian, kita dapat menyatakan setiap anggota \( S_2 \) dalam
        bentuk yang diperlukan agar berada dalam \( S_1 \).
      \partsitem Kedua himpunan ini sama.

        Untuk membuktikan bahwa \( S_1\subseteq S_2 \), kita harus dapat
        menyelesaikan
        \begin{equation*}
          \begin{linsys}{2}
           2m  &+  &4n  &=  &1t\hfill  \\
           4m  &+  &8n  &=  &2t\hfill  
          \end{linsys}
        \end{equation*}
        untuk \( m \) dan \( n \) dalam bentuk \( t \).
        Terapkan reduksi Gauss
        \begin{equation*}
          \begin{amat}{2}
            2  &4   &1t  \\
            4  &8   &2t
          \end{amat}
          \onegrstep{-2\rho_1+\rho_2}
          \begin{amat}{2}
            2  &4   &1t  \\
            0  &0   &0
          \end{amat}
        \end{equation*}
        untuk menyimpulkan bahwa setiap pasangan \( m,n \) yang memenuhi
        \( 2m+4n=t \) dapat digunakan.
        Sebagai contoh,
        \begin{equation*}
          \colvec[r]{1 \\ 2}t
          =\colvec[r]{2 \\ 4}((1/2)t)
          +\colvec[r]{4 \\ 8}(0)
        \end{equation*}
        atau
        \begin{equation*}
          \colvec[r]{1 \\ 2}t
          =\colvec[r]{2 \\ 4}((-3/2)t)
          +\colvec[r]{4 \\ 8}(t).
        \end{equation*}
        Jadi, \( S_1\subseteq S_2 \).

        Untuk menunjukkan \( S_2\subseteq S_1 \), kita selesaikan
        \begin{equation*}
          \begin{linsys}{2}
           1t  &=  &2m+4n\hfill  \\
           2t  &=  &4m+8n\hfill  
          \end{linsys}
        \end{equation*}
        dengan Metode Gauss
        \begin{equation*}
          \begin{amat}{1}
            1  &2m+4n  \\
            2  &4m+8n
          \end{amat}
          \grstep{-2\rho_1+\rho_2}
          \begin{amat}{1}
            1  &2m+4n  \\
            0  &0
          \end{amat}
        \end{equation*}
        untuk menyimpulkan bahwa setiap vektor dalam \( S_2 \) juga berada
        dalam \( S_1 \).
        \begin{equation*}
          \colvec[r]{2 \\ 4}m
          +\colvec[r]{4 \\ 8}n
          =\colvec[r]{1 \\ 2}(2m+4n)\in S_1.
        \end{equation*}
      \partsitem Tidak satu pun dari kedua himpunan merupakan himpunan bagian
        dari yang lain.

        Agar \( S_1\subseteq S_2 \), kita harus dapat menyelesaikan
        \begin{equation*}
          \begin{linsys}{2}
          -1m  &+  &0n   &=  &1s-1t\hfill  \\
           1m  &+  &1n   &=  &0s+1t\hfill  \\
           1m  &+  &3n   &=  &2s+0t\hfill  
          \end{linsys}
        \end{equation*}
        untuk \( m \) dan \( n \) dalam bentuk \( t \) dan \( s \).
        Metode Gauss
        \begin{align*}
          \begin{amat}{2}
           -1  &0   &1s-1t  \\
            1  &1   &0s+1t  \\
            1  &3   &2s+0t
          \end{amat}
          &\grstep[\rho_1+\rho_3]{\rho_1+\rho_2}
          \begin{amat}{2}
           -1  &0   &1s-1t  \\
            0  &1   &1s+0t  \\
            0  &3   &3s-1t
          \end{amat}                        \\
          &\grstep{-3\rho_2+\rho_3}
          \begin{amat}{2}
           -1  &0   &1s-1t  \\
            0  &1   &1s+0t  \\
            0  &0   &0s-1t
          \end{amat}
        \end{align*}
        menunjukkan bahwa kita hanya dapat menemukan pasangan \( m,n \) yang
        sesuai ketika \( t=0 \).
        Artinya,
        \begin{equation*}
          \colvec[r]{-1 \\ 1 \\ 0}
        \end{equation*}
        tidak dapat dinyatakan dalam bentuk
        \begin{equation*}
           \colvec[r]{-1 \\ 1 \\ 1}m+\colvec[r]{0 \\ 1 \\ 3}n.
        \end{equation*}

        Setelah menunjukkan bahwa \( S_1 \) bukan himpunan bagian dari
        \( S_2 \), kita mengetahui bahwa \( S_1\neq S_2 \), sehingga sebenarnya
        kita tidak perlu melanjutkan.
        Namun, kita juga akan menunjukkan bahwa \( S_2 \) bukan himpunan bagian
        dari \( S_1 \).

        Agar \( S_2\subseteq S_1 \), kita harus dapat menyelesaikan
        \begin{equation*}
          \begin{linsys}{2}
           1s  &-  &1t   &=  &-1m+0n\hfill  \\
           0s  &+  &1t   &=  &1m+1n\hfill  \\
           2s  &+  &0t   &=  &1m+3n\hfill  
          \end{linsys}
        \end{equation*}
        untuk \( s \) dan \( t \).
        Terapkan reduksi baris
        \begin{align*}
          \begin{amat}{2}
            1  &-1  &-1m+0n  \\
            0  &1   &1m+1n  \\
            2  &0   &1m+3n
          \end{amat}
          &\grstep{-2\rho_1+\rho_3}
          \begin{amat}{2}
            1  &-1  &-1m+0n  \\
            0  &1   &1m+1n  \\
            0  &2   &3m+3n
          \end{amat}                                    \\
          &\grstep{-2\rho_2+\rho_3}
          \begin{amat}{2}
            1  &-1  &-1m+0n  \\
            0  &1   &1m+1n  \\
            0  &0   &1m+1n
          \end{amat}
        \end{align*}
        untuk menyimpulkan bahwa satu-satunya vektor dari \( S_2 \) yang juga
        berada dalam \( S_1 \) berbentuk
        \begin{equation*}
          \colvec[r]{-1 \\ 1 \\ 1}m
          +\colvec[r]{0 \\ 1 \\ 3}(-m).
        \end{equation*}
        Sebagai contoh,
        \begin{equation*}
          \colvec{-1 \\ 1 \\ 1}
        \end{equation*}
        berada dalam \( S_2 \), tetapi tidak dalam \( S_1 \).
      \partsitem Kedua himpunan ini sama.

        Pertama-tama, kita ubah parameter-parameternya:
        \begin{equation*}
          S_2=\set{\colvec[r]{3 \\ 7 \\ 7}m
                   +\colvec[r]{1 \\ 3 \\ 1}n
                   \suchthat m,n\in\Re}.
        \end{equation*}

        Selanjutnya, untuk menunjukkan bahwa \( S_1\subseteq S_2 \), kita
        selesaikan
        \begin{equation*}
          \begin{linsys}{2}
           3m  &+  &1n   &=  &1t+2s\hfill  \\
           7m  &+  &3n   &=  &3t+4s\hfill  \\
           7m  &+  &1n   &=  &1t+6s\hfill  
          \end{linsys}
        \end{equation*}
        dengan Metode Gauss
        \begin{align*}
          \begin{amat}{2}
            3  &1   &1t+2s  \\
            7  &3   &3t+4s  \\
            7  &1   &1t+6s
          \end{amat}
          &\grstep[(-7/3)\rho_1+\rho_3]{(-7/3)\rho_1+\rho_2}
          \begin{amat}{2}
            3  &1    &1t+2s  \\
            0  &2/3  &(2/3)t-(2/3)s  \\
            0  &-4/3 &(-4/3)t+(4/3)s
          \end{amat}                                    \\
          &\grstep{2\rho_2+\rho_3}
          \begin{amat}{2}
            3  &1    &1t+2s  \\
            0  &2/3  &(2/3)t-(2/3)s  \\
            0  &0    &0
          \end{amat}
        \end{align*}
        untuk memperoleh
        \begin{equation*}
          \colvec[r]{1 \\ 3 \\ 1}t
          +\colvec[r]{2 \\ 4 \\ 6}s
          =\colvec[r]{3 \\ 7 \\ 7}(s)
          +\colvec[r]{1 \\ 3 \\ 1}(t-s)
        \end{equation*}
        sehingga \( S_1\subseteq S_2 \).

        Pembuktian bahwa \( S_2\subseteq S_1 \) mengharuskan kita menyelesaikan
        \begin{equation*}
          \begin{linsys}{2}
           1t  &+  &2s   &=  &3m+1n\hfill  \\
           3t  &+  &4s   &=  &7m+3n\hfill  \\
           1t  &+  &6s   &=  &7m+1n\hfill  
          \end{linsys}
        \end{equation*}
        dengan reduksi Gauss
        \begin{align*}
          \begin{amat}{2}
            1  &2   &3m+1n  \\
            3  &4   &7m+3n  \\
            1  &6   &7m+1n
          \end{amat}
          &\grstep[-\rho_1+\rho_3]{-3\rho_1+\rho_2}
          \begin{amat}{2}
            1  &2   &3m+1n  \\
            0  &-2  &-2m    \\
            0  &4   &4m
          \end{amat}                                \\
          &\grstep{2\rho_2+\rho_3}
          \begin{amat}{2}
            1  &2   &3m+1n  \\
            0  &-2  &-2m    \\
            0  &0   &0
          \end{amat}
        \end{align*}
        untuk menyimpulkan bahwa
        \begin{equation*}
          \colvec[r]{3 \\ 7 \\ 7}m
          +\colvec[r]{1 \\ 3 \\ 1}n
          =\colvec[r]{1 \\ 3 \\ 1}(m+n)
          +\colvec[r]{2 \\ 4 \\ 6}(m)
        \end{equation*}
        sehingga setiap vektor dalam \( S_2 \) juga berada dalam \( S_1 \).
    \end{exparts}  
   \end{answer}
\end{exercises}
